% Fonctions logarithmes népérien — Mathématiques Tle Tech — Cours signé OnBuch+
% Programme officiel MINESEC. Compiler avec : tectonic N12-fonctions-logarithmes-neperien.tex
\newcommand{\DOCMATIERE}{Mathématiques}
\newcommand{\DOCNIVEAU}{Tle Tech}
\newcommand{\DOCMODULE}{Mathématiques – classe de Terminale technique}
\newcommand{\DOCLECON}{N12}
\newcommand{\DOCTITRE}{Fonctions logarithmes népérien}
% OnBuch+ — préambule « cours signés » ENRICHI (compilé avec tectonic / XeLaTeX)
% Système visuel « Pop » d'OnBuch+ : fond crème (jamais blanc), bordures encre
% épaisses, ombres « dures » décalées, titres Archivo Black en capitales.
% Version MATHÉMATIQUES : graphes (pgfplots/tikz), boîtes pédagogiques
% (définition, propriété, méthode, attention, le savais-tu, exercice résolu,
% exercice, TP, bilan), page de garde et signature OnBuch+ ancrée en bas.
%
% Macros attendues dans main.tex AVANT \input{preamble} :
%   \DOCMATIERE  \DOCNIVEAU  \DOCMODULE  \DOCLECON (numéro)  \DOCTITRE
\documentclass[11pt]{article}

\usepackage{fontspec}
\usepackage[french]{babel}
\usepackage[a4paper,margin=2cm,top=3.4cm,bottom=2.8cm,headheight=1.4cm,footskip=1.3cm]{geometry}
\usepackage{amsmath,amssymb,amsfonts}
\usepackage{mathtools} % \xmapsto, \coloneqq... souvent utilisés par les modèles
\usepackage{cancel} % simplifications barrées
% \abs{z} (module, valeur absolue) et \norm{u} (norme), utilisés par les modèles
\providecommand{\abs}{}\let\abs\relax\DeclarePairedDelimiter{\abs}{\lvert}{\rvert}
\providecommand{\norm}{}\let\norm\relax\DeclarePairedDelimiter{\norm}{\lVert}{\rVert}
\usepackage[table]{xcolor} % [table] : fournit aussi \rowcolors (alternance de lignes)
\usepackage{pagecolor}
\usepackage{tikz}
\usetikzlibrary{arrows.meta,positioning,calc,shapes.geometric,shapes.misc,shapes.arrows,decorations.pathmorphing,decorations.pathreplacing,decorations.markings,patterns,shadows,mindmap,trees,fit,backgrounds,matrix,intersections,angles,quotes}
\usepackage{pgfplots}
\usepackage{pgf-pie} % diagrammes circulaires \pie (statistiques)
\pgfplotsset{compat=1.18}
\usepgfplotslibrary{fillbetween,groupplots} % aires entre courbes (calcul intégral)
\usepackage{enumitem}
\usepackage{fancyhdr}
% [most] charge toutes les bibliothèques tcolorbox (listings, minted, raster,
% documentation...) et gonfle inutilement la mémoire TeX sur un document long
% (~300 boîtes) : on ne charge que ce qui est réellement utilisé.
\usepackage[skins,breakable]{tcolorbox}
\usepackage{graphicx}
\usepackage{colortbl}
\usepackage{booktabs}
\usepackage{tabularx}
\usepackage{diagbox} % case d'en-tête diagonale des tableaux à double entrée
% Colonnes extensibles centrée / gauche / droite (C, L, R), souvent utilisées par les modèles
\newcolumntype{C}{>{\centering\arraybackslash}X}
\newcolumntype{L}{>{\raggedright\arraybackslash}X}
\newcolumntype{R}{>{\raggedleft\arraybackslash}X}
\usepackage{array}
\usepackage{multicol}
\usepackage{siunitx}
% parse-numbers=false : accepte aussi \SI{2,5 \times 10^{-3}}{...}
\sisetup{locale=FR,output-decimal-marker={,},group-minimum-digits=4,parse-numbers=false}
\DeclareSIUnit\ohm{\ensuremath{\symup{\Omega}}} % U+2126 absent de la police
\usepackage{qrcode}
% \degree : commande fréquemment utilisée par les modèles pour un angle ou une
% température en dehors de \SI{}{} (non fournie par les paquets chargés).
\providecommand{\degree}{^{\circ}}
% \qed (fin de démonstration), utilisé par les modèles sans amsthm
\providecommand{\qed}{\hfill$\square$}
% \barre : double barre des valeurs interdites dans un tableau de variations (tkz-tab)
\providecommand{\barre}{\ensuremath{\|}}
% environnement proof (amsthm non chargé) utilisé par les modèles
\ifdefined\proof\else\newenvironment{proof}[1][Démonstration]{\par\noindent\textit{#1.}\ }{\qed\par}\fi
\usepackage{lastpage}
\usepackage{hyperref}
\hypersetup{hidelinks}

% ============================================================
% Palette « Pop » OnBuch+ — mêmes hex que la classe P (plus_ui.dart)
% ============================================================
\definecolor{poporange}{HTML}{F59321}
\definecolor{popdark}{HTML}{E07A0C}
\definecolor{popcream}{HTML}{FFF6E8}
\definecolor{popink}{HTML}{1C1714}
\definecolor{poppurple}{HTML}{7A5AE0}
\definecolor{popgreen}{HTML}{2E9E6A}
\definecolor{poppink}{HTML}{E0457B}
\definecolor{popblue}{HTML}{2F7DE1}
\definecolor{popgold}{HTML}{C98A12}
\definecolor{popmuted}{HTML}{8A7F76}
\definecolor{popline}{HTML}{EADFD2}
\definecolor{popgray}{HTML}{8A7F76}
\colorlet{popgrey}{popgray}
% Alias tolérants (le modèle nomme parfois les couleurs différemment)
\colorlet{popwhite}{white}
\colorlet{popblack}{popink}
\colorlet{popred}{poppink}
\colorlet{popyellow}{popgold}
% Teintes claires pour fonds de tableaux / zones
\colorlet{popblueL}{popblue!10!white}
\colorlet{poporangeL}{poporange!14!white}
\colorlet{popgreenL}{popgreen!12!white}
\colorlet{poppurpleL}{poppurple!12!white}
\colorlet{poppinkL}{poppink!12!white}
\colorlet{popgoldL}{popgold!14!white}

\pagecolor{popcream}

% ============================================================
% Typographie
% ============================================================
\defaultfontfeatures{Ligatures=TeX}
\setmainfont{PlusJakartaSans-Medium.ttf}[Path=../../../fonts/,BoldFont=PlusJakartaSans-Bold.ttf]
% Maths en sans-serif (Fira Math) avec chiffres et lettres droites de Plus
% Jakarta, pour que les formules se fondent dans le texte.
\usepackage[math-style=ISO,bold-style=ISO]{unicode-math}
\setmathfont{FiraMath-Regular.otf}[Path=../../../fonts/]
\setmathfont{PlusJakartaSans-Medium.ttf}[Path=../../../fonts/,range={up/num,up/{Latin,latin},\mathup}]
\usepackage{icomma} % virgule décimale sans espace en mode math
% Caractères mathématiques Unicode absents des polices de texte (Plus Jakarta,
% Archivo Black) : rendus par leur équivalent mathématique, en texte comme en
% formule — sinon ils s'affichent en carré vide (ex. « Arithmétique dans ℤ »).
\usepackage{newunicodechar}
\newunicodechar{ℝ}{\ensuremath{\mathbb{R}}}
\newunicodechar{ℤ}{\ensuremath{\mathbb{Z}}}
\newunicodechar{ℂ}{\ensuremath{\mathbb{C}}}
\newunicodechar{ℕ}{\ensuremath{\mathbb{N}}}
\newunicodechar{ℚ}{\ensuremath{\mathbb{Q}}}
\newunicodechar{𝔻}{\ensuremath{\mathbb{D}}}
\newunicodechar{α}{\ensuremath{\alpha}}
\newunicodechar{β}{\ensuremath{\beta}}
\newunicodechar{γ}{\ensuremath{\gamma}}
\newunicodechar{δ}{\ensuremath{\delta}}
\newunicodechar{ε}{\ensuremath{\varepsilon}}
\newunicodechar{θ}{\ensuremath{\theta}}
\newunicodechar{λ}{\ensuremath{\lambda}}
\newunicodechar{μ}{\ensuremath{\mu}}
\newunicodechar{σ}{\ensuremath{\sigma}}
\newunicodechar{τ}{\ensuremath{\tau}}
\newunicodechar{φ}{\ensuremath{\varphi}}
\newunicodechar{ψ}{\ensuremath{\psi}}
\newunicodechar{ω}{\ensuremath{\omega}}
\newunicodechar{Σ}{\ensuremath{\Sigma}}
% majuscules produites par \MakeUppercase dans les titres (α-aminés -> Α)
\newunicodechar{Α}{\ensuremath{\alpha}}
\newunicodechar{Β}{\ensuremath{\beta}}
\newunicodechar{Γ}{\ensuremath{\Gamma}}
\newunicodechar{Δ}{\ensuremath{\Delta}}
\newunicodechar{Ε}{\ensuremath{\varepsilon}}
\newunicodechar{Θ}{\ensuremath{\Theta}}
\newunicodechar{Λ}{\ensuremath{\Lambda}}
\newunicodechar{Μ}{\ensuremath{\mu}}
\newunicodechar{Τ}{\ensuremath{\tau}}
\newunicodechar{Φ}{\ensuremath{\Phi}}
\newunicodechar{Ψ}{\ensuremath{\Psi}}
\newunicodechar{Ω}{\ensuremath{\Omega}}
\newunicodechar{↦}{\ensuremath{\mapsto}}
\newunicodechar{∈}{\ensuremath{\in}}
\newunicodechar{∉}{\ensuremath{\notin}}
\newunicodechar{∧}{\ensuremath{\wedge}}
\newunicodechar{⇔}{\ensuremath{\Leftrightarrow}}
\newunicodechar{⟺}{\ensuremath{\Longleftrightarrow}}
\newunicodechar{⇒}{\ensuremath{\Rightarrow}}
\newunicodechar{⟹}{\ensuremath{\Longrightarrow}}
\newunicodechar{∘}{\ensuremath{\circ}}
\newunicodechar{∪}{\ensuremath{\cup}}
\newunicodechar{∩}{\ensuremath{\cap}}
\newunicodechar{≡}{\ensuremath{\equiv}}
\newunicodechar{⊂}{\ensuremath{\subset}}
\newunicodechar{⊥}{\ensuremath{\perp}}
\newunicodechar{∃}{\ensuremath{\exists}}
\newunicodechar{∀}{\ensuremath{\forall}}
\newunicodechar{∛}{\ensuremath{\sqrt[3]{\,}}}
\newunicodechar{∜}{\ensuremath{\sqrt[4]{\,}}}
\newunicodechar{⃗}{\ensuremath{{}^{\to}}}
\newunicodechar{ⁿ}{\ifmmode^{n}\else\textsuperscript{\ensuremath{n}}\fi}
\newunicodechar{ˣ}{\ifmmode^{x}\else\textsuperscript{\ensuremath{x}}\fi}
\newunicodechar{ᵇ}{\ifmmode^{b}\else\textsuperscript{\ensuremath{b}}\fi}
\newunicodechar{ʳ}{\ifmmode^{r}\else\textsuperscript{\ensuremath{r}}\fi}
\newunicodechar{ᵘ}{\ifmmode^{u}\else\textsuperscript{\ensuremath{u}}\fi}
\newunicodechar{ʸ}{\ifmmode^{y}\else\textsuperscript{\ensuremath{y}}\fi}
\newunicodechar{ᵃ}{\ifmmode^{a}\else\textsuperscript{\ensuremath{a}}\fi}
\newunicodechar{ᵉ}{\ifmmode^{e}\else\textsuperscript{\ensuremath{e}}\fi}
\newunicodechar{ᵏ}{\ifmmode^{k}\else\textsuperscript{\ensuremath{k}}\fi}
\newunicodechar{ᵖ}{\ifmmode^{p}\else\textsuperscript{\ensuremath{p}}\fi}
\newunicodechar{ᴺ}{\ifmmode^{N}\else\textsuperscript{\ensuremath{N}}\fi}
\newunicodechar{ᶜ}{\ifmmode^{c}\else\textsuperscript{\ensuremath{c}}\fi}
\newunicodechar{ʰ}{\ifmmode^{h}\else\textsuperscript{\ensuremath{h}}\fi}
\newunicodechar{ᵠ}{\ifmmode^{q}\else\textsuperscript{\ensuremath{q}}\fi}
\newunicodechar{ₙ}{\ifmmode_{n}\else\textsubscript{\ensuremath{n}}\fi}
\newunicodechar{ₐ}{\ifmmode_{a}\else\textsubscript{\ensuremath{a}}\fi}
\newunicodechar{ᵢ}{\ifmmode_{i}\else\textsubscript{\ensuremath{i}}\fi}
\newunicodechar{ᵣ}{\ifmmode_{r}\else\textsubscript{\ensuremath{r}}\fi}
\newunicodechar{ₖ}{\ifmmode_{k}\else\textsubscript{\ensuremath{k}}\fi}
\newunicodechar{ᵦ}{\ifmmode_{\beta}\else\textsubscript{\ensuremath{\beta}}\fi}
\newunicodechar{ₓ}{\ifmmode_{x}\else\textsubscript{\ensuremath{x}}\fi}
\newfontfamily\popdisplay{ArchivoBlack-Regular.ttf}[Path=../../../fonts/]
\newfontfamily\poplogo{SpaceGrotesk-Bold.ttf}[Path=../../../fonts/]
\newfontfamily\popsemi{PlusJakartaSans-SemiBold.ttf}[Path=../../../fonts/]
\newfontfamily\popxbold{PlusJakartaSans-ExtraBold.ttf}[Path=../../../fonts/]
\newfontfamily\popxboldls{PlusJakartaSans-ExtraBold.ttf}[Path=../../../fonts/,LetterSpace=3.0]

\setlist[enumerate]{leftmargin=1.7em,itemsep=5pt,topsep=4pt}
\setlist[itemize]{leftmargin=1.7em,itemsep=4pt,topsep=4pt,label={\color{poporange}\textbullet}}
\setlength{\parindent}{0pt}
\setlength{\parskip}{5pt}
\renewcommand{\arraystretch}{1.25}

% pgfplots : style Pop par défaut
\pgfplotsset{
  every axis/.append style={axis line style={popink,line width=1pt},tick style={popink},
    grid style={popline,line width=0.5pt},label style={font=\small},tick label style={font=\footnotesize}},
  popcurve/.style={poporange,line width=1.8pt,no markers,smooth},
}
% Même style disponible dans un \draw TikZ simple (hors pgfplots)
\tikzset{popcurve/.style={poporange,line width=1.8pt}}

% ============================================================
% Pill / squiggle
% ============================================================
\newcommand{\poppill}[3]{%
  \tikz[baseline=-0.55ex]\node[draw=popink,line width=1.2pt,rounded corners=2.6pt,%
    fill=#2,inner xsep=6.5pt,inner ysep=3pt]{%
    \popxboldls\fontsize{7}{7}\selectfont\color{#3}\MakeUppercase{#1}};%
}
\newcommand{\popsquiggle}[1]{%
  \tikz[baseline=-0.4ex]{
    \draw[#1,line width=2.6pt,line cap=round]
      plot [smooth] coordinates {(0,0.09) (0.34,0.01) (0.68,0.21) (1.02,0.01) (1.36,0.10)};
  }%
}
% Mot-clé surligné (marqueur orange sous le texte)
\newcommand{\cle}[1]{{\popxbold\color{popdark}#1}}

% ============================================================
% En-tête / pied
% ============================================================
\pagestyle{fancy}
\fancyhf{}
\renewcommand{\headrulewidth}{0pt}
\renewcommand{\footrulewidth}{0pt}
\lhead{\raisebox{-0.15ex}{\poplogo\fontsize{13}{13}\selectfont\color{popink}OnBuch\color{poporange}+}}
\rhead{\poppill{\DOCMATIERE\ \textbullet\ \DOCNIVEAU\ \textbullet\ Leçon \DOCLECON}{white}{popink}}
\lfoot{\popxbold\fontsize{7.6}{7.6}\selectfont\color{popmuted}\MakeUppercase{Cours signé OnBuch+}\ \textbullet\ {\popsemi\DOCTITRE}}
\rfoot{\tikz[baseline=-0.6ex]\node[rounded corners=8.5pt,fill=popink,minimum height=17pt,minimum width=17pt,inner xsep=5pt,inner sep=0]{\popxbold\fontsize{8}{8}\selectfont\color{popcream}\thepage\,/\,\pageref*{LastPage}};}

% ============================================================
% Titres
% ============================================================
\newcommand{\chaptitre}[2]{%
  \begin{tcolorbox}[enhanced,colback=poporange,colframe=popink,boxrule=2.6pt,arc=11pt,
      drop shadow=popink,left=15pt,right=15pt,top=12pt,bottom=11pt,before skip=3pt,after skip=18pt]
    \poppill{#2}{popink}{popcream}\par\vspace{9pt}
    {\raggedright\hyphenpenalty=10000\exhyphenpenalty=10000\popdisplay\fontsize{22}{24}\selectfont\color{popcream}\MakeUppercase{#1}\par}
  \end{tcolorbox}
}
% Section numérotée : pastille orange avec le numéro + titre Archivo
\newcounter{popsec}
\newcommand{\coursec}[1]{%
  \stepcounter{popsec}\setcounter{popsub}{0}%
  \par\vspace{16pt}\noindent
  \begin{minipage}{\linewidth}
  \tikz[baseline=-0.7ex]\node[draw=popink,line width=1.6pt,fill=poporange,rounded corners=4pt,minimum size=22pt,inner sep=0]{\popdisplay\fontsize{12}{12}\selectfont\color{popcream}\thepopsec};%
  \hspace{8pt}{\popdisplay\fontsize{15}{17}\selectfont\color{popink}\MakeUppercase{#1}}\par
  \vspace{4pt}\noindent{\color{poporange}\rule{36pt}{3.6pt}}
  \end{minipage}\par\nopagebreak\vspace{8pt}
}
\newcounter{popsub}
\newcommand{\courssub}[1]{%
  \stepcounter{popsub}%
  \par\vspace{9pt}\noindent{\popxbold\fontsize{11.5}{13}\selectfont\color{popink}{\color{poporange}\thepopsec.\thepopsub}\hspace{6pt}#1}\par\nopagebreak\vspace{4pt}
}

% ============================================================
% Boîtes pédagogiques — même recette : fond blanc, bordure épaisse colorée,
% ombre dure encre, badge plein. Titre optionnel [#1].
% ============================================================
\tcbset{popbase/.style={breakable,enhanced,colback=white,boxrule=2.3pt,arc=9pt,
  shadow={3pt}{-3pt}{0pt}{popink},left=14pt,right=14pt,top=11pt,bottom=11pt,before skip=11pt,after skip=15pt,
  fontupper=\popsemi\fontsize{10.6}{14.5}\selectfont\color{popink},
  fontlower=\popsemi\fontsize{10.6}{14.5}\selectfont\color{popink}}}
% Coche dessinée (le glyphe ✓ n'existe pas dans les polices du document)
\newcommand{\popcheck}[1]{\tikz[baseline=-0.5ex]\draw[#1,line width=1.6pt,line cap=round,line join=round](-0.13,0.0)--(-0.04,-0.09)--(0.13,0.1);}
\newcommand{\popboxhead}[3]{\poppill{#1}{#2}{white}\ifx\relax#3\relax\else\hspace{6pt}{\popxbold\fontsize{10.6}{12}\selectfont\color{popink}#3}\fi\par\vspace{7pt}}

\newtcolorbox{aretenir}[1][]{popbase,colframe=popgreen,before upper={\popboxhead{À retenir}{popgreen}{#1}}}
\newtcolorbox{exemplebox}[1][]{popbase,colframe=poporange,before upper={\popboxhead{Exemple}{poporange}{#1}}}
\newtcolorbox{definition}[1][]{popbase,colframe=popblue,before upper={\popboxhead{Définition}{popblue}{#1}}}
\newtcolorbox{propriete}[1][]{popbase,colframe=poppurple,before upper={\popboxhead{Propriété}{poppurple}{#1}}}
\newtcolorbox{methode}[1][]{popbase,colframe=popgold,before upper={\popboxhead{Méthode}{popgold}{#1}}}
\newtcolorbox{attention}[1][]{popbase,colframe=poppink,before upper={\popboxhead{Attention}{poppink}{#1}}}
\newtcolorbox{experience}[1][]{popbase,colframe=popblue,colback=popblueL,before upper={\popboxhead{Expérience}{popblue}{#1}}}
% « Le savais-tu ? » : carte violette pleine (contexte, culture, Cameroun)
\newtcolorbox{savaistu}[1][]{popbase,colback=poppurple,colframe=popink,
  fontupper=\popsemi\fontsize{10.4}{14.2}\selectfont\color{white},
  before upper={\poppill{Le savais-tu\,?}{white}{poppurple}\ifx\relax#1\relax\else\hspace{6pt}{\popxbold\color{white}#1}\fi\par\vspace{7pt}}}
% Exercice résolu : énoncé en haut, \tcblower puis la solution sur fond crème
\newcounter{popexo}\newcounter{popexr}
\newtcolorbox{exoresolu}[1][]{popbase,colframe=poporange,colbacklower=popcream,
  segmentation style={popink,line width=1.2pt,dashed},
  before upper={\stepcounter{popexr}\popboxhead{Exercice résolu \thepopexr}{poporange}{#1}},
  before lower={\poppill{Solution}{popink}{popcream}\par\vspace{6pt}}}
% Exercice (non résolu en place) — la correction est dans la partie Corrigés
\newtcolorbox{exercice}[1][]{popbase,colframe=popink,
  before upper={\stepcounter{popexo}\popboxhead{Exercice \thepopexo}{popink}{#1}}}
% Corrigé
\newtcolorbox{corrige}[1][]{popbase,colframe=popgreen,colback=popgreenL,
  before upper={\popboxhead{Corrigé}{popgreen}{#1}}}
% Objectifs / prérequis / bilan
\newtcolorbox{objectifs}{popbase,colframe=popink,colback=poporangeL,
  before upper={\popboxhead{Objectifs de la leçon}{popink}{}}}
\newtcolorbox{prerequis}{popbase,colframe=popink,colback=popblueL,
  before upper={\popboxhead{Prérequis}{popblue}{}}}
\newtcolorbox{bilan}{popbase,colframe=popink,colback=white,boxrule=2.8pt,
  before upper={\popboxhead{Fiche bilan}{poporange}{L'essentiel à connaître pour l'examen}}}
% Figure encadrée avec légende
\newtcolorbox{popfigure}[1][]{enhanced,breakable=false,colback=white,colframe=popline,boxrule=1.4pt,arc=8pt,
  left=10pt,right=10pt,top=10pt,bottom=8pt,before skip=10pt,after skip=12pt,halign=center,
  lower separated=false,halign lower=center,
  fontlower=\popsemi\fontsize{9}{11.5}\selectfont\color{popmuted},
  #1}
\newcommand{\legende}[1]{\par\vspace{6pt}{\popsemi\fontsize{9}{11.5}\selectfont\color{popmuted}#1}}

% ============================================================
% Signature OnBuch+
% ============================================================
\newcommand{\poplogoinline}[1]{{\poplogo\fontsize{#1}{#1}\selectfont\color{popink}OnBuch\color{poporange}+}}
\newcommand{\signatureonbuch}{%
  \begin{tcolorbox}[enhanced,colback=popink,colframe=popink,boxrule=2.2pt,arc=10pt,
      drop shadow=poporange,left=14pt,right=14pt,top=10pt,bottom=10pt,before skip=0pt,after skip=0pt]
    \tikz[baseline=-0.7ex]\node[circle,fill=popgreen,draw=popcream,line width=1.3pt,minimum size=22pt,inner sep=0]{\popcheck{white}};%
    \hspace{9pt}\begin{minipage}[c]{0.62\linewidth}
      {\popxbold\fontsize{12}{13}\selectfont\color{popcream}Signé\ \textbullet\ {\poplogo OnBuch\color{poporange}+}}\par\vspace{2pt}
      {\raggedright\popsemi\fontsize{8}{10.5}\selectfont\color{popline}\MakeUppercase{Équipe pédagogique OnBuch+}\\\MakeUppercase{Conforme au programme officiel MINESEC}\par}
    \end{minipage}\hfill
    {\poplogo\fontsize{22}{22}\selectfont\color{popcream}OnBuch\color{poporange}+}
  \end{tcolorbox}%
}

% ============================================================
% Page de garde
% #1 titre · #2 matière · #3 niveau · #4 module · #5 n° leçon · #6 durée ·
% #7 description / objectifs (texte ou itemize)
% ============================================================
\newcommand{\pagegarde}[7]{%
  \thispagestyle{empty}
  \newgeometry{a4paper,margin=1.7cm,top=1.6cm,bottom=1.4cm}%
  \begin{tikzpicture}[remember picture,overlay]
    \fill[poporange!18] ([xshift=-3.2cm,yshift=-3.6cm]current page.north east) circle (4.6cm);
    \fill[poppurple!14] ([xshift=2.4cm,yshift=7.6cm]current page.south west) circle (3.8cm);
  \end{tikzpicture}%
  {\poplogo\fontsize{30}{30}\selectfont\color{popink}OnBuch\color{poporange}+}\hfill
  \raisebox{0.6ex}{\poppill{Cours signé \textbullet\ Premium}{popink}{popcream}}\par
  \vspace{1pt}\popsquiggle{poppurple}\par
  \vspace{8pt}
  \begin{tcolorbox}[enhanced,colback=poporange,colframe=popink,boxrule=2.8pt,arc=13pt,
      drop shadow=popink,left=18pt,right=18pt,top=12pt,bottom=13pt,after skip=10pt]
    \poppill{#2 \textbullet\ #3}{popink}{popcream}\hspace{5pt}\poppill{#4}{white}{popink}\par\vspace{8pt}
    {\popdisplay\fontsize{14}{14}\selectfont\color{popink}LEÇON #5}\par\vspace{4pt}
    {\raggedright\hyphenpenalty=10000\exhyphenpenalty=10000\popdisplay\fontsize{25}{27}\selectfont\color{popcream}\MakeUppercase{#1}\par}
    \vspace{8pt}
    {\popxbold\fontsize{9.5}{9.5}\selectfont\color{popink}\MakeUppercase{Durée conseillée : #6}}
  \end{tcolorbox}
  \begin{tcolorbox}[enhanced,colback=white,colframe=popink,boxrule=2.2pt,arc=10pt,
      drop shadow=popink,left=16pt,right=16pt,top=10pt,bottom=10pt,after skip=8pt]
    \poppill{Dans cette leçon}{poppurple}{white}\par\vspace{5pt}
    \popsemi\fontsize{9.3}{12.6}\selectfont\color{popink}#7
  \end{tcolorbox}
  \noindent\begin{minipage}[t]{0.487\linewidth}
    \begin{tcolorbox}[enhanced,colback=popgreen,colframe=popink,boxrule=2.2pt,arc=10pt,
        drop shadow=popink,left=11pt,right=11pt,top=10pt,bottom=10pt,height=3.1cm,valign=center]
      \tcbox[colback=white,colframe=popink,boxrule=1.3pt,arc=5pt,boxsep=0pt,nobeforeafter,
        left=3pt,right=3pt,top=3pt,bottom=3pt,valign=center]{\qrcode[height=1.9cm]{https://play.google.com/store/apps/details?id=com.luvvix.onbuch&hl=fr}}%
      \hspace{8pt}\begin{minipage}[c]{0.5\linewidth}
        {\popdisplay\fontsize{11}{12.5}\selectfont\color{white}\MakeUppercase{Télécharge OnBuch}}\par\vspace{4pt}
        {\raggedright\popsemi\fontsize{8.4}{11}\selectfont\color{white}Gratuit \textbullet\ Google Play\\Tuteur IA, annales, concours, résultats\par}
      \end{minipage}
    \end{tcolorbox}
  \end{minipage}\hfill
  \begin{minipage}[t]{0.487\linewidth}
    \begin{tcolorbox}[enhanced,colback=poppurple,colframe=popink,boxrule=2.2pt,arc=10pt,
        drop shadow=popink,left=11pt,right=11pt,top=10pt,bottom=10pt,height=3.1cm,valign=center]
      \tikz[baseline=-0.7ex]\node[circle,fill=white,draw=popink,line width=1.3pt,minimum size=1.6cm,inner sep=0]{\popdisplay\fontsize{24}{24}\selectfont\color{poppurple}+};%
      \hspace{8pt}\begin{minipage}[c]{0.6\linewidth}
        {\popdisplay\fontsize{11}{12.5}\selectfont\color{white}\MakeUppercase{Passe à OnBuch+}}\par\vspace{4pt}
        {\raggedright\popsemi\fontsize{8.4}{11}\selectfont\color{white}Tous les cours signés, Tuteur IA illimité, fiches PDF — onglet OnBuch+ de l'app\par}
      \end{minipage}
    \end{tcolorbox}
  \end{minipage}
  \vspace{8pt}
  \popcontenu
  \vfill
  \signatureonbuch
  \restoregeometry
  \newpage
}

% Rangée « Ce cours contient » (page de garde)
\newcommand{\poptile}[3]{% #1 couleur #2 gros texte #3 légende
  \begin{tcolorbox}[enhanced,colback=#1,colframe=popink,boxrule=2pt,arc=9pt,shadow={2.5pt}{-2.5pt}{0pt}{popink},
      left=6pt,right=6pt,top=9pt,bottom=9pt,halign=center,valign=center,height=1.75cm,width=\linewidth,nobeforeafter]
    {\popdisplay\fontsize{12.5}{13}\selectfont\color{white}\MakeUppercase{#2}}\par\vspace{3pt}
    {\centering\popsemi\fontsize{7.8}{9.5}\selectfont\color{white}#3\par}
  \end{tcolorbox}}
\newcommand{\popcontenu}{%
  \par\noindent{\popxboldls\fontsize{7.5}{7.5}\selectfont\color{popmuted}CE COURS SIGNÉ CONTIENT}\par\vspace{7pt}
  \noindent\begin{minipage}[t]{0.235\linewidth}\poptile{popblue}{Cours}{complet, illustré, pas à pas}\end{minipage}\hfill
  \begin{minipage}[t]{0.235\linewidth}\poptile{poporange}{Méthodes}{et exercices résolus}\end{minipage}\hfill
  \begin{minipage}[t]{0.235\linewidth}\poptile{poppink}{Exercices}{type examen + corrigés}\end{minipage}\hfill
  \begin{minipage}[t]{0.235\linewidth}\poptile{popgreen}{Fiche bilan}{l'essentiel à retenir}\end{minipage}\par}

% Sommaire
\newtcolorbox{sommaire}{enhanced,colback=white,colframe=popink,boxrule=2.2pt,arc=10pt,
  drop shadow=popink,left=18pt,right=18pt,top=13pt,bottom=13pt,before skip=6pt,after skip=20pt,
  before upper={\poppill{Sommaire}{poppurple}{white}\par\vspace{9pt}\popsemi\fontsize{10.6}{14.5}\selectfont\color{popink}}}

% Signature de fin de cours — poussée tout en bas de la dernière page
\newcommand{\findecours}{%
  \par\vfill\nopagebreak
  \begin{center}\popsquiggle{poporange}\end{center}\vspace{4pt}
  \signatureonbuch
}
\AtBeginDocument{\def\llbracket{\mathopen{[\mkern-3.5mu[}}\def\rrbracket{\mathclose{]\mkern-3.5mu]}}} % intervalles d'entiers (glyphe ⟦ absent de la police)
% Glyphes absents de Fira Math : redéfinis à partir de glyphes disponibles
\AtBeginDocument{%
  \def\setminus{\mathbin{\backslash}}\let\smallsetminus\setminus
  \def\vdots{\mathord{\vbox{\baselineskip3.2pt\lineskiplimit0pt\kern4pt\hbox{.}\hbox{.}\hbox{.}}}}%
  \def\ddots{\mathinner{\mkern1mu\raise6.5pt\hbox{.}\mkern2mu\raise3.6pt\hbox{.}\mkern2mu\raise0.7pt\hbox{.}\mkern1mu}}%
  \def\triangle{\mathord{\scalebox{0.9}{$\symup{\Delta}$}}}%
  \def\nabla{\mathord{\raisebox{\depth}{\scalebox{1}[-1]{$\symup{\Delta}$}}}}%
  \def\checkmark{\popcheck{popdark}}%
  \def\star{\mathbin{\ast}}\def\bigstar{\mathord{\ast}}%
  \def\diamond{\mathbin{\scalebox{0.6}{$\Diamond$}}}%
  \def\bigcup{\mathop{\mathchoice{\raisebox{-0.3em}{\scalebox{1.7}{$\cup$}}}{\scalebox{1.25}{$\cup$}}{\cup}{\cup}}\displaylimits}%
  \def\bigcap{\mathop{\mathchoice{\raisebox{-0.3em}{\scalebox{1.7}{$\cap$}}}{\scalebox{1.25}{$\cap$}}{\cap}{\cap}}\displaylimits}%
  \def\lesseqqgtr{\mathrel{\lessgtr}}\def\bigoplus{\mathop{\oplus}\displaylimits}%
}
\providecommand{\ssi}{si et seulement si} % abréviation fréquente
\AtBeginDocument{\providecommand{\wideparen}{\overparen}} % arc de cercle

%% FIGFIX-BEGIN v3 — correctifs globaux des figures (docs/onbuch-generation/tools/figfix)
\usepackage{adjustbox}\usepackage{etoolbox}
% toute figure TikZ/pgfplots plus large que la ligne est réduite à la largeur disponible
\BeforeBeginEnvironment{tikzpicture}{\begin{adjustbox}{max width=\linewidth}}
\AfterEndEnvironment{tikzpicture}{\end{adjustbox}}
% légendes pgfplots lisibles par-dessus les courbes
\pgfplotsset{every axis/.append style={legend style={fill=white,fill opacity=0.92,text opacity=1,draw=black!15,inner sep=3pt}}}
% étiquettes d'arêtes (syntaxe edge["..."]) avec fond blanc
\tikzset{every edge quotes/.append style={fill=white,inner sep=1.2pt}}
% bulles de cartes mentales : texte à la ligne dans la bulle, bulle un peu plus grande
\tikzset{every concept/.append style={text width=2.0cm,align=center,minimum size=2.3cm,inner sep=2pt}}
%% FIGFIX-END
\begin{document}

\pagegarde{Fonctions logarithmes népérien}{Mathématiques}{Tle Tech}{Mathématiques – classe de Terminale technique}{N12}{5 séances (fiche de progression harmonisée)}{Cette leçon introduit la fonction logarithme népérien ln comme réciproque de l'exponentielle, établit ses propriétés algébriques, ses limites, sa dérivée et les primitives de la forme u'/u. Elle aboutit à l'étude complète de fonctions faisant intervenir ln et à la résolution de problèmes concrets de croissance et de modélisation.
\vspace{4pt}
\begin{multicols}{2}\small
\begin{itemize}
\item À la fin de cette leçon, je sais définir la fonction ln comme réciproque de l'exponentielle et déterminer l'ensemble de définition d'une fonction contenant ln
\item À la fin de cette leçon, je sais utiliser les propriétés algébriques de ln (produit, quotient, puissance) pour transformer et simplifier des expressions
\item À la fin de cette leçon, je sais résoudre des équations et inéquations contenant ln ou e\^{}x
\item À la fin de cette leçon, je sais calculer les limites usuelles faisant intervenir ln, y compris les croissances comparées
\item À la fin de cette leçon, je sais dériver les fonctions définies à l'aide de ln, en particulier ln(u) et u·ln
\item À la fin de cette leçon, je sais déterminer les primitives des fonctions de la forme u'/u
\end{itemize}\end{multicols}}

\begin{sommaire}
\begin{multicols}{2}
\begin{enumerate}
\item Définition et premières propriétés de la fonction ln
\item Propriétés algébriques du logarithme népérien
\item Limites faisant intervenir ln
\item Dérivation des fonctions définies à l'aide de ln
\item Primitives des fonctions de la forme u'/u
\item Étude de fonctions faisant intervenir ln et applications concrètes
\item Activité d'intégration
\item Méthodes et exercices résolus
\item Exercices
\item Corrigés des exercices
\item Fiche bilan
\end{enumerate}
\end{multicols}
\end{sommaire}

\begin{objectifs}
\begin{itemize}
\item À la fin de cette leçon, je sais définir la fonction ln comme réciproque de l'exponentielle et déterminer l'ensemble de définition d'une fonction contenant ln
\item À la fin de cette leçon, je sais utiliser les propriétés algébriques de ln (produit, quotient, puissance) pour transformer et simplifier des expressions
\item À la fin de cette leçon, je sais résoudre des équations et inéquations contenant ln ou e\^{}x
\item À la fin de cette leçon, je sais calculer les limites usuelles faisant intervenir ln, y compris les croissances comparées
\item À la fin de cette leçon, je sais dériver les fonctions définies à l'aide de ln, en particulier ln(u) et u·ln
\item À la fin de cette leçon, je sais déterminer les primitives des fonctions de la forme u'/u
\item À la fin de cette leçon, je sais dresser le tableau de variations complet et tracer la courbe d'une fonction faisant intervenir ln
\item À la fin de cette leçon, je sais modéliser et résoudre une situation concrète de la vie courante (croissance, pH, capital) à l'aide du logarithme népérien
\item À la fin de cette leçon, je sais rédiger et justifier rigoureusement une démarche, une réponse ou une conclusion
\end{itemize}
\end{objectifs}

\begin{prerequis}
\begin{itemize}
\item Fonction exponentielle e\^{}x : définition, propriétés, dérivée, courbe (leçon précédente de Terminale)
\item Notion de fonction réciproque et de bijection (composée et symétrie des courbes par rapport à y = x)
\item Limites de fonctions et opérations sur les limites (Première et début de Terminale)
\item Dérivation : dérivée des fonctions composées, produit, quotient (Première)
\item Primitives usuelles et linéarité de l'intégration (début de Terminale)
\end{itemize}
\end{prerequis}

\newpage

\coursec{Définition et premières propriétés de la fonction ln}

Monsieur Nkoulou, promoteur d'une porcherie moderne à Obala, suit de près la fermentation de ses déchets organiques pour produire du biogaz. Il constate que la population bactérienne dans sa cuve double toutes les 3 heures. Modélisée par une exponentielle $N(t) = N_0 e^{kt}$, cette croissance le conduit à se poser une question cruciale pour dimensionner son installation : \emph{au bout de combien de temps la population initiale sera-t-elle multipliée par 100 ? Par 1000 ?} Résoudre $e^{kt} = 100$ ou $e^{kt} = 1000$ revient à « inverser » la fonction exponentielle. C'est précisément le rôle d'une nouvelle fonction, le \textbf{logarithme népérien}, notée $\ln$. Au-delà de l'élevage, cette fonction permet de mesurer le pH d'une solution en chimie ($\text{pH} = -\log_{10}[H^+] = -\frac{\ln[H^+]}{\ln 10}$), l'intensité d'un séisme (échelle de Richter) ou le temps de doublement d'un capital placé à la microfinance. Dans cette section, nous la définissons rigoureusement, en déduisons ses premières propriétés et apprenons à l'utiliser pour résoudre équations et inéquations.

\courssub{Rappel : la fonction exponentielle, une bijection}

En classe de Première, nous avons établi que la fonction exponentielle $\exp : \mathbb{R} \to ]0 ; +\infty[$, $x \mapsto e^x$ (souvent notée $e^x$) est :
\begin{itemize}
    \item \textbf{strictement croissante} sur $\mathbb{R}$ (si $x_1 < x_2$ alors $e^{x_1} < e^{x_2}$) ;
    \item \textbf{continue} sur $\mathbb{R}$ ;
    \item \textbf{bijective} de $\mathbb{R}$ sur $]0 ; +\infty[$ : pour tout $y > 0$, il existe un unique $x \in \mathbb{R}$ tel que $e^x = y$.
\end{itemize}
Cette bijectivité est la clé qui permet de définir sa fonction réciproque.

\begin{popfigure}
\begin{tikzpicture}[scale=0.9, every node/.style={font=\small}]
    % Axes
    \draw[->] (-3.5,0) -- (3.5,0) node[right] {$x$};
    \draw[->] (0,-1.5) -- (0,3.5) node[above] {$y$};
    % Line y=x
    \draw[popmuted, dashed, thick] (-3,-3) -- (3,3) node[above right, pos=0.9] {$y=x$};
    % Exp curve
    \draw[popblue, thick, domain=-3:1.1, samples=100] plot (\x, {exp(\x)});
    \node[popblue, anchor=south west] at (1, 2.7) {$y = e^x$};
    % Ln curve (reflection of exp across y=x)
    \draw[poporange, thick, domain=0.05:3, samples=100] plot ({ln(\x)}, \x);
    \node[poporange, anchor=north east] at (2.5, 1) {$y = \ln x$};
    % Points
    \fill[popblue] (0,1) circle (2pt) node[above left] {$(0,1)$};
    \fill[poporange] (1,0) circle (2pt) node[below right] {$(1,0)$};
    \fill[popblue] (1,2.718) circle (2pt) node[above right] {$(1,e)$};
    \fill[poporange] (2.718,1) circle (2pt) node[below left] {$(e,1)$};
    % Symmetry arrows
    \draw[poppurple, <->, thick] (0.5, 1.65) -- (1.65, 0.5);
    \node[poppurple, rotate=-45] at (1.3, 1.3) {symétrie};
    % Ticks
    \foreach \x/\label in {1/1, 2.718/e} \draw (\x, 0.05) -- (\x, -0.05) node[below] {$\label$};
    \foreach \y/\label in {1/1, 2.718/e} \draw (0.05, \y) -- (-0.05, \y) node[left] {$\label$};
\end{tikzpicture}
\legende{Courbes de $y = e^x$ et $y = \ln x$ : symétriques par rapport à la droite $y = x$.}
\end{popfigure}

\courssub{Définition du logarithme népérien}

\begin{definition}[Fonction logarithme népérien]
La fonction \textbf{logarithme népérien}, notée $\ln$, est la \textbf{bijection réciproque} de la fonction exponentielle $\exp$. Elle est définie par :
\[
\ln : ]0 ; +\infty[ \longrightarrow \mathbb{R}, \quad y \longmapsto x \quad \text{si et seulement si} \quad e^x = y.
\]
En d'autres termes, pour tout $x > 0$, $\ln x$ est l'unique nombre réel tel que $e^{\ln x} = x$.
\end{definition}

\begin{popfigure}
\begin{tikzpicture}[node distance=1.5cm, >=Stealth, every node/.style={font=\small}]
    % Sets
    \node[draw, rounded corners, fill=popblueL, minimum width=3cm, minimum height=2cm] (R) at (0,0) {$\mathbb{R}$};
    \node[draw, rounded corners, fill=poporangeL, minimum width=3cm, minimum height=2cm] (Rplus) at (6,0) {$]0 ; +\infty[$};
    % Elements in R
    \foreach \y/\label in {0.6/x, 0/-1, -0.6/0, -1.2/2} {
        \fill (0.5,\y) circle (1.5pt);
        \node[left] at (-0.8,\y) {$\label$};
    }
    % Elements in R+
    \foreach \y/\label in {0.6/e^x, 0/0.37, -0.6/1, -1.2/7.39} {
        \fill (5.5,\y) circle (1.5pt);
        \node[right] at (7.3,\y) {$\label$};
    }
    % Arrows exp
    \draw[popblue, thick, ->] (1.5,0.6) -- (4.5,0.6) node[midway, above] {$\exp : x \mapsto e^x$};
    \draw[popblue, thick, ->] (1.5,-0.6) -- (4.5,-0.6);
    \draw[popblue, thick, ->] (1.5,-1.2) -- (4.5,-1.2);
    % Arrows ln
    \draw[poporange, thick, <-] (1.5,0.6) -- (4.5,0.6);
    \draw[poporange, thick, <-] (1.5,-0.6) -- (4.5,-0.6) node[midway, below] {$\ln : y \mapsto \ln y$};
    \draw[poporange, thick, <-] (1.5,-1.2) -- (4.5,-1.2);
    % Labels
    \node[popblue, above] at (0,1.3) {Domaine de départ : $\mathbb{R}$};
    \node[poporange, above] at (6,1.3) {Ensemble d'arrivée : $]0;+\infty[$};
\end{tikzpicture}
\legende{Schéma de la bijection réciproque : $\exp$ et $\ln$ s'annulent mutuellement.}
\end{popfigure}

\noindent De cette définition découlent immédiatement les \textbf{formules d'inversion} (ou formules fondamentales) :

\begin{propriete}[Formules d'inversion]
Pour tout réel $x$ et pour tout réel strictement positif $y$ :
\begin{align*}
    e^{\ln y} &= y \qquad (y > 0) \\
    \ln(e^x) &= x \qquad (x \in \mathbb{R})
\end{align*}
Ces égalités traduisent le fait que $\exp$ et $\ln$ sont bijections réciproques l'une de l'autre.
\end{propriete}

\noindent Deux valeurs remarquables s'en déduisent sans calcul :
\begin{itemize}
    \item $\ln 1 = 0$ car $e^0 = 1$.
    \item $\ln e = 1$ car $e^1 = e$.
\end{itemize}

\begin{savaistu}[Origine de la notation et du nombre $e$]
La notation $\ln$ vient du latin \emph{logarithmus naturalis} (« logarithme naturel »). Le nombre $e \approx 2,718\,281\,828\ldots$ a été étudié par Jacob Bernoulli (intérêts composés) puis formalisé par Leonhard \textbf{Euler} au XVIIIᵉ siècle, qui lui donna la lettre $e$ (initiale de \emph{exponential}). C'est la base « naturelle » car la dérivée de $e^x$ est $e^x$ elle-même, ce qui simplifie infiniment le calcul différentiel et intégral.
\end{savaistu}

\courssub{Variations et injectivité de $\ln$}

Puisque $\ln$ est la réciproque d'une fonction strictement croissante, elle hérite de cette monotonie.

\begin{propriete}[Variations et injectivité de $\ln$ (admise ici, démontrée en Section 4 via la dérivée)]
La fonction $\ln$ est \textbf{continue} et \textbf{strictement croissante} sur son ensemble de définition $]0 ; +\infty[$.
\end{propriete}

\noindent Cette monotonie stricte entraîne l'équivalence fondamentale pour résoudre équations et inéquations :

\begin{propriete}[Équivalences d'ordre et d'égalité]
Pour tous réels strictement positifs $a$ et $b$ :
\begin{align*}
    \ln a = \ln b &\iff a = b \\
    \ln a < \ln b &\iff a < b \\
    \ln a > \ln b &\iff a > b
\end{align*}
\emph{Attention :} Ces équivalences ne sont valides que si $a>0$ et $b>0$ (condition d'existence de $\ln$).
\end{propriete}

\courssub{Ensemble de définition d'une fonction contenant $\ln(u(x))$}

Toute expression $\ln(u(x))$ n'a de sens que si son argument est strictement positif.

\begin{methode}[Déterminer l'ensemble de définition de $f(x) = \ln(u(x))$]
\begin{enumerate}
    \item Écrire la condition d'existence : $u(x) > 0$ (inéquation \textbf{stricte}).
    \item Résoudre cette inéquation dans $\mathbb{R}$.
    \item L'ensemble des solutions est l'ensemble de définition $\mathcal{D}_f$.
\end{enumerate}
\end{methode}

\begin{exemplebox}[Ensemble de définition]
Déterminer $\mathcal{D}_f$ pour $f(x) = \ln(2x - 4)$.
\begin{enumerate}
    \item Condition : $2x - 4 > 0$.
    \item Résolution : $2x > 4 \iff x > 2$.
    \item $\mathcal{D}_f = ]2 ; +\infty[$.
\end{enumerate}
\end{exemplebox}

\begin{exemplebox}[Ensemble de définition (cas quadratique)]
Déterminer $\mathcal{D}_g$ pour $g(x) = \ln(x^2 - 4)$.
\begin{enumerate}
    \item Condition : $x^2 - 4 > 0 \iff (x-2)(x+2) > 0$.
    \item Résolution : $x < -2$ ou $x > 2$.
    \item $\mathcal{D}_g = ]-\infty ; -2[ \cup ]2 ; +\infty[$.
\end{enumerate}
\end{exemplebox}

\courssub{Résolution d'équations et d'inéquations faisant intervenir $\ln$ et $\exp$}

La stratégie repose sur l'utilisation des formules d'inversion, \textbf{en vérifiant systématiquement les conditions d'existence}.

\begin{methode}[Résoudre une équation ou inéquation avec $\ln$ et $\exp$]
\begin{enumerate}
    \item \textbf{Identifier les conditions d'existence (C.E.)} : tout argument de $\ln$ doit être $>0$.
    \item \textbf{Transformer l'équation/inéquation} en utilisant :
          \begin{itemize}
              \item $\ln A = k \iff A = e^k$ (avec $A>0$) ;
              \item $e^X = k \iff X = \ln k$ (avec $k>0$) ;
              \item $\ln A = \ln B \iff A = B$ (avec $A>0, B>0$).
          \end{itemize}
    \item \textbf{Résoudre} l'équation ou l'inéquation algébrique obtenue.
    \item \textbf{Sélectionner} les solutions qui satisfont les C.E. de l'étape 1.
    \item \textbf{Conclure} en donnant l'ensemble solution.
\end{enumerate}
\end{methode}

\begin{exoresolu}[Résolution d'une équation logarithmique]
Résoudre dans $\mathbb{R}$ : $\ln(2x - 4) = 3$.
\medskip
\noindent\textbf{Solution détaillée.}
\begin{enumerate}
    \item \textbf{C.E.} : $2x - 4 > 0 \iff x > 2$.
    \item \textbf{Transformation} : $\ln(2x - 4) = 3 \iff 2x - 4 = e^3$ (car $e^3 > 0$, la C.E. est automatiquement satisfaite par la solution).
    \item \textbf{Résolution} : $2x = e^3 + 4 \iff x = \dfrac{e^3 + 4}{2}$.
    \item \textbf{Vérification C.E.} : $e^3 \approx 20,085$, donc $x \approx \frac{24,085}{2} \approx 12,04 > 2$. La solution est admissible.
    \item \textbf{Conclusion} : $S = \left\{ \dfrac{e^3 + 4}{2} \right\} \approx \{12,04\}$.
\end{enumerate}
\end{exoresolu}

\begin{exoresolu}[Résolution d'une équation mixte $\ln$ et $\exp$]
Résoudre dans $\mathbb{R}$ : $e^{2x} - 3e^x + 2 = 0$.
\medskip
\noindent\textbf{Solution détaillée.}
\begin{enumerate}
    \item Pas de $\ln$ dans l'équation, pas de C.E. initiale. Posons $X = e^x$ ($X > 0$).
    \item $X^2 - 3X + 2 = 0 \iff (X-1)(X-2) = 0 \iff X = 1$ ou $X = 2$.
    \item Retour à $x$ : $e^x = 1 \iff x = \ln 1 = 0$ ; $e^x = 2 \iff x = \ln 2$.
    \item Les deux valeurs sont dans $\mathbb{R}$.
    \item $S = \{0 ; \ln 2\}$.
\end{enumerate}
\end{exoresolu}

\begin{exemplebox}[Résolution d'une inéquation]
Résoudre : $\ln(x + 3) \le 1$.
\begin{enumerate}
    \item \textbf{C.E.} : $x + 3 > 0 \iff x > -3$.
    \item $\ln(x+3) \le 1 \iff \ln(x+3) \le \ln e \iff x+3 \le e$ (car $\ln$ est croissante).
    \item $x \le e - 3 \approx -0,28$.
    \item Intersection avec C.E. : $-3 < x \le e - 3$.
    \item $S = ]-3 ; e - 3]$.
\end{enumerate}
\end{exemplebox}

\begin{attention}[Pièges classiques à éviter]
\begin{itemize}
    \item \textbf{Oublier les conditions d'existence} : écrire $\ln(-2) = \dots$ ou résoudre $\ln(x-1)=0$ en trouvant $x=1$ (or $\ln 0$ n'existe pas).
    \item \textbf{Confondre $\ln x = k$ et $\ln x = \ln k$} : $\ln x = k \iff x = e^k$ ; $\ln x = \ln k \iff x = k$ (si $k>0$).
    \item \textbf{Appliquer $\ln$ des deux côtés sans vérifier la positivité} : $e^x = -2$ n'a \textbf{pas} de solution (car $e^x > 0$), on n'écrit pas $x = \ln(-2)$.
    \item \textbf{Notation} : $\ln^2 x$ signifie $(\ln x)^2$, pas $\ln(\ln x)$.
    \item \textbf{Changement de variable} : Pour $e^{2x} - 3e^x + 2 = 0$, poser $X=e^x$ impose $X>0$. Ne pas oublier de rejeter les racines négatives du trinôme.
\end{itemize}
\end{attention}

\begin{savaistu}[Applications concrètes au Cameroun et ailleurs]
\begin{itemize}
    \item \textbf{Chimie / Environnement} : Le pH d'une eau de forage ou d'un effluent industriel se calcule par $\text{pH} = -\log_{10}[H^+] = -\frac{\ln[H^+]}{\ln 10}$. Une eau de pH 4 a une concentration en $H^+$ 100 fois plus grande qu'une eau de pH 6.
    \item \textbf{Finance / Microfinance} : Pour un capital $C_0$ placé à un taux $t$ (composé continu), le temps pour doubler est $T = \frac{\ln 2}{t}$. À 6\%/an, $T \approx \frac{0,693}{0,06} \approx 11,5$ ans.
    \item \textbf{Sismologie} : L'échelle de Richter utilise un logarithme (base 10) : une magnitude 6 libère 10 fois plus d'énergie qu'une magnitude 5.
    \item \textbf{Acoustique} : Le niveau sonore en décibels : $L = 10 \log_{10}(I/I_0) = \frac{10}{\ln 10} \ln(I/I_0)$.
\end{itemize}
\end{savaistu}

\coursec{Propriétés algébriques du logarithme népérien}

Les propriétés algébriques du logarithme népérien découlent directement des propriétés de l'exponentielle et des formules d'inversion. Elles sont indispensables pour simplifier les expressions, dériver et intégrer.

\begin{propriete}[Propriétés algébriques fondamentales]
Pour tous réels strictement positifs $a$ et $b$, et pour tout réel $r$ :
\begin{align*}
    \ln(ab) &= \ln a + \ln b \quad &\text{(logarithme d'un produit)} \\
    \ln\left(\frac{a}{b}\right) &= \ln a - \ln b \quad &\text{(logarithme d'un quotient)} \\
    \ln(a^r) &= r \ln a \quad &\text{(logarithme d'une puissance)}
\end{align*}
\end{propriete}

\begin{experience}[Démonstration guidée de $\ln(ab) = \ln a + \ln b$]
\begin{enumerate}
    \item Soient $a>0, b>0$. Posons $A = \ln a$ et $B = \ln b$. Par définition, $e^A = a$ et $e^B = b$.
    \item Calculons $e^{A+B}$ : $e^{A+B} = e^A \cdot e^B = a \cdot b = ab$.
    \item Appliquons $\ln$ des deux côtés : $\ln(e^{A+B}) = \ln(ab)$.
    \item Par formule d'inversion, $A+B = \ln(ab)$.
    \item En remplaçant $A$ et $B$ : $\ln a + \ln b = \ln(ab)$. \qed
\end{enumerate}
\end{experience}

\begin{exemplebox}[Simplification d'expressions]
Simplifier : $E = \ln 8 - \ln 2 + 2\ln 3 - \ln 6$.
\begin{align*}
    E &= \ln\left(\frac{8}{2}\right) + \ln(3^2) - \ln 6 \\
      &= \ln 4 + \ln 9 - \ln 6 \\
      &= \ln\left(\frac{4 \times 9}{6}\right) = \ln\left(\frac{36}{6}\right) = \ln 6.
\end{align*}
\end{exemplebox}

\begin{attention}[Erreurs fréquentes sur les propriétés algébriques]
\begin{itemize}
    \item $\ln(a+b) \neq \ln a + \ln b$ (pas de distributivité sur la somme).
    \item $\ln(a-b) \neq \ln a - \ln b$.
    \item $\frac{\ln a}{\ln b} \neq \ln\left(\frac{a}{b}\right)$.
    \item $\ln(a^r) = r\ln a$ nécessite $a>0$. Pour $a<0$ et $r$ entier, on utilise $\ln|a|$ (voir primitives).
\end{itemize}
\end{attention}

\coursec{Limites faisant intervenir ln}

La connaissance des limites de $\ln$ aux bornes de son ensemble de définition et à l'infini est cruciale pour l'étude de fonctions.

\begin{propriete}[Limites remarquables de $\ln$]
\begin{align*}
    \lim_{x \to 0^+} \ln x &= -\infty \\
    \lim_{x \to +\infty} \ln x &= +\infty
\end{align*}
\end{propriete}

\begin{propriete}[Croissance comparée (admise au Probatoire)]
Pour tout $\alpha > 0$ :
\begin{align*}
    \lim_{x \to +\infty} \frac{\ln x}{x^\alpha} &= 0 \quad \text{(\emph{« toute puissance bat le logarithme » à $+\infty$})} \\
    \lim_{x \to 0^+} x^\alpha \ln x &= 0 \quad \text{(\emph{« toute puissance bat le logarithme » en $0^+$})}
\end{align*}
\end{propriete}

\begin{methode}[Calcul de limites de fonctions composées avec $\ln$]
\begin{enumerate}
    \item Calculer la limite de l'expression intérieure $u(x)$.
    \item Si $\lim u(x) = \ell \in ]0;+\infty[$ : $\lim \ln(u(x)) = \ln \ell$ (continuité).
    \item Si $\lim u(x) = 0^+$ : $\lim \ln(u(x)) = -\infty$.
    \item Si $\lim u(x) = +\infty$ : $\lim \ln(u(x)) = +\infty$.
    \item Si forme indéterminée (ex: $\frac{\ln x}{x}$), utiliser les croissances comparées ou l'Hôpital (si au programme, sinon raisonner par encadrement).
\end{enumerate}
\end{methode}

\begin{exoresolu}[Calcul de limites]
Calculer les limites en $0^+$ et $+\infty$ de $f(x) = x \ln x$ et $g(x) = \frac{\ln x}{x}$.
\medskip
\noindent\textbf{Solution détaillée.}
\begin{itemize}
    \item \textbf{$f(x) = x \ln x$} :
          \begin{itemize}
              \item $x \to 0^+$ : forme $0 \times (-\infty)$. Posons $X = 1/x \to +\infty$. $f(x) = -\frac{\ln X}{X} \to 0$ (croissance comparée). $\lim_{0^+} f = 0$.
              \item $x \to +\infty$ : $+\infty \times +\infty = +\infty$.
          \end{itemize}
    \item \textbf{$g(x) = \frac{\ln x}{x}$} :
          \begin{itemize}
              \item $x \to 0^+$ : $\frac{-\infty}{0^+} = -\infty$.
              \item $x \to +\infty$ : forme $\frac{+\infty}{+\infty}$. Croissance comparée ($\alpha=1$) : $\lim_{+\infty} g = 0$.
          \end{itemize}
\end{itemize}
\end{exoresolu}

\begin{exemplebox}[Limite de $\ln(u(x))$]
$\lim_{x \to 2} \ln(x^2 - 3)$. $u(x)=x^2-3 \to 1$. $\ln 1 = 0$.
$\lim_{x \to 1^+} \ln(x-1)$. $u(x)=x-1 \to 0^+$. $\ln(0^+) = -\infty$.
\end{exemplebox}

\coursec{Dérivation des fonctions définies à l'aide de ln}

La dérivée de $\ln$ est la pierre angulaire du calcul différentiel sur cette fonction.

\begin{propriete}[Dérivée de la fonction logarithme népérien]
La fonction $\ln$ est dérivable sur $]0;+\infty[$ et pour tout $x>0$ :
\[
(\ln x)' = \frac{1}{x}.
\]
\emph{Preuve (exigible au Probatoire) :} Par dérivation de la bijection réciproque. $y = \ln x \iff x = e^y$. En dérivant $x = e^y$ par rapport à $x$ : $1 = e^y \cdot y' = x \cdot y' \implies y' = 1/x$.
\end{propriete}

\begin{propriete}[Dérivée de $\ln(u(x))$ — Formule de la chaîne]
Soit $u$ une fonction dérivable et strictement positive sur un intervalle $I$. La fonction $x \mapsto \ln(u(x))$ est dérivable sur $I$ et :
\[
(\ln(u(x)))' = \frac{u'(x)}{u(x)}.
\]
\end{propriete}

\begin{methode}[Dériver une fonction faisant intervenir $\ln$]
\begin{enumerate}
    \item Identifier la forme : $\ln(u(x))$, produit, quotient, somme faisant intervenir $\ln$.
    \item Appliquer les règles de dérivation (somme, produit, quotient, chaîne).
    \item Simplifier l'expression obtenue (factoriser, mettre au même dénominateur).
    \item Préciser le domaine de la dérivée (identique au domaine de la fonction).
\end{enumerate}
\end{methode}

\begin{exoresolu}[Dérivation]
Soit $f(x) = \ln(3x^2 + 1)$ sur $\mathbb{R}$.
$u(x) = 3x^2+1 > 0$, $u'(x) = 6x$.
$f'(x) = \frac{6x}{3x^2+1}$.

Soit $g(x) = x \ln x$ sur $]0;+\infty[$.
$g'(x) = 1 \cdot \ln x + x \cdot \frac{1}{x} = \ln x + 1$.

Soit $h(x) = \frac{\ln x}{x}$ sur $]0;+\infty[$.
$h'(x) = \frac{\frac{1}{x} \cdot x - \ln x \cdot 1}{x^2} = \frac{1 - \ln x}{x^2}$.
\end{exoresolu}

\begin{exemplebox}[Recherche d'extremum via la dérivée]
Étudier les variations de $f(x) = \frac{\ln x}{x}$ sur $]0;+\infty[$.
$f'(x) = \frac{1-\ln x}{x^2}$. Le dénominateur $>0$. Signe de $f'$ = signe de $1-\ln x$.
$f'(x) = 0 \iff \ln x = 1 \iff x = e$.
\begin{itemize}
    \item $]0;e[$ : $\ln x < 1 \implies f'>0$ : $f$ croissante.
    \item $]e;+\infty[$ : $\ln x > 1 \implies f'<0$ : $f$ décroissante.
\end{itemize}
Maximum en $x=e$, $f(e) = \frac{1}{e}$.
\begin{center}
\begin{tabular}{|c|ccccc|}\hline
$x$ & $0^+$ & & $e$ & $+\infty$ \\ \hline
$f'(x)$ & & $+$ & $0$ & $-$ \\ \hline
$f$ & $-\infty$ & $\nearrow$ & $\frac{1}{e}$ & $\searrow$ & $0$ \\ \hline
\end{tabular}
\end{center}
\end{exemplebox}

\coursec{Primitives des fonctions de la forme u'/u}

L'intégration est l'opération inverse de la dérivation. La formule $(\ln u)' = u'/u$ donne immédiatement une primitive.

\begin{propriete}[Primitive de $\frac{u'}{u}$]
Soit $u$ une fonction dérivable et strictement positive sur un intervalle $I$. Une primitive de $\frac{u'}{u}$ sur $I$ est $\ln(u(x))$.
Plus généralement, si $u$ est dérivable et ne s'annule pas sur $I$, une primitive de $\frac{u'}{u}$ est $\ln|u(x)|$.
\end{propriete}

\begin{attention}[Valeur absolue]
La primitive $\ln|u(x)|$ est valide sur tout intervalle où $u$ ne s'annule pas. Si l'intervalle n'est pas précisé ou si $u$ change de signe, on utilise $\ln|u|$. Au Probatoire, on précise souvent le domaine pour lever l'ambiguïté.
\end{attention}

\begin{methode}[Reconnaître et intégrer une forme $\frac{u'}{u}$]
\begin{enumerate}
    \item Identifier le numérateur comme la dérivée (ou un multiple constant de la dérivée) du dénominateur.
    \item Si $f(x) = k \frac{u'(x)}{u(x)}$, alors une primitive est $k \ln|u(x)| + C$.
    \item Vérifier par dérivation.
\end{enumerate}
\end{methode}

\begin{exoresolu}[Calcul de primitives]
\begin{enumerate}
    \item $\int \frac{2x}{x^2+1}\,dx$. $u=x^2+1, u'=2x$. Forme $\frac{u'}{u}$. Primitive : $\ln(x^2+1) + C$ (pas de valeur absolue car $x^2+1>0$).
    \item $\int \frac{\cos x}{\sin x}\,dx$ sur $]0;\pi[$. $u=\sin x, u'=\cos x$. Primitive : $\ln(\sin x) + C$ ($\sin x > 0$ sur $]0;\pi[$).
    \item $\int \frac{3}{2x-1}\,dx = \frac{3}{2} \int \frac{2}{2x-1}\,dx = \frac{3}{2} \ln|2x-1| + C$.
    \item $\int \tan x\,dx = \int \frac{\sin x}{\cos x}\,dx = -\ln|\cos x| + C = \ln|\sec x| + C$.
\end{enumerate}
\end{exoresolu}

\begin{exemplebox}[Intégrale définie et aire]
Calculer l'aire sous la courbe $y = \frac{1}{x}$ entre $x=1$ et $x=e$.
$\mathcal{A} = \int_1^e \frac{1}{x}\,dx = [\ln x]_1^e = \ln e - \ln 1 = 1 - 0 = 1$ unité d'aire.
\end{exemplebox}

\coursec{Étude de fonctions faisant intervenir ln et applications concrètes}

Cette section synthétise toutes les notions précédentes pour mener une étude complète de fonction (domaine, limites, dérivée, variations, tableau, courbe) et résoudre des problèmes concrets.

\begin{methode}[Étude complète d'une fonction $f$ faisant intervenir $\ln$]
\begin{enumerate}
    \item \textbf{Domaine de définition $\mathcal{D}_f$} : résoudre $u(x)>0$ pour chaque $\ln(u(x))$.
    \item \textbf{Parité / Périodicité} : généralement ni pair ni impair, non périodique.
    \item \textbf{Limites aux bornes de $\mathcal{D}_f$} : utiliser les limites de $\ln$ et croissances comparées. Chercher les asymptotes (verticales, horizontales, obliques).
    \item \textbf{Dérivée $f'$} : calculer et simplifier (souvent forme $\frac{P(x)}{x Q(x)}$ ou $\frac{P(\ln x)}{x}$).
    \item \textbf{Signe de $f'$ et Tableau de variations} : résoudre $f'(x)=0$, étudier le signe sur $\mathcal{D}_f$.
    \item \textbf{Valeurs remarquables / Intersections} : $f(x)=0$, $f(1)$ si $1\in\mathcal{D}_f$.
    \item \textbf{Tracé de la courbe} : repère orthonormé, asymptotes, points caractéristiques.
\end{enumerate}
\end{methode}

\begin{exoresolu}[{Étude complète : $f(x) = x \ln x$ sur $]0;+\infty[$}]
\medskip
\noindent\textbf{Solution détaillée.}
\begin{enumerate}
    \item $\mathcal{D}_f = ]0;+\infty[$.
    \item \textbf{Limites} :
          $\lim_{0^+} x\ln x = 0$ (croissance comparée). $\lim_{+\infty} x\ln x = +\infty$.
          Asymptote en $0$ ? Non, limite finie. Pas d'asymptote horizontale ni oblique en $+\infty$ (croissance super-linéaire).
    \item \textbf{Dérivée} : $f'(x) = \ln x + 1$.
    \item \textbf{Variations} : $f'(x)=0 \iff \ln x = -1 \iff x = e^{-1} = 1/e \approx 0,368$.
          \begin{itemize}
              \item $]0; 1/e[$ : $\ln x < -1 \implies f'<0$ : $f$ décroissante.
              \item $]1/e; +\infty[$ : $\ln x > -1 \implies f'>0$ : $f$ croissante.
          \end{itemize}
          Minimum en $x=1/e$, $f(1/e) = -1/e \approx -0,368$.
    \item \textbf{Zéros} : $x\ln x = 0 \iff x=1$ (car $x>0$).
    \item \textbf{Tableau de variations} :
\begin{center}
\begin{tabular}{|c|ccccccc|}\hline
$x$ & $0^+$ & & $1/e$ & & $1$ & & $+\infty$ \\ \hline
$f'(x)$ & & $-$ & $0$ & $+$ & $+$ & $+$ & \\ \hline
$f$ & $0$ & $\searrow$ & $-1/e$ & $\nearrow$ & $0$ & $\nearrow$ & $+\infty$ \\ \hline
\end{tabular}
\end{center}
\end{enumerate}
\end{exoresolu}

\begin{exoresolu}[Problème concret : Croissance bactérienne (Reprise de l'introduction)]
La population bactérienne de la cuve de M. Nkoulou suit le modèle $N(t) = N_0 e^{kt}$ avec $t$ en heures. On sait que la population double en 3 heures ($N(3) = 2N_0$).
\begin{enumerate}
    \item Déterminer $k$.
    \item Au bout de combien de temps la population est-elle multipliée par 100 ? Par 1000 ?
    \item Tracer la courbe de $t \mapsto \frac{N(t)}{N_0}$.
\end{enumerate}
\medskip
\noindent\textbf{Solution détaillée.}
\begin{enumerate}
    \item $N(3) = N_0 e^{3k} = 2N_0 \iff e^{3k} = 2 \iff 3k = \ln 2 \iff k = \frac{\ln 2}{3} \approx 0,231\,\text{h}^{-1}$.
    \item Multipliée par 100 : $N(t) = 100 N_0 \iff e^{kt} = 100 \iff kt = \ln 100 = 2\ln 10 \iff t = \frac{2\ln 10}{k} = \frac{2\ln 10}{\ln 2/3} = 6 \frac{\ln 10}{\ln 2} \approx 19,93\,\text{h}$.
    Multipliée par 1000 : $t = \frac{\ln 1000}{k} = \frac{3\ln 10}{\ln 2/3} = 9 \frac{\ln 10}{\ln 2} \approx 29,9\,\text{h}$.
    \item Courbe exponentielle passant par $(0,1)$, $(3,2)$, $(19,93, 100)$, $(29,9, 1000)$. Croissance de plus en plus rapide.
\end{enumerate}
\end{exoresolu}

\begin{exoresolu}[Application : pH d'une solution]
La concentration en ions hydrogène $[H^+]$ (en mol/L) d'une solution est $3,2 \times 10^{-4}$. Calculer son pH.
$\text{pH} = -\log_{10}[H^+] = -\frac{\ln(3,2 \times 10^{-4})}{\ln 10} = -\frac{\ln 3,2 + \ln 10^{-4}}{\ln 10} = -\frac{\ln 3,2 - 4\ln 10}{\ln 10} = 4 - \frac{\ln 3,2}{\ln 10}$.
$\ln 3,2 \approx 1,163$, $\ln 10 \approx 2,303$. $\text{pH} \approx 4 - 0,505 = 3,495$. La solution est acide (pH < 7).
\end{exoresolu}

\begin{exoresolu}[Application : Finance - Intérêts composés continus]
Un capital de 500 000 FCFA est placé à un taux annuel de 5\% composé en continu. Quelle somme au bout de 10 ans ? Au bout de combien de temps le capital double-t-il ?
Formule : $C(t) = C_0 e^{rt}$ avec $r=0,05$.
\begin{enumerate}
    \item $C(10) = 500\,000 \times e^{0,05 \times 10} = 500\,000 \times e^{0,5} \approx 500\,000 \times 1,6487 = 824\,360\,\text{FCFA}$.
    \item Doublement : $C(t) = 2C_0 \iff e^{0,05t} = 2 \iff 0,05t = \ln 2 \iff t = \frac{\ln 2}{0,05} \approx \frac{0,693}{0,05} = 13,86\,\text{ans}$.
\end{enumerate}
\end{exoresolu}

\begin{exercice}[Entraînement au Probatoire]
\begin{enumerate}
    \item Résoudre dans $\mathbb{R}$ : $\ln(x^2 - 5x + 6) = 0$.
    \item Étudier la fonction $f(x) = \ln(x^2 + 1)$ sur $\mathbb{R}$ (domaine, limites, dérivée, variations, tableau, courbe).
    \item Calculer les primitives sur $]1;+\infty[$ de $g(x) = \frac{2x-1}{x^2-x}$.
    \item Une culture bactérienne suit $P(t) = 1000 e^{0,2t}$ ($t$ en heures). Déterminer le temps nécessaire pour atteindre 5000 bactéries.
    \item Soit $h(x) = \frac{\ln x}{x^2}$ sur $]0;+\infty[$. Montrer que $h$ admet un maximum et le déterminer.
\end{enumerate}
\end{exercice}

\begin{corrige}[Exercice n°1]
$\ln(x^2-5x+6)=0 \iff x^2-5x+6 = 1 \iff x^2-5x+5=0$. $\Delta = 25-20=5$. $x = \frac{5\pm\sqrt{5}}{2}$.
C.E. : $x^2-5x+6 > 0 \iff (x-2)(x-3)>0 \iff x<2$ ou $x>3$.
$\frac{5-\sqrt{5}}{2} \approx 1,38 < 2$ (OK). $\frac{5+\sqrt{5}}{2} \approx 3,62 > 3$ (OK).
$S = \left\{ \frac{5-\sqrt{5}}{2} ; \frac{5+\sqrt{5}}{2} \right\}$.
\end{corrige}

\begin{corrige}[Exercice n°2]
$f(x)=\ln(x^2+1)$. $\mathcal{D}_f=\mathbb{R}$ ($x^2+1\ge 1>0$).
Pair : $f(-x)=f(x)$.
Limites : $+\infty \to +\infty$.
$f'(x) = \frac{2x}{x^2+1}$. Signe de $f'$ = signe de $x$.
Décroissante sur $]-\infty;0]$, croissante sur $[0;+\infty[$. Minimum en $0$, $f(0)=0$.
Tableau : $-\infty \searrow 0 \nearrow +\infty$.
Courbe : forme de U, tangente horizontale en $(0,0)$.
\end{corrige}

\begin{corrige}[Exercice n°3]
$g(x) = \frac{2x-1}{x^2-x} = \frac{(x^2-x)'}{x^2-x}$. $u=x^2-x>0$ sur $]1;+\infty[$.
Primitive : $G(x) = \ln(x^2-x) + C = \ln(x(x-1)) + C$.
\end{corrige}

\begin{corrige}[Exercice n°4]
$1000 e^{0,2t} = 5000 \iff e^{0,2t} = 5 \iff 0,2t = \ln 5 \iff t = 5\ln 5 \approx 5 \times 1,609 = 8,05\,\text{h}$.
\end{corrige}

\begin{corrige}[Exercice n°5]
$h(x) = \frac{\ln x}{x^2}$. $h'(x) = \frac{\frac{1}{x}x^2 - \ln x \cdot 2x}{x^4} = \frac{x(1 - 2\ln x)}{x^4} = \frac{1-2\ln x}{x^3}$.
$h'(x)=0 \iff 1-2\ln x=0 \iff \ln x = 1/2 \iff x = \sqrt{e}$.
$]0;\sqrt{e}[$ : $h'>0$ (croissante). $]\sqrt{e};+\infty[$ : $h'<0$ (décroissante).
Maximum en $x=\sqrt{e}$, $h(\sqrt{e}) = \frac{1/2}{e} = \frac{1}{2e}$.
\end{corrige}

\coursec{Propriétés algébriques du logarithme népérien}

Dans la section précédente, nous avons défini la fonction \cle{logarithme népérien} $\ln$ comme la \emph{réciproque} de la fonction exponentielle : pour tout $x>0$, $\ln x$ est l'unique nombre $y$ tel que $e^{y}=x$. Nous avons retenu les deux relations fondamentales
\[ e^{\ln x}=x\quad (x>0) \qquad \text{et} \qquad \ln(e^{x})=x\quad (x\in\mathbb{R}). \]
Nous allons maintenant découvrir ce qui fait toute la puissance du logarithme : il \textbf{transforme les multiplications en additions}, les quotients en soustractions et les puissances en produits. C'est exactement cette idée qui a rendu les calculs possibles avant l'invention des calculatrices, et c'est elle qui servira, dans les sections suivantes, à dériver, à intégrer et à résoudre des équations.

\courssub{La propriété fondamentale : $\ln(ab)=\ln a+\ln b$}

\textbf{Intuition.} Dire que $\ln x=y$ signifie \og $x$ est $e$ à la puissance $y$ \fg{}. Multiplier deux nombres revient à \emph{additionner leurs puissances de $e$} : c'est la règle bien connue $e^{A}\times e^{B}=e^{A+B}$. Le logarithme, qui \og lit \fg{} les puissances, transforme donc naturellement un produit en somme.

\begin{propriete}[Propriété fondamentale (théorème)]
Pour tous réels $a>0$ et $b>0$ :
\[ \ln(ab)=\ln a+\ln b. \]
Cette propriété est exigible au Probatoire et au Baccalauréat technique ; sa démonstration est courte et peut être demandée.
\end{propriete}

\begin{experience}[Démonstration guidée de $\ln(ab)=\ln a+\ln b$]
Suivons la démarche pas à pas, en utilisant uniquement la définition de $\ln$ comme réciproque de l'exponentielle.
\begin{enumerate}
\item \textbf{On pose} $A=\ln a$ et $B=\ln b$. Par définition du logarithme, cela signifie :
\[ a=e^{A} \qquad \text{et} \qquad b=e^{B}. \]
\item \textbf{On calcule le produit} $ab$ en utilisant la propriété des exponentielles :
\[ ab=e^{A}\times e^{B}=e^{A+B}. \]
\item \textbf{On applique $\ln$ aux deux membres.} Comme $\ln(e^{x})=x$ pour tout réel $x$ :
\[ \ln(ab)=\ln\!\left(e^{A+B}\right)=A+B. \]
\item \textbf{Conclusion.} En remplaçant $A$ et $B$ par leurs valeurs :
\[ \ln(ab)=\ln a+\ln b. \qquad \blacksquare \]
\end{enumerate}
L'idée-clé est donc : \emph{le produit $ab$ s'écrit $e^{A+B}$, et le logarithme récupère l'exposant $A+B$.}
\end{experience}

\begin{exemplebox}[Vérification numérique]
Prenons $a=2$ et $b=8$. Alors $ab=16$. Avec la calculatrice :
\[ \ln 2\approx 0{,}693, \qquad \ln 8\approx 2{,}079, \qquad \ln 16\approx 2{,}773. \]
On vérifie : $0{,}693+2{,}079=2{,}772\approx\ln 16$. La formule fonctionne.
\end{exemplebox}

\begin{attention}[Le piège de la somme : $\ln(a+b)$ ne se simplifie pas]
L'erreur la plus fréquente consiste à écrire $\ln(a+b)=\ln a+\ln b$. C'est \textbf{faux} ! Par exemple :
\[ \ln(2+3)=\ln 5\approx 1{,}609 \qquad \text{alors que} \qquad \ln 2+\ln 3=\ln 6\approx 1{,}792. \]
Le logarithme possède des propriétés pour les \textbf{produits}, les \textbf{quotients} et les \textbf{puissances}, mais \emph{aucune} propriété simple pour une \textbf{somme}. De même, $\ln(a-b)$ ne se transforme pas. Retenir : on ne peut \og casser \fg{} un logarithme que sur $\times$, $\div$ et les puissances.
\end{attention}

\courssub{Conséquences : inverse, quotient et puissances}

Toutes les formules de cette partie se \emph{déduisent} de la propriété fondamentale. Il faut savoir refaire ces déductions rapidement.

\begin{propriete}[Logarithme d'un inverse et d'un quotient]
Pour tous réels $a>0$ et $b>0$ :
\[ \ln\!\left(\frac{1}{a}\right)=-\ln a \qquad \text{et} \qquad \ln\!\left(\frac{a}{b}\right)=\ln a-\ln b. \]
\end{propriete}

\textbf{Démonstration.} Pour l'inverse, on part de $a\times\dfrac{1}{a}=1$. En appliquant $\ln$ :
\[ \ln a+\ln\!\left(\frac{1}{a}\right)=\ln 1=0, \qquad \text{donc} \qquad \ln\!\left(\frac{1}{a}\right)=-\ln a. \]
Pour le quotient, on écrit $\dfrac{a}{b}=a\times\dfrac{1}{b}$ et on applique la propriété fondamentale :
\[ \ln\!\left(\frac{a}{b}\right)=\ln a+\ln\!\left(\frac{1}{b}\right)=\ln a-\ln b. \qquad \blacksquare \]

\begin{propriete}[Logarithme d'une puissance]
Pour tout réel $a>0$ et tout entier naturel $n$ :
\[ \ln(a^{n})=n\,\ln a. \]
En particulier, pour la racine carrée : $\ln(\sqrt{a})=\dfrac{1}{2}\ln a$.
\emph{Complément utile (hors exigibilité stricte mais très pratique)} : la formule se généralise à tout exposant réel $r$ : $\ln(a^{r})=r\ln a$.
\end{propriete}

\textbf{Démonstration (idée).} Pour $n$ entier, $a^{n}=a\times a\times\cdots\times a$ ($n$ facteurs). En appliquant la propriété fondamentale $n-1$ fois, on obtient une somme de $n$ termes égaux à $\ln a$, donc $n\ln a$. Pour la racine carrée, on remarque que $\sqrt{a}\times\sqrt{a}=a$, donc $2\ln(\sqrt{a})=\ln a$, d'où le résultat. $\blacksquare$

\begin{exemplebox}[Applications directes]
\begin{itemize}
\item $\ln 8=\ln(2^{3})=3\ln 2$ ;
\item $\ln\!\left(\dfrac{1}{5}\right)=-\ln 5$ ;
\item $\ln(\sqrt{3})=\dfrac{1}{2}\ln 3$ ;
\item $\ln(e^{2})=2\ln e=2$ (on retrouve une propriété de la section 1).
\end{itemize}
\end{exemplebox}

\begin{aretenir}[Tableau récapitulatif des formules algébriques de $\ln$]
Pour tous $a>0$, $b>0$ et tout entier $n$ :
\begin{center}
\rowcolors{1}{popblueL}{white}
\begin{tabularx}{\linewidth}{|l|X|}
\hline
\rowcolor{popblue!40} \textbf{Opération} & \textbf{Formule} \\
\hline
Produit & $\ln(ab)=\ln a+\ln b$ \\
\hline
Inverse & $\ln\!\left(\dfrac{1}{a}\right)=-\ln a$ \\
\hline
Quotient & $\ln\!\left(\dfrac{a}{b}\right)=\ln a-\ln b$ \\
\hline
Puissance & $\ln(a^{n})=n\ln a$ \\
\hline
Racine carrée & $\ln(\sqrt{a})=\dfrac{1}{2}\ln a$ \\
\hline
Valeurs repères & $\ln 1=0$ ; \quad $\ln e=1$ \\
\hline
Équivalence clé & $\ln A=\ln B \iff A=B$ \quad ($A,B>0$) \\
\hline
\end{tabularx}
\end{center}
\end{aretenir}

\begin{savaistu}[Les tables de logarithmes : trois siècles de calculs]
Avant l'invention des calculatrices électroniques (années 1970), multiplier de grands nombres à la main était long et source d'erreurs. Grâce à la formule $\ln(ab)=\ln a+\ln b$, on pouvait remplacer une \emph{multiplication} par une simple \emph{addition} : on cherchait $\ln a$ et $\ln b$ dans une \textbf{table de logarithmes}, on les additionnait, puis on relisait le résultat dans la table en sens inverse. Inventés par John Napier au début du XVII\textsuperscript{e} siècle, les logarithmes ont servi aux astronomes, aux navigateurs et aux ingénieurs pendant plus de trois siècles. La \textbf{règle à calcul}, outil emblématique des techniciens jusqu'aux années 1970, fonctionne exactement sur ce principe.
\end{savaistu}

\courssub{Simplification d'expressions : écrire sous la forme d'un seul logarithme}

\begin{methode}[Regrouper plusieurs logarithmes en un seul]
Pour simplifier une expression contenant plusieurs logarithmes :
\begin{enumerate}
\item On \textbf{élimine les coefficients} devant les logarithmes en les faisant \og monter \fg{} en exposant : $k\ln a=\ln(a^{k})$.
\item On \textbf{regroupe} : les sommes deviennent des produits, les différences deviennent des quotients :
\[ \ln a+\ln b=\ln(ab), \qquad \ln a-\ln b=\ln\!\left(\frac{a}{b}\right). \]
\item On \textbf{simplifie} l'écriture obtenue, et si possible on factorise à nouveau (par exemple $\ln 9=2\ln 3$).
\end{enumerate}
\end{methode}

\begin{exoresolu}[Simplifier une somme de logarithmes]
\textbf{Énoncé.} Écrire sous la forme d'un seul logarithme, puis simplifier : $E=\ln 12+\ln 3-\ln 4$.
\tcblower
\textbf{Solution détaillée.}
\begin{align*}
E &= \ln 12+\ln 3-\ln 4 \\
&= \ln(12\times 3)-\ln 4 && \text{(la somme devient un produit)} \\
&= \ln\!\left(\frac{12\times 3}{4}\right) && \text{(la différence devient un quotient)} \\
&= \ln\!\left(\frac{36}{4}\right)=\ln 9.
\end{align*}
Comme $9=3^{2}$, on peut encore écrire : $E=\ln(3^{2})=2\ln 3$.
\textbf{Conclusion :} $E=\ln 9=2\ln 3$.
\end{exoresolu}

\begin{exoresolu}[Avec des coefficients]
\textbf{Énoncé.} Simplifier $F=2\ln 5+\dfrac{1}{2}\ln 16-\ln 10$.
\tcblower
\textbf{Solution détaillée.} On commence par faire monter les coefficients en exposants :
\[ 2\ln 5=\ln(5^{2})=\ln 25, \qquad \frac{1}{2}\ln 16=\ln(16^{1/2})=\ln\sqrt{16}=\ln 4. \]
Puis on regroupe :
\[ F=\ln 25+\ln 4-\ln 10=\ln\!\left(\frac{25\times 4}{10}\right)=\ln 10. \]
\textbf{Conclusion :} $F=\ln 10$.
\end{exoresolu}

\courssub{Lien avec les suites géométriques}

Voici une application remarquable, très classique au Baccalauréat technique.

\begin{propriete}[Du géométrique vers l'arithmétique]
Soit $(u_{n})$ une suite géométrique de premier terme $u_{0}>0$ et de raison $q>0$, avec $u_{n}>0$ pour tout $n$. Alors la suite $(v_{n})$ définie par $v_{n}=\ln(u_{n})$ est une suite \textbf{arithmétique} de raison $r=\ln q$ et de premier terme $v_{0}=\ln u_{0}$.
\end{propriete}

\textbf{Démonstration.} Comme $(u_{n})$ est géométrique de raison $q$, on a $u_{n+1}=q\,u_{n}$. Alors :
\[ v_{n+1}=\ln(u_{n+1})=\ln(q\,u_{n})=\ln q+\ln(u_{n})=v_{n}+\ln q. \]
La suite $(v_{n})$ passe donc d'un terme au suivant en ajoutant la constante $\ln q$ : c'est une suite arithmétique de raison $\ln q$. $\blacksquare$

\begin{popfigure}
\begin{center}
\begin{tikzpicture}[scale=1, every node/.style={font=\small}]
% Suite géométrique
\node[font=\bfseries, popblue] at (-1.2,1.6) {$(u_n)$};
\foreach \x/\val in {0/{2}, 2.2/{6}, 4.4/{18}, 6.6/{54}} {
  \node[draw=popblue, fill=popblueL, rounded corners=2pt, minimum width=1.1cm] at (\x,1.6) {$\val$};
}
\foreach \x in {1.1, 3.3, 5.5} {
  \node[popdark] at (\x,1.95) {$\times 3$};
  \draw[->, popblue, thick] (\x-0.55,1.6) -- (\x+0.55,1.6);
}
% Flèches ln
\foreach \x in {0, 2.2, 4.4, 6.6} {
  \draw[->, poporange, thick] (\x,1.15) -- (\x,0.35);
}
\node[poporange, font=\bfseries] at (7.6,0.75) {$\ln$};
% Suite arithmétique
\node[font=\bfseries, popgreen!60!black] at (-1.2,-0.1) {$(v_n)$};
\foreach \x/\val in {0/{$\ln 2$}, 2.2/{$\ln 6$}, 4.4/{$\ln 18$}, 6.6/{$\ln 54$}} {
  \node[draw=popgreen!60!black, fill=popgreenL, rounded corners=2pt, minimum width=1.1cm] at (\x,-0.1) {\val};
}
\foreach \x in {1.1, 3.3, 5.5} {
  \node[popdark] at (\x,0.28) {$+\ln 3$};
  \draw[->, popgreen!60!black, thick] (\x-0.62,-0.1) -- (\x+0.62,-0.1);
}
\end{tikzpicture}
\end{center}
\legende{Le logarithme transforme une suite géométrique (raison $q=3$, on multiplie) en une suite arithmétique (raison $\ln 3$, on ajoute).}
\end{popfigure}

\begin{exemplebox}[Application : un capital qui croît]
Un technicien place un capital qui augmente de $5\,\%$ par an. Si $u_{n}$ est le capital (en milliers de FCFA) après $n$ années, alors $u_{n+1}=1{,}05\,u_{n}$ : la suite est géométrique de raison $q=1{,}05$. La suite $v_{n}=\ln(u_{n})$ est arithmétique de raison $\ln(1{,}05)\approx 0{,}0488$. C'est ce passage à une suite arithmétique qui permet de répondre facilement à la question : \og au bout de combien d'années le capital aura-t-il doublé ? \fg{} (voir plus bas).
\end{exemplebox}

\courssub{Résolution d'équations contenant des logarithmes}

\begin{methode}[Résoudre une équation avec plusieurs logarithmes]
\begin{enumerate}
\item \textbf{Conditions d'existence} : on exige que chaque quantité sous un $\ln$ soit \emph{strictement positive}. On résout ce système d'inéquations \emph{avant} tout calcul.
\item \textbf{Regroupement} : on transforme les sommes et différences de logarithmes en un seul logarithme.
\item \textbf{Élimination du $\ln$} : on utilise l'équivalence
\[ \ln A=\ln B \iff A=B \qquad (A>0,\ B>0), \]
ou bien $\ln A=\ln c \iff A=c$.
\item \textbf{Vérification} : on ne garde que les solutions qui respectent les conditions d'existence de l'étape 1.
\end{enumerate}
\end{methode}

\begin{exoresolu}[Résoudre $\ln(x)+\ln(x-1)=\ln 6$]
\textbf{Énoncé.} Résoudre dans $\mathbb{R}$ l'équation $\ln(x)+\ln(x-1)=\ln 6$.
\tcblower
\textbf{Solution détaillée.}
\textbf{Étape 1 : conditions d'existence.} Il faut simultanément :
\[ x>0 \qquad \text{et} \qquad x-1>0 \iff x>1. \]
On cherchera donc des solutions dans $]1\,;\,+\infty[$.

\textbf{Étape 2 : regroupement.} La somme devient un produit :
\[ \ln\!\big(x(x-1)\big)=\ln 6. \]

\textbf{Étape 3 : élimination du $\ln$.} Par l'équivalence $\ln A=\ln B\iff A=B$ :
\[ x(x-1)=6 \iff x^{2}-x-6=0. \]
Discriminant : $\Delta=(-1)^{2}-4\times 1\times(-6)=1+24=25$, donc $\sqrt{\Delta}=5$. Les solutions de l'équation du second degré sont :
\[ x_{1}=\frac{1-5}{2}=-2 \qquad \text{et} \qquad x_{2}=\frac{1+5}{2}=3. \]

\textbf{Étape 4 : vérification.} La condition est $x>1$. La valeur $x_{1}=-2$ est \textbf{rejetée} (elle ne vérifie pas les conditions d'existence). La valeur $x_{2}=3$ convient : $\ln 3+\ln 2=\ln 6$. $\checkmark$

\textbf{Conclusion :} l'équation admet une unique solution : $S=\{3\}$.
\end{exoresolu}

\begin{attention}[Ne jamais oublier les conditions d'existence]
Dans l'exercice précédent, sans la vérification finale, on aurait donné \emph{deux} solutions, dont une fausse : $\ln(-2)$ n'existe pas ! Toute valeur trouvée doit être confrontée aux conditions posées au départ. C'est le premier point contrôlé par le correcteur au Baccalauréat.
\end{attention}

\courssub{Application concrète : le temps de doublement}

\textbf{Situation.} De nombreuses grandeurs techniques et économiques évoluent de façon multiplicative : un capital placé à intérêts composés, la population d'une ville, la charge d'un condensateur, le nombre de bactéries dans une culture. Si la quantité est multipliée par $q>1$ à chaque période, sa valeur au temps $t$ est :
\[ u(t)=a\,q^{t}, \]
où $a$ est la valeur initiale. Une question naturelle se pose : \textbf{au bout de combien de temps la quantité aura-t-elle doublé ?}

\begin{exoresolu}[Temps de doublement]
\textbf{Énoncé.} Une quantité évolue selon $u(t)=a\,q^{t}$ avec $q>1$. Montrer que le temps de doublement est $t=\dfrac{\ln 2}{\ln q}$. Application : un capital placé à $4\,\%$ l'an ($q=1{,}04$) ; au bout de combien d'années a-t-il doublé ?
\tcblower
\textbf{Solution détaillée.} On cherche $t$ tel que $u(t)=2a$ :
\begin{align*}
a\,q^{t} &= 2a \\
q^{t} &= 2 && \text{(on divise par } a>0\text{)} \\
\ln(q^{t}) &= \ln 2 && \text{(on applique } \ln\text{)} \\
t\,\ln q &= \ln 2 && \text{(propriété des puissances)} \\
t &= \frac{\ln 2}{\ln q}.
\end{align*}
Remarquons que le résultat \textbf{ne dépend pas de $a$} : peu importe la somme de départ, le temps de doublement est le même.

\textbf{Application numérique.} Avec $q=1{,}04$ :
\[ t=\frac{\ln 2}{\ln(1{,}04)}\approx\frac{0{,}693}{0{,}0392}\approx 17{,}7. \]
\textbf{Conclusion :} le capital double en environ $18$ ans. C'est ce type de calcul qu'utilisent les banquiers et les gestionnaires pour comparer des placements.
\end{exoresolu}

\begin{aretenir}[Ce qu'il faut retenir de cette section]
\begin{itemize}
\item La propriété fondamentale $\ln(ab)=\ln a+\ln b$ se démontre en posant $a=e^{A}$, $b=e^{B}$ et en utilisant $e^{A}e^{B}=e^{A+B}$.
\item Toutes les autres formules (inverse, quotient, puissance, racine) en découlent.
\item $\ln(a+b)$ ne se simplifie \emph{jamais} : c'est le piège numéro 1.
\item Pour résoudre une équation avec des logarithmes : conditions d'existence, regroupement, équivalence $\ln A=\ln B\iff A=B$, vérification.
\item Le logarithme transforme une suite géométrique en suite arithmétique : c'est l'outil-clé pour les problèmes de doublement, avec $t=\dfrac{\ln 2}{\ln q}$.
\end{itemize}
\end{aretenir}

\begin{exercice}[Pour s'entraîner]
\begin{enumerate}
\item Simplifier : $A=\ln 18+\ln 2-\ln 4$ ; \quad $B=3\ln 2+2\ln 3-\ln 12$.
\item Résoudre dans $\mathbb{R}$ : $\ln(x+2)+\ln(x)=\ln 3$.
\item Une population de bactéries double toutes les 20 minutes. Sachant qu'elle est multipliée par $q$ chaque minute, déterminer $q$ (on posera $q^{20}=2$).
\item Soit $(u_{n})$ la suite géométrique de premier terme $u_{0}=5$ et de raison $q=1{,}1$. Montrer que $v_{n}=\ln(u_{n})$ est arithmétique et préciser sa raison.
\end{enumerate}
\end{exercice}

\begin{corrige}[Exercice n°1]
\textbf{1.} $A=\ln\!\left(\dfrac{18\times 2}{4}\right)=\ln 9=2\ln 3$. \quad $B=\ln 8+\ln 9-\ln 12=\ln\!\left(\dfrac{72}{12}\right)=\ln 6$.

\textbf{2.} Conditions : $x+2>0$ et $x>0$, donc $x>0$. L'équation devient $\ln\!\big(x(x+2)\big)=\ln 3$, soit $x^{2}+2x-3=0$. $\Delta=4+12=16$, $x_{1}=\dfrac{-2-4}{2}=-3$ (rejetée), $x_{2}=\dfrac{-2+4}{2}=1$. $S=\{1\}$.

\textbf{3.} $q^{20}=2$ donne $20\ln q=\ln 2$, donc $\ln q=\dfrac{\ln 2}{20}\approx 0{,}0347$, puis $q=e^{0{,}0347}\approx 1{,}0353$. La population augmente d'environ $3{,}5\,\%$ par minute.

\textbf{4.} $v_{n+1}=\ln(u_{n+1})=\ln(1{,}1\,u_{n})=\ln(1{,}1)+v_{n}$ : $(v_{n})$ est arithmétique de raison $r=\ln(1{,}1)\approx 0{,}0953$ et de premier terme $v_{0}=\ln 5$.
\end{corrige}

\coursec{Limites faisant intervenir ln}

Après avoir maîtrisé les propriétés algébriques du logarithme népérien, nous abordons maintenant son comportement aux frontières de son ensemble de définition et sa croissance comparée aux fonctions puissances. Ces limites sont les outils indispensables pour étudier les fonctions composées avec $\ln$, calculer des asymptotes et résoudre des indéterminations classiques.

\courssub{Limites aux bornes de l'intervalle de définition}

Rappelons que la fonction $\ln$ est définie et strictement croissante sur $]0\,;+\infty[$. Son comportement aux bornes de cet intervalle est fondamental.

\begin{propriete}[{Limites de la fonction $\ln$ aux bornes de $]0\,;+\infty[$}]
La fonction logarithme népérien admet les limites suivantes :
\[
\lim_{x\to 0^+} \ln x = -\infty \qquad\text{et}\qquad \lim_{x\to +\infty} \ln x = +\infty.
\]
Ces limites sont \textbf{admises} au Probatoire. Elles traduisent graphiquement :
\begin{itemize}
    \item l'existence d'une \textbf{asymptote verticale} d'équation $x=0$ (l'axe des ordonnées) pour la courbe $\mathcal{C}_{\ln}$ ;
    \item une croissance sans limite, mais \emph{très lente}, vers $+\infty$ lorsque $x\to +\infty$.
\end{itemize}
\end{propriete}

\begin{popfigure}
\begin{tikzpicture}[scale=0.85]
\begin{axis}[
    axis lines=middle,
    xlabel=$x$, ylabel=$y$,
    xmin=-0.5, xmax=5.5,
    ymin=-3.5, ymax=3.5,
    xtick={0,1,2,3,4,5},
    ytick={-3,-2,-1,0,1,2,3},
    xticklabels={$0$,$1$,$2$,$3$,$4$,$5$},
    yticklabels={$-3$,$-2$,$-1$,$0$,$1$,$2$,$3$},
    clip=false,
    domain=0.01:5.5,
    samples=300,
    width=\linewidth,
    height=0.6\linewidth,
    axis line style={-Stealth, thick},
    tick style={thick},
    xlabel style={anchor=west},
    ylabel style={anchor=south},
]
% Asymptote verticale x=0
\draw[popline, dashed, thick] (axis cs:0,-3.5) -- (axis cs:0,3.5) node[above left, popmuted] {$x=0$};
% Courbe ln x
\addplot[popcurve, very thick, domain=0.01:5.5] {ln(x)};
% Point (1,0)
\addplot[only marks, mark=*, mark size=2.5, poporange] coordinates {(1,0)};
\node[above right, poporange] at (axis cs:1,0) {$(1,0)$};
% Flèches indicatives
\draw[popblue, -Stealth, thick] (axis cs:0.1,-2.3) -- (axis cs:0.02,-3);
\node[popblue, font=\small] at (axis cs:0.5,-3.2) {$-\infty$};
\draw[popblue, -Stealth, thick] (axis cs:4.5,1.5) -- (axis cs:5.2,1.7);
\node[popblue, font=\small] at (axis cs:5.5,1.8) {$+\infty$};
\end{axis}
\end{tikzpicture}
\legende{Courbe de $y=\ln x$ : asymptote verticale $x=0$ et croissance lente vers $+\infty$.}
\end{popfigure}

\begin{exemplebox}[Application directe]
Calculer les limites suivantes :
\begin{enumerate}
    \item $\displaystyle\lim_{x\to 0^+} \ln(x^2)$
    \item $\displaystyle\lim_{x\to +\infty} \ln(\sqrt{x})$
    \item $\displaystyle\lim_{x\to 1^-} \ln(1-x)$
\end{enumerate}
\textbf{Solution.}
\begin{enumerate}
    \item Pour $x\to 0^+$, $x^2\to 0^+$. Par composition : $\ln(x^2)\to -\infty$.
    \item Pour $x\to +\infty$, $\sqrt{x}\to +\infty$. Donc $\ln(\sqrt{x})\to +\infty$.
    \item Pour $x\to 1^-$, $1-x\to 0^+$. Donc $\ln(1-x)\to -\infty$.
\end{enumerate}
\end{exemplebox}

\courssub{Croissances comparées : limites de référence}

La fonction $\ln x$ tend vers $+\infty$ quand $x\to +\infty$, mais \emph{plus lentement que toute fonction puissance $x^\alpha$ avec $\alpha>0$}. C'est le cœur du théorème des croissances comparées.

\begin{propriete}[Théorème des croissances comparées (admis)]
\begin{enumerate}
    \item $\displaystyle\lim_{x\to +\infty} \frac{\ln x}{x} = 0$.
    \item Plus généralement, pour tout entier $n\ge 1$ : $\displaystyle\lim_{x\to +\infty} \frac{\ln x}{x^n} = 0$.
    \item En $0^+$ : $\displaystyle\lim_{x\to 0^+} x\ln x = 0$.
\end{enumerate}
\textbf{Interprétation :} « Toute puissance positive de $x$, aussi petite soit-elle, l'emporte finalement sur $\ln x$. »
\end{propriete}

\begin{attention}[Piège fréquent : la « vitesse » de $\ln$]
Ne croyez \textbf{pas} que $\ln x$ tend vers $+\infty$ « vite ». Au contraire, sa croissance est \emph{extrêmement lente} :
\[
\ln 10 \approx 2,3,\quad \ln 100 \approx 4,6,\quad \ln 1000 \approx 6,9,\quad \ln 10^6 \approx 13,8.
\]
Pour atteindre $20$, il faut $x \approx e^{20} \approx 4,85\times 10^8$. C'est pourquoi $\frac{\ln x}{x}\to 0$ : le numérateur rampe pendant que le dénominateur file.
\end{attention}

\begin{popfigure}
\begin{tikzpicture}[scale=0.85]
\begin{axis}[
    axis lines=middle,
    xlabel=$x$, ylabel=$y$,
    xmin=0, xmax=50,
    ymin=0, ymax=50,
    xtick={0,10,20,30,40,50},
    ytick={0,10,20,30,40,50},
    clip=false,
    domain=0.1:50,
    samples=300,
    width=\linewidth,
    height=0.65\linewidth,
    axis line style={-Stealth, thick},
    tick style={thick},
    legend style={at={(0.02,0.98)}, anchor=north west, fill=white, fill opacity=0.9, draw=popline},
]
\addplot[popblue, thick] {ln(x)};
\addlegendentry{$y=\ln x$}
\addplot[poporange, thick] {sqrt(max(0,x))};
\addlegendentry{$y=\sqrt{x}$}
\addplot[popgreen, thick] {x};
\addlegendentry{$y=x$}
% Marqueurs à x=10, 50
\addplot[only marks, mark=*, mark size=2, popblue] coordinates {(10,2.3026) (50,3.912)};
\addplot[only marks, mark=*, mark size=2, poporange] coordinates {(10,3.162) (50,7.071)};
\addplot[only marks, mark=*, mark size=2, popgreen] coordinates {(10,10) (50,50)};
\end{axis}
\end{tikzpicture}
\legende{Comparaison graphique sur $[0;50]$ : $\ln x$ (bleu) est largement dominée par $\sqrt{x}$ (orange) et $x$ (vert).}
\end{popfigure}

\begin{exemplebox}[Limite de $x - \ln x$ en $+\infty$]
On cherche $\displaystyle\lim_{x\to +\infty} (x - \ln x)$.
\begin{align*}
x - \ln x &= x\left(1 - \frac{\ln x}{x}\right) \\
\lim_{x\to +\infty} \frac{\ln x}{x} &= 0 \quad\text{(croissance comparée)} \\
\Rightarrow\quad 1 - \frac{\ln x}{x} &\to 1 \\
\text{Donc } x\left(1 - \frac{\ln x}{x}\right) &\to +\infty \times 1 = +\infty.
\end{align*}
\emph{Conclusion :} $\displaystyle\lim_{x\to +\infty} (x - \ln x) = +\infty$.
\end{exemplebox}

\begin{experience}[Démonstration guidée de $\lim_{x\to 0^+} x\ln x = 0$]
On pose $X = \frac{1}{x}$. Alors $x\to 0^+ \iff X\to +\infty$.
\[
x\ln x = \frac{1}{X}\ln\left(\frac{1}{X}\right) = -\frac{\ln X}{X}.
\]
Or $\displaystyle\lim_{X\to +\infty} \frac{\ln X}{X} = 0$ (théorème des croissances comparées).
Donc $\displaystyle\lim_{x\to 0^+} x\ln x = -0 = 0$.
\end{experience}

\courssub{Taux d'accroissement et limite remarquable en $1$}

Le comportement de $\ln$ au voisinage de $1$ est lié à sa dérivabilité en ce point.

\begin{propriete}[Limite du taux d'accroissement de $\ln$ en $1$]
\[
\lim_{h\to 0} \frac{\ln(1+h)}{h} = 1.
\]
Cette limite est \textbf{admise} (elle exprime que $(\ln)'(1)=1$). Elle permet de traiter les formes indéterminées $\frac{0}{0}$ faisant intervenir $\ln(1+u)$ avec $u\to 0$.
\end{propriete}

\begin{exemplebox}[Utilisation de la limite remarquable]
Calculer $\displaystyle\lim_{x\to 0} \frac{\ln(1+3x)}{x}$.
\[
\frac{\ln(1+3x)}{x} = 3 \times \frac{\ln(1+3x)}{3x}.
\]
Posons $h = 3x$. Quand $x\to 0$, $h\to 0$.
\[
\lim_{x\to 0} \frac{\ln(1+3x)}{x} = 3 \times \lim_{h\to 0} \frac{\ln(1+h)}{h} = 3 \times 1 = 3.
\]
\end{exemplebox}

\courssub{Techniques de levée d'indétermination}

Face à une forme indéterminée mettant en jeu $\ln$, deux stratégies principales s'offrent à nous.

\begin{methode}[Résolution d'indéterminations avec $\ln$]
\textbf{Étape 1 : Identifier la forme.} Les cas fréquents sont :
$\infty - \infty$ (ex. $x - \ln x$), $\frac{\infty}{\infty}$ (ex. $\frac{\ln x}{x}$), $0\times\infty$ (ex. $x\ln x$), $\frac{0}{0}$ (ex. $\frac{\ln(1+x)}{x}$).

\textbf{Étape 2 : Choisir la technique.}
\begin{enumerate}
    \item \textbf{Factorisation par le terme dominant} (pour $\infty-\infty$ ou $\frac{\infty}{\infty}$) : factoriser la partie qui tend vers $\infty$ le plus « fort » (généralement la puissance de $x$).
    \item \textbf{Changement de variable} :
    \begin{itemize}
        \item $X = \ln x$ (ramène $x\to 0^+$ à $X\to -\infty$ et $x\to +\infty$ à $X\to +\infty$) ;
        \item $X = \frac{1}{x}$ (échange $0^+$ et $+\infty$, utile pour $x\ln x$).
    \end{itemize}
    \item \textbf{Utilisation des limites de référence} : $\frac{\ln x}{x}\to 0$, $x\ln x\to 0$, $\frac{\ln(1+h)}{h}\to 1$.
    \item \textbf{Propriétés algébriques} : transformer $\ln A - \ln B$ en $\ln\frac{A}{B}$ pour faire apparaître $1+u$.
\end{enumerate}

\textbf{Étape 3 : Conclure rigoureusement} en citant les théorèmes utilisés.
\end{methode}

\begin{exoresolu}[Exercices résolus : levée d'indéterminations]
\textbf{Exercice 1.} $\displaystyle\lim_{x\to +\infty} \frac{\ln(x^2+1)}{x}$.
\\
\textbf{Solution.}
Pour $x\to +\infty$, $x^2+1 \sim x^2$, donc $\ln(x^2+1) \sim \ln(x^2) = 2\ln x$.
\[
\frac{\ln(x^2+1)}{x} \sim \frac{2\ln x}{x} = 2\frac{\ln x}{x} \to 0.
\]
\emph{Réponse :} $0$.

\tcblower

\textbf{Exercice 2.} $\displaystyle\lim_{x\to 0^+} x^2 \ln x$.
\\
\textbf{Solution.}
Forme $0 \times (-\infty)$. Posons $X = \frac{1}{x}$ ($X\to +\infty$).
\[
x^2 \ln x = \frac{1}{X^2} \ln\left(\frac{1}{X}\right) = -\frac{\ln X}{X^2}.
\]
Or $\displaystyle\lim_{X\to +\infty} \frac{\ln X}{X^2} = 0$ (croissance comparée, $n=2$).
\emph{Réponse :} $0$.

\tcblower

\textbf{Exercice 3.} $\displaystyle\lim_{x\to +\infty} \bigl(\ln(x+1) - \ln x\bigr)$.
\\
\textbf{Solution.}
Forme $\infty - \infty$. Utilisons les propriétés algébriques :
\[
\ln(x+1) - \ln x = \ln\left(\frac{x+1}{x}\right) = \ln\left(1+\frac{1}{x}\right).
\]
Quand $x\to +\infty$, $\frac{1}{x}\to 0$. Par la limite remarquable :
\[
\ln\left(1+\frac{1}{x}\right) \sim \frac{1}{x} \to 0.
\]
\emph{Réponse :} $0$.
\end{exoresolu}

\courssub{Limites de fonctions composées avec $\ln$}

La règle de composition des limites s'applique naturellement, sous réserve de respecter l'ensemble de définition de $\ln$ ($]0,+\infty[$).

\begin{propriete}[Limite de $\ln(u(x))$]
Soit $u$ une fonction définie sur un intervalle $I$ et $a \in \overline{I}$ (borne finie ou infinie).
\begin{enumerate}
    \item Si $\displaystyle\lim_{x\to a} u(x) = \ell \in ]0,+\infty[$, alors $\displaystyle\lim_{x\to a} \ln(u(x)) = \ln \ell$.
    \item Si $\displaystyle\lim_{x\to a} u(x) = 0^+$, alors $\displaystyle\lim_{x\to a} \ln(u(x)) = -\infty$.
    \item Si $\displaystyle\lim_{x\to a} u(x) = +\infty$, alors $\displaystyle\lim_{x\to a} \ln(u(x)) = +\infty$.
    \item Si $\displaystyle\lim_{x\to a} u(x) \le 0$, la limite de $\ln(u(x))$ \textbf{n'existe pas} (fonction non définie au voisinage de $a$).
\end{enumerate}
\end{propriete}

\begin{exemplebox}[Limites composées]
\begin{enumerate}
    \item $\displaystyle\lim_{x\to 0} \ln(2+x) = \ln 2$ (car $2+x\to 2>0$).
    \item $\displaystyle\lim_{x\to 0^+} \ln(\sin x) = -\infty$ (car $\sin x \to 0^+$).
    \item $\displaystyle\lim_{x\to +\infty} \ln(e^x+1) = +\infty$ (car $e^x+1\to +\infty$).
    \item $\displaystyle\lim_{x\to 0} \ln(x^2) = -\infty$ (car $x^2\to 0^+$).
    \item $\displaystyle\lim_{x\to 1} \ln(x-1)$ \textbf{n'existe pas} : pour $x<1$, $x-1<0$ et $\ln$ n'est pas défini.
\end{enumerate}
\end{exemplebox}

\begin{aretenir}[Les cinq limites de référence à connaître par cœur]
\begin{enumerate}
    \item $\displaystyle\lim_{x\to 0^+} \ln x = -\infty$
    \item $\displaystyle\lim_{x\to +\infty} \ln x = +\infty$
    \item $\displaystyle\lim_{x\to +\infty} \frac{\ln x}{x} = 0$ \quad (et $\frac{\ln x}{x^n}=0$ pour $n\in\mathbb{N}^*$)
    \item $\displaystyle\lim_{x\to 0^+} x\ln x = 0$
    \item $\displaystyle\lim_{h\to 0} \frac{\ln(1+h)}{h} = 1$
\end{enumerate}
\end{aretenir}

\courssub{Applications concrètes : modélisation technique}

Les limites de $\ln$ interviennent dans de nombreux modèles techniques : décharge d'un condensateur, refroidissement newtonien, loi de désintégration radioactive, calcul de temps de réponse, etc.

\begin{exoresolu}[Temps de décharge d'un condensateur — Électricité]
Un condensateur de capacité $C$ se décharge à travers une résistance $R$. La tension $u(t)$ aux bornes du condensateur suit la loi :
\[
u(t) = U_0 e^{-t/(RC)},\quad t\ge 0,
\]
où $U_0$ est la tension initiale. On cherche le temps $t_{1\%}$ nécessaire pour que la tension descende à $1\%$ de $U_0$.

\textbf{Solution.}
On résout $u(t) = 0,01\,U_0$ :
\[
U_0 e^{-t/(RC)} = 0,01\,U_0 \;\Longrightarrow\; e^{-t/(RC)} = 0,01.
\]
On applique $\ln$ des deux côtés :
\[
-\frac{t}{RC} = \ln(0,01) = \ln(10^{-2}) = -2\ln 10.
\]
Donc $t_{1\%} = 2RC\ln 10 \approx 4,605\,RC$.

\textbf{Interprétation par les limites :}
\begin{itemize}
    \item $\displaystyle\lim_{t\to +\infty} u(t) = 0$ (la tension tend vers $0$ mais ne l'atteint jamais en temps fini) ;
    \item $\displaystyle\lim_{t\to 0^+} u(t) = U_0$ (continuité à l'origine).
\end{itemize}
Ce modèle illustre la croissance lente de $\ln$ : pour diviser la tension par $10$, il faut ajouter $RC\ln 10 \approx 2,3\,RC$ au temps ; pour la diviser par $100$, il faut $2\times 2,3\,RC$, etc. (propriété $\ln(ab)=\ln a+\ln b$).
\end{exoresolu}

\begin{savaistu}[Échelle logarithmique et décibels — Acoustique / Électronique]
L'oreille humaine perçoit les intensités sonores sur une échelle \emph{logarithmique}. Le niveau sonore en décibels (dB) est défini par :
\[
L = 10 \log_{10}\left(\frac{I}{I_0}\right) = \frac{10}{\ln 10} \ln\left(\frac{I}{I_0}\right),
\]
où $I$ est l'intensité et $I_0=10^{-12}\,\text{W/m}^2$ le seuil d'audibilité.
\begin{itemize}
    \item Une multiplication de l'intensité par $10$ ajoute $10\,\text{dB}$.
    \item La limite $\lim_{I\to 0^+} \ln I = -\infty$ signifie qu'un son infiniment faible correspond à $-\infty\,\text{dB}$ (inaudible).
    \item La limite $\lim_{I\to +\infty} \ln I = +\infty$ n'a pas de sens physique réel (seuil de douleur vers $120\text{--}130\,\text{dB}$).
\end{itemize}
C'est une application directe des limites de $\ln$ aux bornes de son domaine.
\end{savaistu}

\begin{exercice}[Entraînement — Limites avec $\ln$]
Calculer les limites suivantes (préciser $+\infty$, $-\infty$, valeur finie ou « n'existe pas ») :
\begin{enumerate}
    \item $\displaystyle\lim_{x\to +\infty} \frac{\ln x}{\sqrt{x}}$
    \item $\displaystyle\lim_{x\to 0^+} \sqrt{x}\ln x$
    \item $\displaystyle\lim_{x\to +\infty} \bigl(\ln(x^2+3x) - 2\ln x\bigr)$
    \item $\displaystyle\lim_{x\to 0} \frac{\ln(1+5x)}{x}$
    \item $\displaystyle\lim_{x\to 1} \ln(x-1)$
    \item $\displaystyle\lim_{x\to +\infty} \frac{(\ln x)^2}{x}$
\end{enumerate}
\end{exercice}

\begin{corrige}[Exercice — Limites avec $\ln$]
\begin{enumerate}
    \item $\frac{\ln x}{\sqrt{x}} = \frac{\ln x}{x^{1/2}} \to 0$ (croissance comparée, valable pour toute puissance $\alpha>0$). \textbf{Réponse : $0$.}
    \item Posons $X=1/x$ ($X\to +\infty$) : $\sqrt{x}\ln x = \frac{1}{\sqrt{X}}\ln(1/X) = -\frac{\ln X}{\sqrt{X}} \to 0$. \textbf{Réponse : $0$.}
    \item $\ln(x^2+3x) - 2\ln x = \ln\left(\frac{x^2+3x}{x^2}\right) = \ln\left(1+\frac{3}{x}\right) \sim \frac{3}{x} \to 0$. \textbf{Réponse : $0$.}
    \item $\frac{\ln(1+5x)}{x} = 5\frac{\ln(1+5x)}{5x} \to 5\times 1 = 5$. \textbf{Réponse : $5$.}
    \item Pour $x\to 1^-$, $x-1\to 0^-$ : $\ln$ non défini. Pour $x\to 1^+$, $x-1\to 0^+$ : $\ln(x-1)\to -\infty$. La limite bilatère \textbf{n'existe pas}. \textbf{Réponse : n'existe pas.}
    \item $\frac{(\ln x)^2}{x} = \left(\frac{\ln x}{\sqrt{x}}\right)^2 \to 0^2 = 0$. \textbf{Réponse : $0$.}
\end{enumerate}
\end{corrige}

\coursec{Dérivation des fonctions définies à l'aide de $\ln$}

Après avoir étudié les limites des fonctions faisant intervenir le \cle{logarithme népérien}, nous abordons maintenant leur \cle{dérivation}. Cette section est fondamentale : elle nous donnera les outils pour étudier les variations de ces fonctions (section 5), tracer leurs courbes et résoudre des problèmes d'\cle{optimisation} concrets (section 6). Rappelons que $\ln$ est la bijection réciproque de l'exponentielle : cette relation va nous permettre de trouver sa dérivée de façon élégante.

\courssub{Dérivabilité de la fonction logarithme népérien}

\begin{experience}[Démonstration guidée : dérivée de $\ln$]
On sait que $\ln : ]0;+\infty[ \to \mathbb{R}$ est la bijection réciproque de $\exp : \mathbb{R} \to ]0;+\infty[$. Par définition, pour tout $x>0$ :
\[
e^{\ln x} = x.
\]
Les deux fonctions sont dérivables (l'exponentielle sur $\mathbb{R}$, le logarithme sur $]0;+\infty[$ par le théorème de dérivation des bijections réciproques). Dérivons l'identité membre à membre par rapport à $x$ :
\[
\frac{d}{dx}\bigl(e^{\ln x}\bigr) = \frac{d}{dx}(x).
\]
À gauche, on utilise la \cle{règle de la chaîne} : $(e^u)' = u' e^u$ avec $u(x)=\ln x$. On obtient :
\[
(\ln)'(x) \cdot e^{\ln x} = 1.
\]
Mais $e^{\ln x} = x$, donc :
\[
(\ln)'(x) \cdot x = 1 \quad\Longrightarrow\quad (\ln)'(x) = \frac{1}{x}.
\]
\end{experience}

\begin{propriete}[Dérivée de la fonction logarithme népérien]
La fonction $\ln$ est dérivable sur $]0;+\infty[$ et pour tout $x>0$ :
\[
(\ln)'(x) = \frac{1}{x}.
\]
\end{propriete}

\begin{aretenir}[Conséquence immédiate sur les variations]
Puisque $\frac{1}{x} > 0$ pour tout $x \in ]0;+\infty[$, la fonction $\ln$ est \textbf{strictement croissante} sur $]0;+\infty[$. Cela confirme rigoureusement le sens de variation annoncé en Section 1.
\end{aretenir}

\begin{popfigure}
\begin{tikzpicture}
\begin{axis}[
    width=\linewidth, height=8cm,
    axis lines=middle,
    xlabel=$x$, ylabel=$y$,
    xmin=0, xmax=4,
    ymin=-1.5, ymax=2,
    xtick={1,2,3}, ytick={-1,1,2},
    domain=0.05:4, samples=200,
    clip=false,
    enlargelimits=0.05
]
\addplot[popcurve, thick, domain=0.05:4] {ln(x)};
\addplot[poporange, thick, dashed, domain=0:4] {x-1};
\addplot[only marks, mark=*, mark size=2.5, popdark] coordinates {(1,0)};
\node[above right, popdark] at (axis cs:1,0) {$(1;0)$};
\node[above left, poporange] at (axis cs:3.5,2.5) {$y=x-1$};
\node[above right, popblue] at (axis cs:3.8,1.35) {$y=\ln x$};
\end{axis}
\end{tikzpicture}
\legende{Courbe de $y=\ln x$ et sa tangente au point $(1;0)$ d'équation $y=x-1$. On voit que la courbe est entièrement située \emph{sous} sa tangente (sauf au point de contact).}
\end{popfigure}

\courssub{Dérivée de $\ln(u)$ : formule de dérivation composée}

Soit $u$ une fonction dérivable et \textbf{strictement positive} sur un intervalle $I$. La \cle{fonction composée} $x \mapsto \ln(u(x))$ est alors définie et dérivable sur $I$.

\begin{propriete}[Dérivée de $\ln(u)$]
Si $u$ est dérivable et $u(x) > 0$ sur $I$, alors :
\[
\bigl(\ln(u(x))\bigr)' = \frac{u'(x)}{u(x)} \quad \text{sur } I.
\]
\end{propriete}

\begin{methode}[Dériver une fonction de la forme $f(x)=\ln(u(x))$]
\begin{enumerate}
\item \textbf{Vérifier la condition de positivité :} déterminer l'ensemble $I$ où $u(x) > 0$. C'est l'ensemble de définition de $f$ et le domaine de dérivation.
\item \textbf{Calculer $u'(x)$ :} dériver la fonction intérieure $u$.
\item \textbf{Appliquer la formule :} $f'(x) = \dfrac{u'(x)}{u(x)}$ sur $I$.
\item \textbf{Simplifier si possible.}
\end{enumerate}
\end{methode}

\begin{exemplebox}[Exemple : $f(x) = \ln(x^2+1)$]
\begin{enumerate}
\item $u(x) = x^2+1$. Pour tout $x \in \mathbb{R}$, $x^2+1 \ge 1 > 0$. Donc $f$ est définie et dérivable sur $\mathbb{R}$.
\item $u'(x) = 2x$.
\item $f'(x) = \dfrac{2x}{x^2+1}$ sur $\mathbb{R}$.
\end{enumerate}
\end{exemplebox}

\begin{attention}[Erreurs fréquentes — Pièges à éviter]
\begin{itemize}
\item \textbf{Oublier le facteur $u'$ :} écrire $(\ln u)' = \frac{1}{u}$ au lieu de $\frac{u'}{u}$. C'est l'erreur n°1 ! La \cle{règle de la chaîne} est obligatoire.
\item \textbf{Dériver sans vérifier $u(x)>0$ :} par exemple, dériver $\ln(x-2)$ sur $\mathbb{R}$ alors qu'elle n'est définie que sur $]2;+\infty[$. La dérivée n'existe pas en dehors du domaine de définition.
\item \textbf{Confondre $\ln(u)$ et $\ln|u|$ :} la dérivée de $\ln|u|$ est aussi $\frac{u'}{u}$ mais sur l'ensemble où $u \neq 0$ (utile pour les primitives, voir Section 5).
\end{itemize}
\end{attention}

\courssub{Dérivées de produits et quotients faisant intervenir $\ln$}

On combine la formule $(\ln u)' = \frac{u'}{u}$ avec les règles usuelles : $(uv)' = u'v+uv'$ et $\left(\frac{u}{v}\right)' = \frac{u'v-uv'}{v^2}$.

\begin{exoresolu}[{Dérivée de $f(x) = x \ln x$ sur $]0;+\infty[$}]
$f$ est un produit : $u(x)=x$, $v(x)=\ln x$.
\[
f'(x) = u'(x)v(x) + u(x)v'(x) = 1 \cdot \ln x + x \cdot \frac{1}{x} = \ln x + 1.
\]
\end{exoresolu}

\begin{exoresolu}[{Dérivée de $g(x) = \dfrac{\ln x}{x}$ sur $]0;+\infty[$}]
$g$ est un quotient : $u(x)=\ln x$, $v(x)=x$.
\[
g'(x) = \frac{u'(x)v(x) - u(x)v'(x)}{v(x)^2} = \frac{\frac{1}{x}\cdot x - \ln x \cdot 1}{x^2} = \frac{1 - \ln x}{x^2}.
\]
\end{exoresolu}

\begin{aretenir}[Tableau récapitulatif des dérivées usuelles avec $\ln$]
\begin{center}
\begin{tabularx}{\linewidth}{|l|X|X|}\hline
\rowcolor{popblueL}
\textbf{Fonction $f(x)$} & \textbf{Domaine} & \textbf{Dérivée $f'(x)$} \\ \hline
$\ln x$ & $]0;+\infty[$ & $\dfrac{1}{x}$ \\ \hline
$\ln(u(x))$ & $\{x \mid u(x)>0\}$ & $\dfrac{u'(x)}{u(x)}$ \\ \hline
$x \ln x$ & $]0;+\infty[$ & $\ln x + 1$ \\ \hline
$\dfrac{\ln x}{x}$ & $]0;+\infty[$ & $\dfrac{1-\ln x}{x^2}$ \\ \hline
$\ln(x^2+1)$ & $\mathbb{R}$ & $\dfrac{2x}{x^2+1}$ \\ \hline
$\ln(ax+b)\;(a\neq0)$ & $]-\frac{b}{a};+\infty[$ si $a>0$ & $\dfrac{a}{ax+b}$ \\ \hline
\end{tabularx}
\end{center}
\end{aretenir}

\courssub{Tangente en 1 et inégalité fondamentale}

\begin{experience}[Tangente à la courbe de $\ln$ au point d'abscisse 1]
La courbe $\mathcal{C}$ de la fonction $\ln$ passe par $A(1;0)$ car $\ln 1 = 0$.
La pente de la tangente en $A$ est $f'(1) = \frac{1}{1} = 1$.
L'équation de la tangente $\mathcal{T}$ est :
\[
y - 0 = 1 \cdot (x - 1) \quad\Longleftrightarrow\quad y = x - 1.
\]
\end{experience}

\begin{propriete}[Inégalité fondamentale : $\ln x \le x-1$]
Pour tout $x > 0$, on a $\ln x \le x - 1$, avec égalité \textbf{si et seulement si} $x=1$.
\emph{Preuve :} la fonction $\ln$ est concave (sa dérivée $1/x$ est décroissante). Sa courbe se trouve donc sous sa tangente en tout point, en particulier en $1$.
\end{propriete}

\begin{savaistu}[L'inégalité $\ln(1+x) \le x$ : un outil universel]
En posant $x = 1+u$ ($u > -1$), l'inégalité devient $\ln(1+u) \le u$. Cette forme est omniprésente :
\begin{itemize}
\item \textbf{Probabilités / Statistique :} elle sert à majorer des logarithmes de vraisemblance, à démontrer l'inégalité de Jensen pour l'entropie, ou à borner l'erreur d'approximation de $\ln(1+u)$ par $u$ quand $u \approx 0$.
\item \textbf{Économie / Finance :} pour modéliser l'utilité logarithmique ou approximer des taux de croissance continus ($r \approx \ln(1+r)$ pour $r$ petit).
\item \textbf{Algorithmique :} analyse de complexité (ex. : borne sur la hauteur d'arbres binaires).
\end{itemize}
Au \cle{Probatoire}, cette inégalité peut être demandée en démonstration ou utilisée pour majorer/minorer une expression.
\end{savaistu}

\courssub{Applications : recherche d'extremums de fonctions contenant $\ln$}

La dérivation permet d'étudier les variations et de trouver les \cle{extremums} (maximum/minimum) de fonctions modélisant des situations techniques.

\begin{exoresolu}[Optimisation : rendement d'un procédé chimique]
Le rendement $R$ (en \%) d'une réaction en fonction de la température $x$ (en $^\circ$C) est modélisé par :
\[
R(x) = 100 \cdot \frac{\ln(x+1)}{x+1}, \quad x \in [0; 100].
\]
Trouver la température optimale (maximisant le rendement).
\vspace{0,5em}
\noindent\textbf{Solution.}
\begin{enumerate}
\item $R$ est dérivable sur $[0;100]$ comme quotient de fonctions dérivables (dénominateur $x+1 \neq 0$).
\item Posons $u(x) = \ln(x+1)$, $v(x) = x+1$. Alors $R(x) = 100 \frac{u}{v}$.
\[
u'(x) = \frac{1}{x+1}, \quad v'(x) = 1.
\]
\[
R'(x) = 100 \cdot \frac{u'v - uv'}{v^2} = 100 \cdot \frac{\frac{1}{x+1}(x+1) - \ln(x+1)\cdot 1}{(x+1)^2} = 100 \cdot \frac{1 - \ln(x+1)}{(x+1)^2}.
\]
\item Signe de $R'$ : le dénominateur est $>0$. Le numérateur $1 - \ln(x+1)$ est :
\begin{itemize}
\item $>0$ si $\ln(x+1) < 1 \iff x+1 < e \iff x < e-1 \approx \num{1,718}$,
\item $=0$ si $x = e-1$,
\item $<0$ si $x > e-1$.
\end{itemize}
\item Tableau de variations :
\begin{center}
\begin{tabular}{|c|ccccc|}\hline
$x$ & $0$ & & $e-1$ & & $100$ \\ \hline
$R'(x)$ & & $+$ & $0$ & $-$ & \\ \hline
$R(x)$ & $0$ & $\nearrow$ & $R(e-1)$ & $\searrow$ & $R(100)$ \\ \hline
\end{tabular}
\end{center}
\item \textbf{Conclusion :} le rendement est maximal pour $x = e-1 \approx \num{1,72}\;^\circ\text{C}$ (arrondi à \num{1,7} °C). La valeur maximale est $R(e-1) = 100 \cdot \frac{1}{e} \approx \num{36,8}\%$.
\end{enumerate}
\end{exoresolu}

\begin{exercice}[Entraînement : dériver et étudier]
Soit $f(x) = \ln(x^2 - 4x + 5)$.
\begin{enumerate}
\item Déterminer l'ensemble de définition $D_f$ de $f$.
\item Calculer $f'(x)$ sur $D_f$.
\item Étudier le signe de $f'(x)$ et dresser le \cle{tableau de variations} de $f$.
\item En déduire l'extremum de $f$ et sa valeur.
\end{enumerate}
\end{exercice}

\begin{corrige}[Exercice : dériver et étudier]
\begin{enumerate}
\item $u(x) = x^2 - 4x + 5 = (x-2)^2 + 1 \ge 1 > 0$ pour tout $x \in \mathbb{R}$. Donc $D_f = \mathbb{R}$.
\item $u'(x) = 2x - 4 = 2(x-2)$. Donc $f'(x) = \dfrac{2(x-2)}{x^2-4x+5}$ sur $\mathbb{R}$.
\item Le dénominateur est toujours $>0$. Le signe de $f'$ est celui du numérateur $2(x-2)$ :
$f'(x) < 0$ pour $x < 2$, $f'(2)=0$, $f'(x) > 0$ pour $x > 2$.
\item Tableau de variations :
\begin{center}
\begin{tabular}{|c|ccccc|}\hline
$x$ & $-\infty$ & & $2$ & & $+\infty$ \\ \hline
$f'(x)$ & & $-$ & $0$ & $+$ & \\ \hline
$f(x)$ & $+\infty$ & $\searrow$ & $\ln 1 = 0$ & $\nearrow$ & $+\infty$ \\ \hline
\end{tabular}
\end{center}
$f$ admet un \textbf{minimum global} en $x=2$, de valeur $f(2) = \ln(1) = 0$. Pas de maximum.
\end{enumerate}
\end{corrige}

\medskip
\noindent Nous disposons maintenant de tous les outils de dérivation nécessaires. La prochaine section (5) utilisera la relation $\int \frac{u'}{u} = \ln|u|$ pour calculer des primitives, avant de mener des études complètes de fonctions (section 6).

\coursec{Primitives des fonctions de la forme u'/u}

Après avoir maîtrisé la dérivation des fonctions composées faisant intervenir le logarithme népérien, nous abordons maintenant l'opération inverse : la \emph{primitivation}. C'est une étape fondamentale pour le calcul d'intégrales et la résolution d'équations différentielles simples, outils indispensables en électricité (charge d'un condensateur), en mécanique (lois de frottement) ou en gestion (croissance continue).

\courssub{Rappel et lecture inverse de la dérivation}

\begin{definition}[Primitive d'une fonction]
Soit $f$ une fonction définie sur un intervalle $I$. On appelle \cle{primitive} de $f$ sur $I$ toute fonction $F$ dérivable sur $I$ telle que $F'(x) = f(x)$ pour tout $x \in I$.
\end{definition}

Si $F$ est une primitive de $f$ sur $I$, alors toutes les primitives de $f$ sur $I$ s'écrivent $F(x) + k$ avec $k \in \mathbb{R}$. On note souvent l'ensemble des primitives $\int f(x) \, dx = F(x) + k$.

Rappelons le résultat fondamental de la section précédente : si $u$ est une fonction dérivable et \textbf{strictement positive} sur un intervalle $I$, alors la fonction $\ln(u)$ est dérivable sur $I$ et
\[
\bigl(\ln(u(x))\bigr)' = \frac{u'(x)}{u(x)}.
\]

En lisant cette égalité de droite à gauche, nous obtenons immédiatement un théorème majeur pour la primitivation.

\begin{propriete}[{Primitive de la forme u'/u (cas u > 0)] [Théorème — Démonstration exigible au Probatoire}]
Soit $u$ une fonction dérivable et \textbf{strictement positive} sur un intervalle $I$.
Les primitives de la fonction $f(x) = \dfrac{u'(x)}{u(x)}$ sur $I$ sont les fonctions
\[
F(x) = \ln\bigl(u(x)\bigr) + k \quad (k \in \mathbb{R}).
\]

\textbf{Preuve :} Posons $F(x) = \ln(u(x))$. Puisque $u$ est dérivable et $u(x) > 0$ sur $I$, la fonction composée $\ln \circ u$ est dérivable sur $I$. D'après la formule de dérivation de la composée (section précédente) :
\[
F'(x) = \frac{u'(x)}{u(x)} = f(x).
\]
Donc $F$ est bien une primitive de $f$ sur $I$. L'ensemble des primitives s'obtient en ajoutant une constante $k \in \mathbb{R}$.
\end{propriete}

\begin{exemplebox}[Application directe]
Soit $f(x) = \dfrac{2x}{x^2 + 3}$ définie sur $\mathbb{R}$.
Le dénominateur $u(x) = x^2 + 3$ est strictement positif (car $x^2 \ge 0 \Rightarrow x^2+3 \ge 3 > 0$).
Sa dérivée est $u'(x) = 2x$, qui est exactement le numérateur.
On a bien $f(x) = \dfrac{u'(x)}{u(x)}$.
Une primitive de $f$ sur $\mathbb{R}$ est donc $F(x) = \ln(x^2 + 3)$.
L'ensemble des primitives est $F(x) = \ln(x^2 + 3) + k,\; k \in \mathbb{R}$.
\end{exemplebox}

\courssub{Cas général : la valeur absolue}

Dans la pratique, la fonction $u$ au dénominateur peut prendre des valeurs négatives sur certains intervalles. Le logarithme népérien $\ln(u)$ n'est alors pas défini. Heureusement, on peut étendre le résultat grâce à la valeur absolue.

\begin{propriete}[{Primitive de la forme u'/u (cas général)] [Complément utile — Admis en Terminale}]
Soit $u$ une fonction dérivable qui \textbf{ne s'annule pas} sur un intervalle $I$ (donc $u(x) > 0$ ou $u(x) < 0$ pour tout $x \in I$).
Une primitive de la fonction $f(x) = \dfrac{u'(x)}{u(x)}$ sur $I$ est
\[
F(x) = \ln\bigl|u(x)\bigr|.
\]
L'ensemble des primitives est $\ln|u(x)| + k,\; k \in \mathbb{R}$.
\end{propriete}

\begin{attention}[Piège majeur : le signe de u et le domaine de définition]
\begin{itemize}
\item \textbf{N'écrivez jamais $\ln(u(x))$ si $u(x)$ peut être négatif.} Sur un intervalle où $u(x) < 0$, la primitive est $\ln(-u(x))$ car $|u(x)| = -u(x)$.
\item \textbf{Le domaine de la primitive doit être un intervalle où $u$ ne s'annule pas.} Si $u$ s'annule en $a$, l'ensemble de définition est morcelé (ex: $]-\infty; a[ \cup ]a; +\infty[$). Il existe alors \emph{deux} constantes $k_1$ et $k_2$ potentiellement différentes de part et d'autre de $a$.
\item \textbf{Vérifiez toujours le domaine de définition} de la fonction à primitiver avant d'écrire le résultat.
\end{itemize}
\end{attention}

\begin{exemplebox}[Cas avec valeur absolue]
Chercher les primitives de $f(x) = \dfrac{1}{x-2}$ sur son domaine de définition.
$f$ est définie sur $]-\infty; 2[ \cup ]2; +\infty[$.
On reconnaît $u(x) = x-2$, $u'(x) = 1$. $u$ ne s'annule pas sur chaque intervalle.
Une primitive sur $]2; +\infty[$ (où $u>0$) : $F(x) = \ln(x-2) + k_1$.
Une primitive sur $]-\infty; 2[$ (où $u<0$) : $F(x) = \ln(2-x) + k_2$ (car $|x-2| = 2-x$).
On écrit souvent de façon compacte : $F(x) = \ln|x-2| + k$, en gardant à l'esprit que $k$ peut changer de valeur de part et d'autre de $2$.
\end{exemplebox}

\begin{popfigure}
\begin{tikzpicture}[scale=1.1, every node/.style={font=\small}]
% Tableau de correspondance
\draw[popline, thick] (0,0) grid (12,4);
% En-têtes
\fill[popblue!20] (0,3) rectangle (12,4);
\node[align=center] at (1.5,3.5) {\textbf{Fonction $f(x)$}};
\node[align=center] at (5,3.5) {\textbf{Forme reconnue}};
\node[align=center] at (8.5,3.5) {\textbf{Primitive $F(x)$}};

% Ligne 1
\node[align=center] at (1.5,2.5) {$\dfrac{u'(x)}{u(x)}$};
\node[align=center] at (5,2.5) {$u>0$ sur $I$};
\node[align=center] at (8.5,2.5) {$\ln(u(x)) + k$};

% Ligne 2
\node[align=center] at (1.5,1.5) {$\dfrac{u'(x)}{u(x)}$};
\node[align=center] at (5,1.5) {$u \neq 0$ sur $I$};
\node[align=center] at (8.5,1.5) {$\ln|u(x)| + k$};

% Ligne 3 : Cas particuliers usuels
\fill[popgreen!15] (0,1) rectangle (12,2);
\node[align=center] at (1.5,1.5) {$\dfrac{1}{x}$};
\node[align=center] at (5,1.5) {$u=x,\; u'=1$};
\node[align=center] at (8.5,1.5) {$\ln|x| + k$};

% Ligne 4
\node[align=center] at (1.5,0.5) {$\dfrac{a}{ax+b}$};
\node[align=center] at (5,0.5) {$u=ax+b,\; u'=a$};
\node[align=center] at (8.5,0.5) {$\ln|ax+b| + k$};
\end{tikzpicture}
\legende{Tableau de correspondance dérivation $\leftrightarrow$ primitivation pour la forme $u'/u$.}
\end{popfigure}

\courssub{Primitives usuelles et ajustement du coefficient}

Deux primitives usuelles découlent immédiatement du théorème précédent et doivent être connues par cœur.

\begin{aretenir}[Primitives usuelles à connaître]
\begin{enumerate}
\item Sur $]0; +\infty[$ : $\displaystyle \int \frac{1}{x} \, dx = \ln x + k$.
\item Sur $\mathbb{R}^*$ : $\displaystyle \int \frac{1}{x} \, dx = \ln|x| + k$.
\item Pour $a \neq 0$ : $\displaystyle \int \frac{1}{ax+b} \, dx = \frac{1}{a} \ln|ax+b| + k$.
\end{enumerate}
\end{aretenir}

Souvent, le numérateur n'est pas \emph{exactement} la dérivée du dénominateur, mais en diffère seulement par un facteur constant. Il faut alors \textbf{ajuster le coefficient} (méthode du « multiplier et diviser par la même constante »).

\begin{methode}[Reconnaître et primitiver la forme $u'/u$ (ajustement du coefficient)]
Pour primitiver une fraction rationnelle $\dfrac{N(x)}{D(x)}$ :
\begin{enumerate}
\item Identifier un candidat pour le dénominateur $u(x) = D(x)$.
\item Calculer sa dérivée $u'(x)$.
\item Comparer $N(x)$ et $u'(x)$.
\item Si $N(x) = \lambda \cdot u'(x)$ (avec $\lambda$ constante), alors on écrit :
\[
\frac{N(x)}{D(x)} = \lambda \cdot \frac{u'(x)}{u(x)}.
\]
\item Une primitive est $\lambda \ln|u(x)| + k$.
\item \textbf{Vérifier en dérivant} le résultat trouvé.
\end{enumerate}
\end{methode}

\begin{exoresolu}[Primitive de $f(x) = \dfrac{6x}{x^2 + 1}$]
\textbf{Énoncé :} Déterminer les primitives de $f(x) = \dfrac{6x}{x^2 + 1}$ sur $\mathbb{R}$.

\textbf{Solution détaillée :}
\begin{enumerate}
\item \textbf{Analyse :} Le dénominateur $u(x) = x^2 + 1$ est strictement positif sur $\mathbb{R}$ (pas besoin de valeur absolue, mais on peut la mettre). Sa dérivée est $u'(x) = 2x$.
\item \textbf{Comparaison :} Le numérateur est $6x$. On a $6x = 3 \times (2x) = 3 \cdot u'(x)$.
\item \textbf{Réécriture :}
\[
f(x) = \frac{6x}{x^2+1} = 3 \cdot \frac{2x}{x^2+1} = 3 \cdot \frac{u'(x)}{u(x)}.
\]
\item \textbf{Primitivation :} Une primitive de $\dfrac{u'}{u}$ est $\ln(u(x)) = \ln(x^2+1)$.
Donc une primitive de $f$ est $F(x) = 3 \ln(x^2+1)$.
\item \textbf{Réponse :} L'ensemble des primitives sur $\mathbb{R}$ est
\[
\boxed{F(x) = 3 \ln(x^2+1) + k,\quad k \in \mathbb{R}}.
\]
\item \textbf{Vérification (réflexe obligatoire) :}
$F'(x) = 3 \cdot \dfrac{2x}{x^2+1} = \dfrac{6x}{x^2+1} = f(x)$. \checkmark
\end{enumerate}
\end{exoresolu}

\begin{exoresolu}[Primitive de $f(x) = \dfrac{5}{3x-2}$]
\textbf{Énoncé :} Déterminer les primitives de $f(x) = \dfrac{5}{3x-2}$ sur son domaine de définition.

\textbf{Solution :}
\begin{enumerate}
\item Domaine : $3x-2 \neq 0 \Leftrightarrow x \neq \frac{2}{3}$. Deux intervalles : $]-\infty; \frac{2}{3}[$ et $]\frac{2}{3}; +\infty[$.
\item On pose $u(x) = 3x-2$, donc $u'(x) = 3$.
\item Le numérateur est $5$. On écrit $5 = \frac{5}{3} \times 3 = \frac{5}{3} u'(x)$.
\item $f(x) = \dfrac{5}{3x-2} = \dfrac{5}{3} \cdot \dfrac{3}{3x-2} = \dfrac{5}{3} \cdot \dfrac{u'(x)}{u(x)}$.
\item Primitive : $F(x) = \dfrac{5}{3} \ln|3x-2| + k$.
\item \textbf{Vérification :} $F'(x) = \dfrac{5}{3} \cdot \dfrac{3}{3x-2} = \dfrac{5}{3x-2} = f(x)$. \checkmark
\end{enumerate}
\end{exoresolu}

\courssub{Vérification par dérivation et application au calcul d'aires}

La dérivation est l'opération inverse de la primitivation. \textbf{Il est impératif de systématiquement dériver la primitive trouvée pour s'assurer qu'on retrouve la fonction de départ.} Cela permet de corriger les erreurs de coefficient, de signe ou d'oubli de la valeur absolue.

\begin{attention}[Erreurs fréquentes à corriger par la vérification]
\begin{itemize}
\item Oublier le facteur $\frac{1}{a}$ pour $\int \frac{1}{ax+b} dx$ (écrire $\ln|ax+b|$ au lieu de $\frac{1}{a}\ln|ax+b|$).
\item Oublier la valeur absolue quand le dénominateur peut être négatif.
\item Confondre $\ln(u^2)$ et $2\ln|u|$ (rappel : $\ln(u^2) = 2\ln|u|$).
\end{itemize}
\end{attention}

L'intégration (calcul d'aires) est l'application concrète majeure des primitives. Le logarithme népérien apparaît naturellement pour calculer l'aire sous l'hyperbole $y = \frac{1}{x}$.

\begin{exemplebox}[Aire sous l'hyperbole $y = 1/x$ entre 1 et $e$]
\textbf{Problème :} Calculer l'aire $\mathcal{A}$ du domaine plan borné par la courbe $\mathcal{C}_f$ de $f(x) = \frac{1}{x}$, l'axe des abscisses, et les droites $x=1$ et $x=e$.

\textbf{Résolution :}
\begin{enumerate}
\item $f$ est continue et positive sur $[1; e]$. L'aire est $\mathcal{A} = \int_1^e \frac{1}{x} \, dx$.
\item Une primitive de $\frac{1}{x}$ sur $]0; +\infty[$ est $\ln x$.
\item D'après le théorème fondamental de l'analyse (lien primitives-intégrales) :
\[
\mathcal{A} = \Bigl[ \ln x \Bigr]_1^e = \ln e - \ln 1.
\]
\item Comme $\ln e = 1$ et $\ln 1 = 0$, on obtient $\mathcal{A} = 1 - 0 = 1$.
\end{enumerate}
\textbf{Conclusion :} L'aire sous l'hyperbole $y=1/x$ entre $1$ et $e$ vaut exactement \textbf{1 unité d'aire}. C'est une propriété caractéristique du nombre $e$.
\end{exemplebox}

\begin{popfigure}
\begin{tikzpicture}[scale=1.2, domain=0.2:4]
\begin{axis}[
    axis lines=middle,
    xlabel=$x$, ylabel=$y$,
    xmin=0, xmax=3.5,
    ymin=0, ymax=1.5,
    xtick={1,2.718}, xticklabels={$1$, $e$},
    ytick={1}, yticklabels={$1$},
    clip=false,
    width=\linewidth,
    height=8cm,
    samples=200,
    enlargelimits=0.1
]
% Courbe y=1/x
\addplot[popcurve, thick, domain=0.3:3.3] {1/x};
\node[popdark, anchor=south west] at (axis cs:2.5,0.4) {$y=\frac{1}{x}$};

% Aire hachurée entre 1 et e
\addplot[popblue!30, domain=1:2.718] {1/x} \closedcycle;

% Lignes verticales
\draw[dashed, popmuted] (axis cs:1,0) -- (axis cs:1,1);
\draw[dashed, popmuted] (axis cs:2.718,0) -- (axis cs:2.718,0.368);

% Annotation aire = 1
\node[popblue, align=center, font=\small\bfseries] at (axis cs:1.8,0.6) {Aire $= 1$};
\end{axis}
\end{tikzpicture}
\legende{Représentation graphique de l'aire $\int_1^e \frac{1}{x}dx = 1$.}
\end{popfigure}

\begin{savaistu}[Le nombre $e$ et l'aire de l'hyperbole]
Le nombre $e \approx 2,71828$ est \emph{défini} en analyse comme l'unique nombre réel $>1$ tel que l'aire sous la courbe $y=1/x$ entre $1$ et $e$ soit exactement égale à $1$. C'est une définition géométrique équivalente à la définition par la limite $\lim_{n\to\infty}(1+1/n)^n$ ou à la série $\sum 1/n!$. Cela montre le lien profond entre le logarithme népérien (aire) et l'exponentielle (croissance).
\end{savaistu}

\courssub{Synthèse et exercices d'entraînement}

\begin{aretenir}[Résumé : Primitives de $u'/u$]
\begin{center}
\begin{tabular}{|c|c|c|}\hline
\textbf{Forme de $f(x)$} & \textbf{Condition sur $u$} & \textbf{Primitive $F(x)$} \\ \hline
$\dfrac{u'(x)}{u(x)}$ & $u(x) > 0$ sur $I$ & $\ln(u(x)) + k$ \\ \hline
$\dfrac{u'(x)}{u(x)}$ & $u(x) \neq 0$ sur $I$ & $\ln|u(x)| + k$ \\ \hline
$\dfrac{1}{x}$ & $x > 0$ & $\ln x + k$ \\ \hline
$\dfrac{1}{x}$ & $x \in \mathbb{R}^*$ & $\ln|x| + k$ \\ \hline
$\dfrac{a}{ax+b}\ (a\neq 0)$ & $ax+b \neq 0$ & $\ln|ax+b| + k$ \\ \hline
$\dfrac{\lambda u'(x)}{u(x)}$ & $u(x) \neq 0$ sur $I$ & $\lambda \ln|u(x)| + k$ \\ \hline
\end{tabular}
\end{center}
\textbf{Règle d'or :} Toujours vérifier par la dérivation : $F'(x) \stackrel{?}{=} f(x)$.
\end{aretenir}

\begin{exercice}[Entraînement]
Déterminer les primitives des fonctions suivantes sur leur domaine de définition (ne pas oublier la vérification) :
\begin{enumerate}
\item $f_1(x) = \dfrac{4x^3}{x^4 + 5}$
\item $f_2(x) = \dfrac{3}{2x-1}$
\item $f_3(x) = \dfrac{2x-1}{x^2 - x + 2}$ \quad (indice : dérivée du dénominateur ?)
\item $f_4(x) = \dfrac{\cos x}{\sin x}$ sur $]0; \pi[$
\item $f_5(x) = \dfrac{e^x}{e^x + 1}$ sur $\mathbb{R}$
\end{enumerate}
\end{exercice}

\begin{corrige}[Exercice n°12 - Primitives u'/u]
\begin{enumerate}
\item $u=x^4+5>0,\ u'=4x^3 \Rightarrow F_1(x) = \ln(x^4+5) + k$.
\item $u=2x-1,\ u'=2 \Rightarrow \frac{3}{2x-1} = \frac{3}{2}\frac{2}{2x-1} \Rightarrow F_2(x) = \frac{3}{2}\ln|2x-1| + k$.
\item $u=x^2-x+2$ (discriminant $\Delta = -7 < 0$, toujours $>0$), $u'=2x-1 \Rightarrow F_3(x) = \ln(x^2-x+2) + k$.
\item $u=\sin x > 0$ sur $]0;\pi[$, $u'=\cos x \Rightarrow F_4(x) = \ln(\sin x) + k$.
\item $u=e^x+1 > 0$, $u'=e^x \Rightarrow F_5(x) = \ln(e^x+1) + k$.
\end{enumerate}
\end{corrige}

\coursec{Étude de fonctions faisant intervenir ln et applications concrètes}

Dans la section précédente, nous avons appris à trouver les primitives des fonctions de la forme $\dfrac{u'}{u}$ : c'est le logarithme népérien qui apparaît naturellement dans la réponse, puisque $\ln u$ est une primitive de $\dfrac{u'}{u}$ (pour $u>0$). Nous allons maintenant faire le chemin inverse : utiliser tout ce que nous savons sur $\ln$ (définition, propriétés algébriques, limites, dérivées) pour \cle{étudier complètement} des fonctions qui font intervenir le logarithme népérien. C'est l'exercice le plus classique du Probatoire et du Baccalauréat technique : « Étudier les variations de $f$ et tracer sa courbe représentative. » Nous terminerons par deux applications concrètes : le temps de doublement d'un capital et le pH d'une solution.

\courssub{Plan type d'étude d'une fonction avec ln}

Lorsqu'on étudie une fonction, on ne procède jamais au hasard : on suit toujours le même plan, dans le même ordre. À l'examen, le correcteur attend cette organisation ; la respecter, c'est déjà gagner des points de présentation.

\begin{methode}[Plan d'étude d'une fonction avec ln — 6 étapes à rédiger dans cet ordre]
\begin{enumerate}
\item \textbf{Ensemble de définition.} On cherche les valeurs de $x$ pour lesquelles $f(x)$ existe. Règle d'or : $\ln(\text{quelque chose})$ n'existe que si ce « quelque chose » est \cle{strictement positif}.
\item \textbf{Limites aux bornes} de l'ensemble de définition (en $0^+$ et en $+\infty$ le plus souvent). On utilise les limites de référence : $\lim_{x\to 0^+}\ln x=-\infty$, $\lim_{x\to+\infty}\ln x=+\infty$, $\lim_{x\to+\infty}\dfrac{\ln x}{x}=0$, $\lim_{x\to 0^+}x\ln x=0$.
\item \textbf{Calcul de la dérivée.} On utilise $(\ln x)'=\dfrac{1}{x}$, $(\ln u)'=\dfrac{u'}{u}$ et les règles de dérivation usuelles (produit, quotient).
\item \textbf{Signe de la dérivée.} On factorise $f'(x)$, on résout $f'(x)>0$ et $f'(x)<0$ sur l'ensemble de définition.
\item \textbf{Tableau de variations.} On y reporte les bornes, les limites, le signe de $f'$ et les valeurs remarquables de $f$.
\item \textbf{Courbe représentative.} On place les points remarquables, les tangentes horizontales, les asymptotes éventuelles, puis on trace.
\end{enumerate}
\end{methode}

\begin{attention}[L'erreur qui coûte cher : oublier les limites aux bornes]
Beaucoup d'élèves calculent la dérivée et dressent le tableau de variations \emph{sans avoir étudié les limites}, en particulier la limite en $0^+$. Résultat : le tableau est incomplet et la courbe est fausse près de l'axe des ordonnées. Une fonction avec $\ln$ a presque toujours un comportement spectaculaire en $0^+$ (asymptote verticale ou prolongement par continuité) : \textbf{ne jamais sauter l'étape 2}. De même, ne jamais écrire $\ln 0$ : cette écriture n'a aucun sens.
\end{attention}

\courssub{Étude complète de $f(x)=x-\ln x$ sur $]0\,;\,+\infty[$}

\begin{exemplebox}[Étude de $f(x)=x-\ln x$]
\textbf{1. Ensemble de définition.} $\ln x$ n'existe que pour $x>0$, donc $D_f=]0\,;\,+\infty[$.

\textbf{2. Limites aux bornes.}
En $0^+$ : $x\to 0$ et $\ln x\to-\infty$, donc $-\ln x\to+\infty$ et $\lim_{x\to 0^+}f(x)=+\infty$.
En $+\infty$ : on a une forme indéterminée « $+\infty-\infty$ ». On lève l'indétermination en factorisant par $x$ :
\[ f(x)=x\left(1-\frac{\ln x}{x}\right). \]
Or $\lim_{x\to+\infty}\dfrac{\ln x}{x}=0$, donc la parenthèse tend vers $1$ et, par produit, $\lim_{x\to+\infty}f(x)=+\infty$.

\textbf{3. Dérivée.} $f'(x)=1-\dfrac{1}{x}=\dfrac{x-1}{x}$.

\textbf{4. Signe de la dérivée.} Sur $]0\,;\,+\infty[$, $x>0$, donc $f'(x)$ a le signe de $x-1$ : $f'(x)<0$ sur $]0\,;\,1[$, $f'(1)=0$, $f'(x)>0$ sur $]1\,;\,+\infty[$.

\textbf{5. Tableau de variations.} $f(1)=1-\ln 1=1-0=1$ : c'est le \cle{minimum} de $f$.
\begin{center}
\begin{tabular}{|c|ccccc|}\hline
$x$ & $0$ & & $1$ & & $+\infty$ \\ \hline
$f'(x)$ & & $-$ & $0$ & $+$ & \\ \hline
$f$ & $+\infty$ & $\searrow$ & $1$ & $\nearrow$ & $+\infty$ \\ \hline
\end{tabular}
\end{center}

\textbf{6. Conséquence.} Pour tout $x>0$, $f(x)\geq 1>0$, donc $x-\ln x>0$, c'est-à-dire $\ln x < x$. On retrouve une inégalité classique : la courbe de $\ln$ est toujours \emph{en dessous} de la droite d'équation $y=x$.
\end{exemplebox}

\courssub{Étude de $f(x)=\dfrac{\ln x}{x}$ sur $]0\,;\,+\infty[$}

\begin{exoresolu}[Étude complète de $f(x)=\dfrac{\ln x}{x}$]
Énoncé : étudier les variations de $f(x)=\dfrac{\ln x}{x}$ sur $]0\,;\,+\infty[$, préciser les limites aux bornes et tracer l'allure de la courbe.
\tcblower
\textbf{Solution détaillée.}

\textbf{1. Ensemble de définition :} $D_f=]0\,;\,+\infty[$ (à cause de $\ln x$).

\textbf{2. Limites aux bornes.}
En $0^+$ : $\ln x\to-\infty$ et $\dfrac{1}{x}\to+\infty$, donc par produit $\lim_{x\to 0^+}\dfrac{\ln x}{x}=-\infty$. La droite d'équation $x=0$ (l'axe des ordonnées) est \cle{asymptote verticale} à la courbe.
En $+\infty$ : c'est une limite de référence, $\lim_{x\to+\infty}\dfrac{\ln x}{x}=0$. La droite d'équation $y=0$ (l'axe des abscisses) est \cle{asymptote horizontale} à la courbe en $+\infty$.

\textbf{3. Dérivée} (dérivée d'un quotient $\dfrac{u}{v}$ avec $u=\ln x$, $u'=\dfrac{1}{x}$, $v=x$, $v'=1$) :
\[ f'(x)=\frac{u'v-uv'}{v^2}=\frac{\dfrac{1}{x}\times x-\ln x\times 1}{x^2}=\frac{1-\ln x}{x^2}. \]

\textbf{4. Signe de la dérivée.} Pour tout $x>0$, $x^2>0$, donc $f'(x)$ a le signe de $1-\ln x$ :
\[ 1-\ln x>0 \iff \ln x<1 \iff x<e. \]
Donc $f$ est strictement croissante sur $]0\,;\,e]$ et strictement décroissante sur $[e\,;\,+\infty[$.

\textbf{5. Tableau de variations.} $f(e)=\dfrac{\ln e}{e}=\dfrac{1}{e}\approx 0,37$ : c'est le \cle{maximum} de $f$.
\begin{center}
\begin{tabular}{|c|ccccc|}\hline
$x$ & $0$ & & $e$ & & $+\infty$ \\ \hline
$f'(x)$ & & $+$ & $0$ & $-$ & \\ \hline
$f$ & $-\infty$ & $\nearrow$ & $\dfrac{1}{e}$ & $\searrow$ & $0$ \\ \hline
\end{tabular}
\end{center}

\textbf{6. Points remarquables.} $f(1)=\dfrac{\ln 1}{1}=0$ : la courbe coupe l'axe des abscisses au point $(1\,;\,0)$. Au point $\left(e\,;\,\dfrac{1}{e}\right)$, la tangente est \cle{horizontale} car $f'(e)=0$.
\end{exoresolu}

\begin{popfigure}
\begin{center}
\begin{tikzpicture}
\begin{axis}[
width=0.85\linewidth, height=6.5cm,
axis lines=middle,
xlabel={$x$}, ylabel={$y$},
xmin=-0.5, xmax=10.5, ymin=-1.6, ymax=1.2,
xtick={1,2.718,5,8}, xticklabels={$1$,$e$,$5$,$8$},
ytick={0.368}, yticklabels={$\frac{1}{e}$},
grid=both, grid style={popline!40},
]
\addplot[popcurve,domain=0.35:10,samples=300]{ln(x)/x};
\addplot[poporange,dashed,domain=-0.5:10.5]{0.368};
\addplot[popmuted,dashed] coordinates {(2.718,-1.6) (2.718,0.368)};
\node[popdark,above right] at (axis cs:2.718,0.368) {maximum $\left(e\,;\,\frac{1}{e}\right)$};
\node[popmuted,right] at (axis cs:6,0.07) {asymptote $y=0$};
\end{axis}
\end{tikzpicture}
\end{center}
\legende{Courbe de $f(x)=\dfrac{\ln x}{x}$ : maximum en $\left(e\,;\,\frac{1}{e}\right)$, asymptote horizontale $y=0$ en $+\infty$, asymptote verticale $x=0$.}
\end{popfigure}

\courssub{Étude de $f(x)=x\ln x$ et prolongement par continuité en $0$}

La fonction $f(x)=x\ln x$ est définie sur $]0\,;\,+\infty[$. Mais on sait que $\lim_{x\to 0^+}x\ln x=0$ (limite de référence). On peut donc « boucher le trou » en $0$.

\begin{definition}[Prolongement par continuité]
Soit $f$ une fonction définie sur $]0\,;\,+\infty[$. Si $\lim_{x\to 0^+}f(x)=\ell$ où $\ell$ est un nombre \emph{fini}, alors la fonction $g$ définie par $g(x)=f(x)$ pour $x>0$ et $g(0)=\ell$ est continue en $0$ : on l'appelle le \cle{prolongement par continuité} de $f$ en $0$.
\end{definition}

Ainsi $x\ln x$ se prolonge par continuité en $0$ par la valeur $0$ : la courbe « part » de l'origine.

\begin{exemplebox}[Étude de $f(x)=x\ln x$]
\textbf{1.} $D_f=]0\,;\,+\infty[$, et $f$ se prolonge en $0$ par $f(0)=0$.

\textbf{2. Limites.} En $0^+$ : $\lim_{x\to 0^+}x\ln x=0$ (limite de référence). En $+\infty$ : $\lim_{x\to+\infty}x\ln x=+\infty$ (produit de deux facteurs qui tendent vers $+\infty$).

\textbf{3. Dérivée} (dérivée d'un produit $uv$ avec $u=x$, $v=\ln x$) :
\[ f'(x)=1\times\ln x+x\times\frac{1}{x}=\ln x+1. \]

\textbf{4. Signe.} $f'(x)>0 \iff \ln x>-1 \iff x>e^{-1}=\dfrac{1}{e}\approx 0,37$. Donc $f$ est strictement décroissante sur $\left]0\,;\,\frac{1}{e}\right]$ puis strictement croissante sur $\left[\frac{1}{e}\,;\,+\infty\right[$.

\textbf{5. Minimum.} $f\left(\dfrac{1}{e}\right)=\dfrac{1}{e}\ln\dfrac{1}{e}=\dfrac{1}{e}\times(-1)=-\dfrac{1}{e}\approx -0,37$.
\begin{center}
\begin{tabular}{|c|ccccc|}\hline
$x$ & $0$ & & $\dfrac{1}{e}$ & & $+\infty$ \\ \hline
$f'(x)$ & & $-$ & $0$ & $+$ & \\ \hline
$f$ & $0$ & $\searrow$ & $-\dfrac{1}{e}$ & $\nearrow$ & $+\infty$ \\ \hline
\end{tabular}
\end{center}

\textbf{6. Point remarquable :} $f(1)=1\times\ln 1=0$ : la courbe repasse par le point $(1\,;\,0)$.
\end{exemplebox}

\begin{popfigure}
\begin{center}
\begin{tikzpicture}
\begin{axis}[
width=0.85\linewidth, height=6.5cm,
axis lines=middle,
xlabel={$x$}, ylabel={$y$},
xmin=-0.4, xmax=3.4, ymin=-0.9, ymax=2.6,
xtick={0.368,1,2,3}, xticklabels={$\frac{1}{e}$,$1$,$2$,$3$},
ytick={-0.368}, yticklabels={$-\frac{1}{e}$},
grid=both, grid style={popline!40},
]
\addplot[popcurve,domain=0.012:3,samples=300]{x*ln(x)};
\addplot[poporange,dashed,domain=-0.4:3.4]{-0.368};
\addplot[popmuted,dashed] coordinates {(0.368,-0.368) (0.368,0)};
\node[popdark,below left] at (axis cs:0.55,-0.368) {minimum $\left(\frac{1}{e}\,;\,-\frac{1}{e}\right)$};
\fill[popgreen] (axis cs:0,0) circle (1.5pt);
\node[popgreen,above left] at (axis cs:0,0) {prolongement};
\end{axis}
\end{tikzpicture}
\end{center}
\legende{Courbe de $f(x)=x\ln x$ prolongée en $0$ : minimum $-\dfrac{1}{e}$ atteint en $x=\dfrac{1}{e}$, tangente horizontale en ce point.}
\end{popfigure}

\begin{attention}[Branches infinies et tangentes : ce qu'il faut vérifier]
\begin{itemize}
\item Si $\lim_{x\to 0^+}f(x)=\pm\infty$ : \textbf{asymptote verticale} d'équation $x=0$ (cas de $\dfrac{\ln x}{x}$).
\item Si $\lim_{x\to+\infty}f(x)=\ell$ avec $\ell$ fini : \textbf{asymptote horizontale} d'équation $y=\ell$.
\item Si $f'(x_0)=0$ : \textbf{tangente horizontale} au point d'abscisse $x_0$.
\item Piège : $x\ln x$ tend vers $0$ en $0^+$ (limite finie), mais $f'(x)=\ln x+1$ tend vers $-\infty$ en $0^+$ : la courbe arrive à l'origine avec une demi-tangente \emph{verticale}. Ne pas confondre « limite finie » et « tangente horizontale ».
\end{itemize}
\end{attention}

\courssub{Application 1 : temps de doublement d'un capital à intérêts composés continus}

Un technicien-comptable place un capital $C_0$ à intérêts composés continus au taux annuel $i$ (par exemple $i=0,08$ pour $8\,\%$). Le capital au bout de $t$ années est :
\[ C(t)=C_0\,e^{it}. \]
On cherche le temps $t$ de \cle{doublement} : $C(t)=2C_0$, c'est-à-dire $e^{it}=2$. En appliquant $\ln$ aux deux membres (ce qui est licite car les deux membres sont strictement positifs) :
\[ it=\ln 2 \quad\Longrightarrow\quad t=\dfrac{\ln 2}{i}. \]
Remarque remarquable : ce temps \textbf{ne dépend pas du capital initial} ! Doubler prend autant de temps pour \num{100000} FCFA que pour 10 millions de FCFA.

\begin{exemplebox}[Capital placé à $8\,\%$ en continu]
Un capital croît de $8\,\%$ par an en intérêts composés continus. Temps de doublement :
\[ t=\frac{\ln 2}{0,08}\approx\frac{0,693}{0,08}\approx 8,7 \text{ ans}. \]
Conclusion rédigée : « Le capital double en environ $8$ ans et $8$ mois, quel que soit le montant placé. »
\end{exemplebox}

\courssub{Application 2 : pH d'une solution}

En chimie, l'acidité d'une solution est mesurée par le \cle{pH}, défini à partir de la concentration en ions hydrogène $[H^+]$ (en mol/L) :
\[ \mathrm{pH}=-\frac{\ln[H^+]}{\ln 10}=-\log[H^+], \]
où $\log$ désigne le logarithme décimal (complément signalé : $\log x=\dfrac{\ln x}{\ln 10}$). Une solution est acide si $\mathrm{pH}<7$, neutre si $\mathrm{pH}=7$, basique si $\mathrm{pH}>7$.

\begin{exoresolu}[Calcul de pH]
Énoncé : une solution a une concentration $[H^+]=10^{-4}$ mol/L. Calculer son pH et conclure. Que devient le pH si la concentration est multipliée par $10$ ?
\tcblower
\textbf{Solution.} $\mathrm{pH}=-\dfrac{\ln 10^{-4}}{\ln 10}=-\dfrac{-4\ln 10}{\ln 10}=4$. La solution est \textbf{acide} ($\mathrm{pH}<7$).

Si $[H^+]$ est multipliée par $10$ : la nouvelle concentration est $10^{-3}$ mol/L, d'où $\mathrm{pH}=3$. Grâce aux propriétés algébriques du logarithme ($\ln(10a)=\ln 10+\ln a$), multiplier la concentration par $10$ \textbf{diminue le pH de 1 unité}.
\end{exoresolu}

\begin{savaistu}[Les échelles logarithmiques autour de nous]
L'échelle de Richter des séismes et l'échelle des décibels (bruit) sont des échelles \emph{logarithmiques} : un séisme de magnitude $7$ dégage une énergie bien supérieure à celle d'un séisme de magnitude $6$, et un son de $70$ dB est $10$ fois plus intense qu'un son de $60$ dB. Le logarithme permet d'écraser des grandeurs énormes sur une échelle lisible — c'est exactement la propriété $\ln(10x)=\ln x+\ln 10$ à l'œuvre.
\end{savaistu}

\courssub{Rédaction d'une conclusion de problème concret}

\begin{methode}[Rédiger une conclusion justifiée à l'examen]
\begin{enumerate}
\item \textbf{Traduire} la question concrète en langage mathématique (équation, inéquation, recherche d'extremum).
\item \textbf{Résoudre} en utilisant les outils de la leçon (propriétés de $\ln$, dérivée, tableau de variations).
\item \textbf{Conclure en français}, avec une phrase complète, l'unité et une vérification de cohérence (un temps positif, un pH entre $0$ et $14$, etc.).
\end{enumerate}
Exemple de conclusion attendue : « D'après le tableau de variations, le coût unitaire maximal est atteint pour $x=e$ pièces produites et vaut $2+\frac{1}{e}$ milliers de FCFA, soit environ \num{2368} FCFA par pièce. »
\end{methode}

\begin{exercice}[À faire — synthèse]
Un atelier estime que le coût unitaire de production, en milliers de FCFA, de $x$ pièces ($x>0$) est $C(x)=\dfrac{\ln x}{x}+2$.
\begin{enumerate}
\item Étudier les variations de $C$ sur $]0\,;\,+\infty[$ (limites, dérivée, tableau).
\item En déduire la quantité qui rend le coût unitaire maximal, puis justifier que pour $x\geq 1$, le coût unitaire reste supérieur ou égal à $2$ milliers de FCFA.
\item Un capital est placé à $5\,\%$ en intérêts composés continus : calculer le temps de doublement (arrondir à $0,1$ an).
\end{enumerate}
\end{exercice}

\begin{corrige}[Exercice : synthèse]
\textbf{1.} $C$ est définie sur $]0\,;\,+\infty[$. $C'(x)=\dfrac{1-\ln x}{x^2}$, du signe de $1-\ln x$ car $x^2>0$. Donc $C$ croît sur $]0\,;\,e]$ et décroît sur $[e\,;\,+\infty[$. Limites : $\lim_{x\to 0^+}C(x)=-\infty$ (asymptote verticale $x=0$) ; $\lim_{x\to+\infty}C(x)=2$ car $\dfrac{\ln x}{x}\to 0$ : asymptote horizontale d'équation $y=2$. Maximum : $C(e)=\dfrac{1}{e}+2\approx 2,37$.
\begin{center}
\begin{tabular}{|c|ccccc|}\hline
$x$ & $0$ & & $e$ & & $+\infty$ \\ \hline
$C'(x)$ & & $+$ & $0$ & $-$ & \\ \hline
$C$ & $-\infty$ & $\nearrow$ & $2+\dfrac{1}{e}$ & $\searrow$ & $2$ \\ \hline
\end{tabular}
\end{center}
\textbf{2.} Le coût unitaire maximal est atteint pour $x=e\approx 2,72$, soit environ $3$ pièces. Pour $x\geq 1$, $\ln x\geq 0$ donc $\dfrac{\ln x}{x}\geq 0$ et $C(x)\geq 2$ : le coût unitaire reste supérieur ou égal à $2$ milliers de FCFA.
\textbf{3.} $t=\dfrac{\ln 2}{0,05}\approx\dfrac{0,693}{0,05}\approx 13,9$ ans.
\end{corrige}

\begin{aretenir}[L'essentiel de la section]
\begin{itemize}
\item Plan d'étude en 6 étapes, \textbf{toujours dans l'ordre} : ensemble de définition, limites aux bornes, dérivée, signe de la dérivée, tableau de variations, courbe.
\item Trois études de référence à connaître : $x-\ln x$ (minimum $1$ en $x=1$), $\dfrac{\ln x}{x}$ (maximum $\frac{1}{e}$ en $x=e$, asymptote $y=0$), $x\ln x$ (prolongée par $0$ en $0$, minimum $-\frac{1}{e}$ en $x=\frac{1}{e}$).
\item Temps de doublement d'un capital à intérêts continus : $t=\dfrac{\ln 2}{i}$, indépendant du capital.
\item $\mathrm{pH}=-\dfrac{\ln[H^+]}{\ln 10}$ : multiplier $[H^+]$ par $10$ diminue le pH de $1$ unité.
\end{itemize}
\end{aretenir}

\newpage

\coursec{Activité d'intégration}

\coursec{Fonctions logarithme népérien}

\courssub{Objectifs de l'activité}

À l'issue de cette activité d'intégration, l'élève sera capable de :
\begin{itemize}
\item mobiliser les \cle{propriétés algébriques} du logarithme népérien ($\ln(ab)=\ln a+\ln b$, $\ln(a^n)=n\ln a$, $\ln(\mathrm{e}^x)=x$) pour résoudre des équations ;
\item résoudre des \cle{équations} de la forme $\mathrm{e}^{kt}=a$ à l'aide de la fonction $\ln$ ;
\item calculer la \cle{dérivée} d'une fonction définie à l'aide de l'exponentielle et interpréter un taux de croissance relatif ;
\item comparer deux modèles d'évolution (croissance exponentielle et croissance linéaire) par le calcul et par une lecture graphique ;
\item \cle{rédiger et justifier} une démarche, une réponse et une conclusion dans une note de synthèse à destination d'un décideur.
\end{itemize}

\courssub{Situation}

La coopérative des producteurs de cacao de Sangmélima, dans la région du Sud, a réalisé en fin de campagne un bénéfice exceptionnel. Lors de l'assemblée générale, les membres décident de placer une somme de $2\,000\,000$ FCFA pour préparer l'achat futur d'une nouvelle sécherie solaire.

Le directeur d'une microfinance de la ville propose deux formules de placement.

\textbf{Formule A (intérêts composés).} La valeur du placement, exprimée en millions de FCFA, évolue en fonction du temps $t$ (en années) selon le modèle :
\[ V(t) = 2\,\mathrm{e}^{0{,}07\,t}. \]

\textbf{Formule B (intérêts simples).} La valeur du placement, toujours en millions de FCFA, est donnée par :
\[ V_2(t) = 2 + 0{,}14\,t. \]

Le président de la coopérative, peu à l'aise avec les mathématiques financières, pose les questions suivantes aux élèves de la classe de Terminale technique du lycée voisin :
\begin{itemize}
\item « Au bout de combien d'années notre capital aura-t-il doublé ? Et triplé ? »
\item « Si nous avions placé $5\,000\,000$ FCFA au lieu de $2\,000\,000$, le temps de doublement aurait-il été plus long ? »
\item « Que signifie concrètement le taux de $7\,\%$ annoncé ? »
\item « Quelle formule est la plus avantageuse, et à partir de quand ? »
\end{itemize}

Les élèves disposent d'une calculatrice et d'un tableur ou d'un papier millimétré pour les représentations graphiques. On donne : $\ln 2 \approx 0{,}693$, $\ln 3 \approx 1{,}099$, $\ln 2{,}5 \approx 0{,}916$, $\mathrm{e}^{0{,}7} \approx 2{,}014$.

\courssub{Consignes}

\begin{enumerate}
\item[] \textbf{Partie 1 : Temps de doublement et de triplement.}
\item Résoudre l'équation $V(t) = 4$ et interpréter le résultat : au bout de combien d'années le capital initial aura-t-il doublé ? Arrondir à $0{,}1$ année près.
\item Résoudre de même l'équation $V(t) = 6$ : au bout de combien d'années le capital aura-t-il triplé ?

\item[] \textbf{Partie 2 : Un résultat étonnant.}
\item On suppose maintenant que le capital initial est de $5$ millions de FCFA, placé selon la formule A : $W(t) = 5\,\mathrm{e}^{0{,}07\,t}$. Résoudre l'équation $W(t) = 10$.
\item En utilisant la propriété $\ln(ab) = \ln a + \ln b$, expliquer pourquoi le temps de doublement ne dépend pas du capital initial. Que vaut ce temps de doublement en fonction du seul taux $0{,}07$ ?

\item[] \textbf{Partie 3 : Taux de croissance relatif.}
\item Calculer $V'(t)$, dérivée de la fonction $V$.
\item On définit le taux de croissance relatif $g(t) = \dfrac{V'(t)}{V(t)}$. Montrer que $g(t)$ est une fonction constante et donner sa valeur.
\item Interpréter ce résultat en langage financier : quel lien avec le « taux de $7\,\%$ » annoncé par la microfinance ?

\item[] \textbf{Partie 4 : Comparaison des deux formules.}
\item Sur un même graphique (papier millimétré ou calculatrice), représenter les fonctions $V$ et $V_2$ pour $t$ compris entre $0$ et $20$ ans. Conjecturer à partir de quelle année la formule A devient plus avantageuse.
\item On pose $h(t) = V(t) - V_2(t) = 2\,\mathrm{e}^{0{,}07\,t} - 2 - 0{,}14\,t$. Calculer $h'(t)$, étudier son signe, puis dresser le tableau de variations de $h$ sur $[0\,;\,20]$.
\item En déduire, par le calcul, à partir de quelle année entière le placement composé (formule A) devient strictement plus avantageux que le placement simple (formule B).

\item[] \textbf{Partie 5 : Note de synthèse.}
\item Rédiger une note d'une demi-page à destination du président de la coopérative. La note doit présenter la démarche, les résultats chiffrés (temps de doublement, de triplement, année de basculement entre les deux formules) et se conclure par un conseil argumenté et justifié.
\end{enumerate}

\courssub{Ressources à mobiliser}

\begin{aretenir}[Boîte à outils de la leçon]
\begin{itemize}
\item \cle{Équivalence fondamentale} : pour $a > 0$ et $b$ réel, $\ln a = b \iff a = \mathrm{e}^b$ ; en particulier $\ln(\mathrm{e}^x) = x$ et $\mathrm{e}^{\ln x} = x$ pour $x>0$.
\item \cle{Propriétés algébriques} : $\ln(ab) = \ln a + \ln b$ ; $\ln\left(\dfrac{a}{b}\right) = \ln a - \ln b$ ; $\ln(a^n) = n\ln a$.
\item \cle{Dérivées} : $(\mathrm{e}^{u})' = u'\,\mathrm{e}^{u}$ ; $(\ln u)' = \dfrac{u'}{u}$ pour $u > 0$.
\item \cle{Étude de fonction} : le signe de $h'$ donne le sens de variation de $h$ ; le tableau de variations permet de localiser les solutions d'une équation $h(t) = 0$.
\item \cle{Rédaction} : toute réponse à une situation concrète doit être justifiée par un calcul et formulée dans le langage du contexte.
\end{itemize}
\end{aretenir}

\courssub{Démarche et corrigé commenté}

\begin{exoresolu}[Corrigé complet de l'activité]
\textbf{Partie 1.} 1) Résoudre $V(t)=4$. 2) Résoudre $V(t)=6$.

\textbf{Partie 2.} 3) Résoudre $W(t)=10$ avec $W(t)=5\,\mathrm{e}^{0{,}07t}$. 4) Montrer que le temps de doublement est indépendant du capital initial.

\textbf{Partie 3.} 5) Calculer $V'(t)$. 6) Étudier $g(t)=\dfrac{V'(t)}{V(t)}$. 7) Interpréter.

\textbf{Partie 4.} 8) Conjecture graphique. 9) Étudier $h(t)=2\,\mathrm{e}^{0{,}07t}-2-0{,}14t$. 10) Conclure sur la comparaison.

\textbf{Partie 5.} 11) Rédiger la note de synthèse.
\tcblower
\textbf{Partie 1 : doublement et triplement.}

1) On résout $V(t) = 4$, c'est-à-dire $2\,\mathrm{e}^{0{,}07\,t} = 4$, soit $\mathrm{e}^{0{,}07\,t} = 2$. En appliquant le logarithme népérien aux deux membres (fonction strictement croissante, l'équivalence est conservée) :
\[ 0{,}07\,t = \ln 2 \quad\Longleftrightarrow\quad t = \frac{\ln 2}{0{,}07} \approx \frac{0{,}693}{0{,}07} \approx 9{,}9. \]
Le capital aura doublé au bout d'environ \textbf{9,9 années} (environ 9 ans et 11 mois).

2) De même, $V(t) = 6 \iff \mathrm{e}^{0{,}07\,t} = 3 \iff t = \dfrac{\ln 3}{0{,}07} \approx \dfrac{1{,}099}{0{,}07} \approx 15{,}7$. Le capital aura triplé au bout d'environ \textbf{15,7 années}.

\textbf{Partie 2 : indépendance vis-à-vis du capital initial.}

3) $W(t) = 10 \iff 5\,\mathrm{e}^{0{,}07\,t} = 10 \iff \mathrm{e}^{0{,}07\,t} = 2 \iff t = \dfrac{\ln 2}{0{,}07} \approx 9{,}9$ années. On retrouve exactement le même temps qu'avec $2$ millions !

4) Justification générale : si le capital initial est $C_0$ (en millions), le modèle est $V(t) = C_0\,\mathrm{e}^{0{,}07\,t}$. Le doublement correspond à $V(t) = 2C_0$, soit $\mathrm{e}^{0{,}07\,t} = 2$. En utilisant la propriété algébrique :
\[ \ln\left(C_0\,\mathrm{e}^{0{,}07\,t}\right) = \ln C_0 + 0{,}07\,t = \ln(2C_0) = \ln 2 + \ln C_0, \]
d'où $0{,}07\,t = \ln 2$ : le terme $\ln C_0$ s'élimine. Le temps de doublement $t = \dfrac{\ln 2}{0{,}07}$ ne dépend \textbf{que du taux}, jamais du capital initial. C'est une propriété fondamentale des intérêts composés.

\textbf{Partie 3 : taux de croissance relatif.}

5) $V(t) = 2\,\mathrm{e}^{0{,}07\,t}$ est de la forme $2\,\mathrm{e}^{u(t)}$ avec $u(t) = 0{,}07\,t$, donc $u'(t) = 0{,}07$ et :
\[ V'(t) = 2 \times 0{,}07\,\mathrm{e}^{0{,}07\,t} = 0{,}14\,\mathrm{e}^{0{,}07\,t}. \]

6) On calcule le quotient :
\[ g(t) = \frac{V'(t)}{V(t)} = \frac{0{,}14\,\mathrm{e}^{0{,}07\,t}}{2\,\mathrm{e}^{0{,}07\,t}} = \frac{0{,}14}{2} = 0{,}07. \]
La fonction $g$ est donc \textbf{constante}, égale à $0{,}07$.

7) Interprétation : à chaque instant, la vitesse de croissance du placement représente exactement $7\,\%$ de sa valeur. C'est précisément le sens du « taux de $7\,\%$ » de la formule A : un taux de croissance \emph{relatif} constant. C'est ce qui distingue les intérêts composés : le gain annuel augmente avec le capital, car le taux s'applique à une valeur de plus en plus grande.

\textbf{Partie 4 : comparaison des deux formules.}

8) Tableau de valeurs pour la conjecture graphique :
\begin{center}
\begin{tabular}{|c|cccccc|}
\hline
$t$ (années) & $0$ & $2$ & $5$ & $10$ & $15$ & $20$ \\
\hline
$V(t)$ (millions) & $2$ & $2{,}30$ & $2{,}84$ & $4{,}03$ & $5{,}72$ & $8{,}11$ \\
\hline
$V_2(t)$ (millions) & $2$ & $2{,}28$ & $2{,}70$ & $3{,}40$ & $4{,}10$ & $4{,}80$ \\
\hline
\end{tabular}
\end{center}
Les deux courbes partent du même point $(0\,;\,2)$, restent très proches au début, puis la courbe exponentielle s'écarte nettement. On conjecture que la formule A devient plus avantageuse très tôt, vers la deuxième année.

9) Posons $h(t) = 2\,\mathrm{e}^{0{,}07\,t} - 2 - 0{,}14\,t$ sur $[0\,;\,20]$. Sa dérivée est :
\[ h'(t) = 0{,}14\,\mathrm{e}^{0{,}07\,t} - 0{,}14 = 0{,}14\left(\mathrm{e}^{0{,}07\,t} - 1\right). \]
Pour $t > 0$, on a $0{,}07\,t > 0$, donc $\mathrm{e}^{0{,}07\,t} > \mathrm{e}^0 = 1$, d'où $h'(t) > 0$. La fonction $h$ est donc \textbf{strictement croissante} sur $[0\,;\,20]$.

\begin{center}
\begin{tabular}{|c|ccc|}
\hline
$t$ & $0$ & & $20$ \\
\hline
$h'(t)$ & $0$ & $+$ & $+$ \\
\hline
$h$ & $0$ & $\nearrow$ & $\approx 3{,}31$ \\
\hline
\end{tabular}
\end{center}

10) Comme $h(0) = 2 - 2 - 0 = 0$ et que $h$ est strictement croissante, on a $h(t) > 0$ pour tout $t > 0$. Autrement dit, $V(t) > V_2(t)$ dès que $t > 0$ : le placement composé est strictement plus avantageux \textbf{dès la première année écoulée} (à partir de l'année 1). L'écart devient spectaculaire avec le temps : à 20 ans, $8{,}11$ millions contre $4{,}80$ millions, soit environ $3{,}3$ millions de FCFA d'écart.

\textbf{Partie 5 : éléments de rédaction de la note de synthèse.}

\emph{Note à Monsieur le Président de la coopérative de cacao de Sangmélima.} Notre classe a étudié les deux formules proposées par la microfinance. Avec la formule A (intérêts composés au taux de $7\,\%$), le capital de $2\,000\,000$ FCFA double en environ $9{,}9$ ans et triple en environ $15{,}7$ ans ; ces durées ne dépendent pas de la somme placée, car le temps de doublement vaut toujours $\ln 2 / 0{,}07$. Nous avons vérifié que le taux de croissance relatif du placement est constant et égal à $7\,\%$, ce qui signifie que les intérêts produisent eux-mêmes des intérêts. La comparaison avec la formule B (intérêts simples) montre, par l'étude de la fonction différence $h$, que la formule A est plus avantageuse dès la première année, avec un écart d'environ $3{,}3$ millions de FCFA au bout de 20 ans. \emph{Conseil :} si la coopérative n'a pas besoin de ce capital à court terme, la formule A est nettement préférable pour financer la sécherie solaire.
\end{exoresolu}

\begin{attention}[Pièges fréquents dans cette activité]
\begin{itemize}
\item Ne pas oublier de \textbf{diviser par le coefficient} $2$ avant d'appliquer $\ln$ : de $2\,\mathrm{e}^{0{,}07t}=4$, on tire d'abord $\mathrm{e}^{0{,}07t}=2$, et non $0{,}07\,t = \ln 4$.
\item $\ln(\mathrm{e}^{0{,}07t}) = 0{,}07\,t$, et non $0{,}07$ ni $\mathrm{e}^{0{,}07t}$.
\item Dans la dérivée de $2\,\mathrm{e}^{0{,}07t}$, ne pas oublier le facteur $0{,}07$ issu de la dérivation de l'exposant.
\item Une conclusion financière doit être rédigée en \textbf{phrases complètes}, avec les unités (années, millions de FCFA) et une justification par le calcul.
\end{itemize}
\end{attention}

\courssub{Bilan de l'activité}

Cette activité a permis de mettre en évidence les points suivants :
\begin{itemize}
\item la fonction logarithme népérien est l'\textbf{outil de résolution} des équations où l'inconnue est en exposant : $\mathrm{e}^{kt} = a \iff t = \dfrac{\ln a}{k}$ ;
\item les propriétés algébriques de $\ln$ (notamment $\ln(ab) = \ln a + \ln b$) expliquent un résultat financier remarquable : le \textbf{temps de doublement d'un placement à intérêts composés ne dépend pas du capital initial}, mais uniquement du taux ;
\item le quotient $\dfrac{V'}{V}$ mesure le \textbf{taux de croissance relatif} ; pour un modèle exponentiel $C_0\,\mathrm{e}^{kt}$, ce taux est constant et égal à $k$, ce qui donne un sens concret au taux annoncé par un banquier ;
\item l'\textbf{étude complète d'une fonction} (dérivée, signe, variations) permet de comparer rigoureusement deux modèles d'évolution et de transformer une conjecture graphique en certitude démontrée ;
\item enfin, les mathématiques de la leçon fournissent des arguments chiffrés et justifiés pour \textbf{conseiller une décision économique réelle}, dans le respect du savoir-faire « rédiger et justifier une démarche, une réponse ou une conclusion ».
\end{itemize}

\newpage

\coursec{Méthodes et exercices résolus}

\coursec{Fonctions logarithme népérien}

\courssub{Définition et premières propriétés}

\begin{definition}[Fonction logarithme népérien]
La fonction \cle{logarithme népérien}, notée $\ln$, est la fonction réciproque de la fonction exponentielle $\exp$ (ou $e^x$). Elle est définie sur $]0\,;\,+\infty[$ et prend ses valeurs dans $\mathbb{R}$.
\begin{itemize}
\item Pour tout $x>0$, $\ln x$ est l'unique réel tel que $e^{\ln x}=x$.
\item Pour tout $y\in\mathbb{R}$, $e^y>0$ et $\ln(e^y)=y$.
\item $\ln 1 = 0$ car $e^0=1$ ; $\ln e = 1$ car $e^1=e$.
\item La fonction $\ln$ est \cle{strictement croissante} sur $]0\,;\,+\infty[$.
\item Graphiquement, la courbe de $\ln$ est l'image de la courbe de $\exp$ par symétrie par rapport à la droite $y=x$ (bisectrice du premier quadrant). Elle coupe l'axe des abscisses en $1$ et admet l'axe des ordonnées ($x=0$) comme asymptote verticale.
\end{itemize}
\end{definition}

\begin{propriete}[Propriétés algébriques fondamentales]
Pour tous nombres réels \cle{strictement positifs} $a$ et $b$, et pour tout entier relatif $n$ (au programme de Terminale technique, on l'utilise pour $n$ entier) :
\begin{align*}
\ln(ab) &= \ln a + \ln b & \text{(logarithme d'un produit)} \\
\ln\!\left(\frac{a}{b}\right) &= \ln a - \ln b & \text{(logarithme d'un quotient)} \\
\ln(a^n) &= n\ln a & \text{(logarithme d'une puissance)} \\
\ln\sqrt{a} &= \tfrac{1}{2}\ln a & \text{(cas particulier $n=\frac{1}{2}$)}
\end{align*}
\emph{Attention :} $\ln(a+b)\neq \ln a+\ln b$ et $\dfrac{\ln a}{\ln b}\neq \ln a-\ln b$. Ces erreurs sont sanctionnées.
\end{propriete}

\courssub{Propriétés algébriques du logarithme népérien}

\begin{methode}[Calculer et simplifier avec les propriétés de ln]
Pour transformer une expression contenant des logarithmes népériens :
\begin{enumerate}
\item Vérifier d'abord que tous les nombres sous le logarithme sont \cle{strictement positifs} (condition d'existence).
\item Appliquer les formules fondamentales, pour $a>0$, $b>0$ et $n$ entier :
\[ \ln(ab)=\ln a+\ln b \qquad \ln\!\left(\frac{a}{b}\right)=\ln a-\ln b \qquad \ln(a^n)=n\ln a \qquad \ln\sqrt{a}=\tfrac{1}{2}\ln a \]
\item Se rappeler les valeurs de référence : $\ln 1 = 0$ et $\ln e = 1$.
\item Pour résoudre une équation, regrouper les logarithmes en un seul, puis utiliser : $\ln A = \ln B \iff A = B$ (avec $A>0$, $B>0$), ou $\ln A = k \iff A = e^{k}$.
\end{enumerate}
\end{methode}

\begin{exoresolu}[Simplification et équation avec ln]
\textbf{Énoncé.} On pose $A = \ln 12 + \ln 3 - \ln 4$.
\begin{enumerate}
\item Écrire $A$ sous la forme $\ln a$ où $a$ est un entier, puis en fonction de $\ln 3$.
\item Résoudre dans $\mathbb{R}$ l'équation $\ln(x+1) + \ln(x-1) = \ln 8$.
\end{enumerate}
\tcblower
\textbf{Solution détaillée.}
\begin{enumerate}
\item Tous les nombres $12$, $3$, $4$ sont strictement positifs : les logarithmes existent. On applique les propriétés algébriques :
\begin{align*}
A &= \ln 12 + \ln 3 - \ln 4 = \ln(12\times 3) - \ln 4 = \ln\!\left(\frac{36}{4}\right) = \ln 9.
\end{align*}
Or $9 = 3^2$, donc $A = \ln(3^2) = 2\ln 3$.
\item \textbf{Conditions d'existence :} il faut $x+1>0$ et $x-1>0$, c'est-à-dire $x>1$. On cherche donc des solutions dans $]1\,;\,+\infty[$.

Pour $x>1$, on regroupe : $\ln(x+1)+\ln(x-1)=\ln\big((x+1)(x-1)\big)=\ln(x^2-1)$.
L'équation devient :
\[ \ln(x^2-1)=\ln 8 \iff x^2-1=8 \iff x^2=9 \iff x=3 \ \text{ou}\ x=-3. \]
\textbf{Vérification des conditions :} $x=-3$ n'appartient pas à $]1\,;\,+\infty[$, il est rejeté. $x=3$ convient : $\ln 4 + \ln 2 = \ln 8$. \checkmark

\textbf{Conclusion :} l'équation admet une unique solution, $S=\{3\}$.
\end{enumerate}
\end{exoresolu}

\courssub{Limites faisant intervenir ln}

\begin{propriete}[Limites de référence de la fonction ln]
\begin{align*}
\lim_{x\to 0^+}\ln x &= -\infty \\
\lim_{x\to +\infty}\ln x &= +\infty \\
\lim_{x\to +\infty}\frac{\ln x}{x} &= 0 \quad \text{(croissance de $\ln$ plus lente que toute fonction affine)} \\
\lim_{x\to 0}\frac{\ln(1+x)}{x} &= 1 \quad \text{(limite remarquable, base de la dérivée de $\ln$ en 1)}
\end{align*}
\end{propriete}

\begin{methode}[Lever les indéterminations avec ln]
\begin{enumerate}
\item Connaître par cœur les \cle{limites de référence} ci-dessus.
\item Face à une forme indéterminée « $\infty - \infty$ » ou « $\infty\times 0$ », \textbf{factoriser par le terme dominant} (souvent $x$) pour faire apparaître $\dfrac{\ln x}{x}$.
\item Pour $\ln u(x)$, déterminer d'abord la limite de $u(x)$, puis composer avec la limite de référence.
\item Conclure en interprétant graphiquement si demandé : limite infinie en $0^+$ $\Rightarrow$ asymptote verticale ; limite finie en $+\infty$ $\Rightarrow$ asymptote horizontale.
\end{enumerate}
\end{methode}

\begin{exoresolu}[Limites avec logarithme népérien]
\textbf{Énoncé.} Soit $f$ la fonction définie par $f(x)=x-\ln x$.
\begin{enumerate}
\item Préciser l'ensemble de définition de $f$.
\item Calculer $\displaystyle\lim_{x\to 0^+} f(x)$ et interpréter graphiquement.
\item Calculer $\displaystyle\lim_{x\to +\infty} f(x)$.
\end{enumerate}
\tcblower
\textbf{Solution détaillée.}
\begin{enumerate}
\item $\ln x$ n'existe que pour $x>0$, donc $D_f = ]0\,;\,+\infty[$.
\item Quand $x\to 0^+$ : $x\to 0$ et $\ln x \to -\infty$, donc $-\ln x \to +\infty$. Par somme :
\[ \lim_{x\to 0^+} f(x) = 0 - (-\infty) = +\infty. \]
\textbf{Interprétation :} la droite d'équation $x=0$ (l'axe des ordonnées) est \cle{asymptote verticale} à la courbe de $f$.
\item Forme indéterminée « $+\infty - \infty$ » : on factorise par $x$, le terme dominant :
\[ f(x) = x\left(1-\frac{\ln x}{x}\right). \]
Or $\displaystyle\lim_{x\to+\infty}\frac{\ln x}{x}=0$ (limite de référence), donc $\displaystyle\lim_{x\to+\infty}\left(1-\frac{\ln x}{x}\right)=1$. Par produit :
\[ \lim_{x\to +\infty} f(x) = +\infty \times 1 = +\infty. \]
\textbf{Conclusion :} $\displaystyle\lim_{x\to 0^+}f(x)=+\infty$ (asymptote verticale $x=0$) et $\displaystyle\lim_{x\to +\infty}f(x)=+\infty$.
\end{enumerate}
\end{exoresolu}

\courssub{Dérivation des fonctions définies à l'aide de ln}

\begin{propriete}[Dérivée de la fonction ln et de ses composées]
\begin{itemize}
\item Sur $]0\,;\,+\infty[$, la fonction $\ln$ est dérivable et $(\ln x)' = \dfrac{1}{x}$.
\item Si $u$ est une fonction dérivable et strictement positive sur un intervalle $I$, alors $\ln\circ u$ est dérivable sur $I$ et :
\[ \big(\ln(u(x))\big)' = \frac{u'(x)}{u(x)}. \]
\end{itemize}
\end{propriete}

\begin{methode}[Calculer la dérivée d'une fonction avec ln]
\begin{enumerate}
\item Identifier la forme de la fonction : $\ln x$ seul, produit, quotient, ou composée $\ln(u)$.
\item Formules à retenir :
\[ (\ln x)' = \frac{1}{x} \qquad \big(\ln u\big)' = \frac{u'}{u} \qquad (uv)'=u'v+uv' \qquad \left(\frac{u}{v}\right)'=\frac{u'v-uv'}{v^2}. \]
\item Pour $\ln(u)$ : calculer d'abord $u'$, puis écrire $\dfrac{u'}{u}$ \textbf{sans oublier} le facteur $u'$.
\item Simplifier le résultat et étudier son \cle{signe} en tenant compte du domaine de définition (sur $]0\,;\,+\infty[$, $x>0$ et $\ln x$ a le signe de $x-1$ : $\ln x > 0 \iff x>1$).
\end{enumerate}
\end{methode}

\begin{exoresolu}[Dérivation de fonctions avec ln]
\textbf{Énoncé.} Dériver les fonctions suivantes sur leur ensemble de définition :
\[ f(x)=x\ln x - x \qquad g(x)=\ln(x^2+1) \qquad h(x)=\frac{\ln x}{x}. \]
\tcblower
\textbf{Solution détaillée.}
\begin{itemize}
\item \textbf{Pour $f$ :} $D_f=]0\,;\,+\infty[$. On dérive le produit $x\ln x$ avec $(uv)'=u'v+uv'$ où $u=x$, $v=\ln x$ :
\[ (x\ln x)' = 1\times\ln x + x\times\frac{1}{x} = \ln x + 1. \]
Donc $f'(x) = \ln x + 1 - 1 = \ln x$. Résultat remarquable : $f$ est une primitive de $\ln$.
\item \textbf{Pour $g$ :} $x^2+1>0$ pour tout réel, donc $D_g=\mathbb{R}$. Forme $\ln u$ avec $u=x^2+1$, $u'=2x$ :
\[ g'(x)=\frac{u'}{u}=\frac{2x}{x^2+1}. \]
\item \textbf{Pour $h$ :} $D_h=]0\,;\,+\infty[$. Forme $\dfrac{u}{v}$ avec $u=\ln x$, $u'=\dfrac{1}{x}$, $v=x$, $v'=1$ :
\[ h'(x)=\frac{u'v-uv'}{v^2}=\frac{\dfrac{1}{x}\times x - \ln x \times 1}{x^2}=\frac{1-\ln x}{x^2}. \]
\end{itemize}
\textbf{Conclusion :} $f'(x)=\ln x$ ; $g'(x)=\dfrac{2x}{x^2+1}$ ; $h'(x)=\dfrac{1-\ln x}{x^2}$.
\end{exoresolu}

\courssub{Primitives des fonctions de la forme u'/u}

\begin{propriete}[Primitive de la forme u'/u]
Soit $u$ une fonction dérivable et strictement positive sur un intervalle $I$. Une primitive de $\dfrac{u'}{u}$ sur $I$ est la fonction $\ln u$ (plus généralement $\ln|u|$ si $u$ peut changer de signe, mais au Probatoire on précise l'intervalle où $u>0$).
\[ \int \frac{u'(x)}{u(x)}\,dx = \ln u(x) + C \quad (u(x)>0). \]
Si le numérateur est $k\,u'(x)$ ($k$ constante), une primitive de $\dfrac{k\,u'}{u}$ est $k\ln u$.
\end{propriete}

\begin{methode}[Reconnaître et primitiver u'/u]
\begin{enumerate}
\item Vérifier que la fonction s'écrit (ou peut s'écrire) sous la forme $\dfrac{u'(x)}{u(x)}$ avec $u(x)>0$ sur l'intervalle considéré.
\item Appliquer la formule : une primitive de $\dfrac{u'}{u}$ est $\ln u$.
\item Si le numérateur est $k\,u'$ (à une constante multiplicative près), écrire $\dfrac{u'}{u}=\dfrac{1}{k}\times\dfrac{k\,u'}{u}$ : une primitive de $\dfrac{k\,u'}{u}$ est $k\ln u$.
\item Ne pas oublier la constante d'intégration $C$ si aucune condition n'est donnée, et vérifier en dérivant le résultat.
\end{enumerate}
\end{methode}

\begin{exoresolu}[Primitives de la forme u'/u]
\textbf{Énoncé.}
\begin{enumerate}
\item Déterminer une primitive de $f(x)=\dfrac{3}{x}$ sur $]0\,;\,+\infty[$.
\item Déterminer une primitive de $g(x)=\dfrac{2x}{x^2+4}$ sur $\mathbb{R}$.
\item Déterminer la primitive $F$ de $h(x)=\dfrac{1}{x\ln x}$ sur $]1\,;\,+\infty[$ telle que $F(e)=1$.
\end{enumerate}
\tcblower
\textbf{Solution détaillée.}
\begin{enumerate}
\item $f(x)=3\times\dfrac{1}{x}=3\times\dfrac{u'}{u}$ avec $u(x)=x>0$ sur $]0\,;\,+\infty[$. Une primitive de $\dfrac{1}{x}$ est $\ln x$, donc :
\[ F(x)=3\ln x + C. \]
\item Posons $u(x)=x^2+4$ ; alors $u'(x)=2x$, qui est exactement le numérateur. De plus $u(x)=x^2+4>0$ pour tout réel. Donc :
\[ G(x)=\ln(x^2+4)+C. \]
\item Sur $]1\,;\,+\infty[$, $\ln x > 0$ donc $x\ln x > 0$ : la fonction est bien définie. On remarque que $\dfrac{1}{x\ln x}=\dfrac{\frac{1}{x}}{\ln x}=\dfrac{u'}{u}$ avec $u(x)=\ln x$ et $u'(x)=\dfrac{1}{x}$. Donc :
\[ F(x)=\ln(\ln x)+C. \]
Condition $F(e)=1$ : $\ln(\ln e)+C=\ln 1 + C = 0 + C = 1$, donc $C=1$.
\[ F(x)=\ln(\ln x)+1. \]
\end{enumerate}
\textbf{Conclusion :} $F(x)=3\ln x + C$ ; $G(x)=\ln(x^2+4)+C$ ; $F(x)=\ln(\ln x)+1$.
\end{exoresolu}

\courssub{Étude de fonctions faisant intervenir ln et applications concrètes}

\begin{methode}[Plan d'étude d'une fonction avec ln]
\begin{enumerate}
\item \textbf{Ensemble de définition} : imposer la stricte positivité des arguments des logarithmes.
\item \textbf{Limites aux bornes} du domaine (utiliser les limites de référence) et interprétation graphique (asymptotes).
\item \textbf{Dérivée} : calculer $f'$, la mettre sous forme factorisée ou de quotient.
\item \textbf{Signe de $f'$} : sur $]0\,;\,+\infty[$, un dénominateur comme $x$ ou $x^2$ est positif ; le signe dépend du numérateur. Rappel : $\ln x \geq 0 \iff x\geq 1$ ; $\ln x \geq k \iff x\geq e^{k}$.
\item \textbf{Tableau de variations} complet (bornes, valeurs remarquables, extremums).
\item \textbf{Conclusion} : extremum éventuel, position par rapport à une droite, réponse au problème concret posé.
\end{enumerate}
\end{methode}

\begin{exoresolu}[Étude complète de $f(x)=x-2\ln x$]
\textbf{Énoncé.} Soit $f$ la fonction définie par $f(x)=x-2\ln x$ et $(\mathcal{C})$ sa courbe représentative.
\begin{enumerate}
\item Déterminer l'ensemble de définition de $f$ et calculer les limites aux bornes.
\item Calculer $f'(x)$ et dresser le tableau de variations de $f$.
\item En déduire le minimum de $f$.
\end{enumerate}
\tcblower
\textbf{Solution détaillée.}
\begin{enumerate}
\item $\ln x$ exige $x>0$ : $D_f=]0\,;\,+\infty[$.

En $0^+$ : $x\to 0$ et $-2\ln x \to +\infty$, donc $\displaystyle\lim_{x\to 0^+}f(x)=+\infty$ (asymptote verticale $x=0$).

En $+\infty$ : forme indéterminée, on factorise : $f(x)=x\left(1-\dfrac{2\ln x}{x}\right)$. Comme $\displaystyle\lim_{x\to+\infty}\dfrac{\ln x}{x}=0$, la parenthèse tend vers $1$, donc $\displaystyle\lim_{x\to+\infty}f(x)=+\infty$.
\item $f$ est dérivable sur $]0\,;\,+\infty[$ et :
\[ f'(x)=1-\frac{2}{x}=\frac{x-2}{x}. \]
Sur $]0\,;\,+\infty[$, $x>0$, donc $f'(x)$ a le signe de $x-2$ : $f'(x)<0$ sur $]0\,;\,2[$, $f'(2)=0$, $f'(x)>0$ sur $]2\,;\,+\infty[$.

Valeur remarquable : $f(2)=2-2\ln 2 \approx 0{,}61$.

\begin{center}
\begin{tabular}{|c|ccccc|}\hline
$x$ & $0$ & & $2$ & & $+\infty$ \\ \hline
$f'(x)$ & & $-$ & $0$ & $+$ & \\ \hline
$f$ & $+\infty$ & $\searrow$ & $2-2\ln 2$ & $\nearrow$ & $+\infty$ \\ \hline
\end{tabular}
\end{center}
\item Le tableau montre que $f$ admet en $x=2$ un \cle{minimum} absolu valant $f(2)=2-2\ln 2 \approx 0{,}61$. Comme ce minimum est strictement positif, on en déduit que $f(x)>0$ pour tout $x>0$.
\end{enumerate}
\textbf{Conclusion :} $f$ est décroissante sur $]0\,;\,2]$, croissante sur $[2\,;\,+\infty[$, et admet un minimum strictement positif en $x=2$.
\end{exoresolu}

\begin{methode}[Résoudre un problème concret avec ln]
\begin{enumerate}
\item \textbf{Traduire} l'énoncé (coût, intensité, capital, désintégration, charge/décharge RC...) en équation ou inéquation contenant $\ln$ ou une exponentielle.
\item \textbf{Isoler} le logarithme ou passer au logarithme : $e^{x}=a \iff x=\ln a$ (pour $a>0$) ; $a^{x}=b \iff x\ln a = \ln b \iff x=\dfrac{\ln b}{\ln a}$.
\item Pour une inéquation, utiliser la \cle{croissance} de $\ln$ : $\ln A < \ln B \iff A<B$ (avec $A,B>0$) — le sens de l'inégalité est conservé.
\item \textbf{Conclure} par une phrase rédigée dans le contexte, avec l'unité adaptée et un arrondi cohérent avec la question.
\end{enumerate}
\end{methode}

\begin{exoresolu}[Temps de charge d'un condensateur]
\textbf{Énoncé.} Dans un atelier d'électricité, la tension aux bornes d'un condensateur en charge est modélisée par $u(t)=12\left(1-e^{-0{,}5t}\right)$ (en volts), où $t$ est le temps en secondes. Le condensateur est considéré comme chargé lorsque $u(t)\geq 11$ V.
\begin{enumerate}
\item Montrer que résoudre $u(t)\geq 11$ revient à $e^{-0{,}5t}\leq \dfrac{1}{12}$.
\item En déduire le temps minimal de charge, arrondi au dixième de seconde.
\end{enumerate}
\tcblower
\textbf{Solution détaillée.}
\begin{enumerate}
\item On part de l'inéquation et on isole l'exponentielle :
\begin{align*}
u(t)\geq 11 &\iff 12\left(1-e^{-0{,}5t}\right)\geq 11 \\
&\iff 1-e^{-0{,}5t}\geq \frac{11}{12} \\
&\iff -e^{-0{,}5t}\geq \frac{11}{12}-1 = -\frac{1}{12} \\
&\iff e^{-0{,}5t}\leq \frac{1}{12} \quad (\text{on divise par } -1 < 0 : \text{le sens change}).
\end{align*}
\item La fonction $\ln$ est strictement croissante sur $]0\,;\,+\infty[$, elle conserve donc le sens de l'inégalité (et $\frac{1}{12}>0$, le logarithme existe) :
\begin{align*}
e^{-0{,}5t}\leq \frac{1}{12} &\iff -0{,}5t \leq \ln\!\left(\frac{1}{12}\right) = -\ln 12 \\
&\iff t \geq \frac{\ln 12}{0{,}5} = 2\ln 12.
\end{align*}
Calcul : $2\ln 12 \approx 2\times 2{,}4849 \approx 4{,}97$.
\end{enumerate}
\textbf{Conclusion :} le condensateur est chargé au bout d'environ $5{,}0$ secondes (temps minimal $t = 2\ln 12 \approx 5{,}0$ s).
\end{exoresolu}

\begin{savaistu}[Le logarithme népérien dans la technique et l'histoire]
\begin{itemize}
\item \textbf{Histoire :} John Napier (Neper) publie ses tables de logarithmes en 1614 pour simplifier les calculs astronomiques (multiplications $\to$ additions). Le logarithme « népérien » (base $e \approx 2,718$) s'impose naturellement en analyse car sa dérivée est $1/x$ et il est l'inverse de l'exponentielle $e^x$ (Euler, 1748).
\item \textbf{Applications techniques :}
\begin{itemize}
\item \textbf{Électricité (circuits RC/RL) :} charge/décharge exponentielle, constante de temps $\tau$. $u(t)=U_0(1-e^{-t/\tau})$. Le temps pour atteindre un pourcentage donné se calcule avec $\ln$.
\item \textbf{Chimie (pH) :} $\mathrm{pH} = -\log_{10}[\mathrm{H}^+] = -\frac{\ln[\mathrm{H}^+]}{\ln 10}$. Mesure d'acidité.
\item \textbf{Acoustique (niveau sonore) :} $L = 20\log_{10}(P/P_0)$ décibels (dB). Échelle logarithmique car l'oreille perçoit les rapports de pression.
\item \textbf{Sismologie (échelle de Richter) :} $M = \log_{10}(A/A_0)$. Magnitude d'un séisme.
\item \textbf{Économie/Gestion :} Intérêts composés continus $C(t)=C_0 e^{rt}$, temps de doublement $t = \ln 2 / r$.
\end{itemize}
\end{itemize}
\end{savaistu}

\begin{attention}[Les 5 erreurs qui coûtent des points au Probatoire]
\begin{enumerate}
\item \textbf{Oublier les conditions d'existence.} Avant tout calcul avec $\ln(x-2)$ ou toute équation contenant des logarithmes, écrire le domaine : ici $x>2$. Une solution comme $x=-3$ trouvée pour $\ln(x^2-1)=\ln 8$ doit être \emph{rejetée} si elle ne vérifie pas les conditions.
\item \textbf{Inventer des formules fausses.} $\ln(a+b)\neq \ln a + \ln b$ et $\dfrac{\ln a}{\ln b}\neq \ln a - \ln b$. Seules les formules du cours sont valables : $\ln(ab)=\ln a+\ln b$, $\ln\!\left(\dfrac{a}{b}\right)=\ln a-\ln b$, $\ln(a^n)=n\ln a$.
\item \textbf{Oublier le $u'$ dans la dérivée de $\ln u$.} La dérivée de $\ln(x^2+1)$ est $\dfrac{2x}{x^2+1}$, et non $\dfrac{1}{x^2+1}$. De même, une primitive de $\dfrac{1}{2x+1}$ est $\dfrac{1}{2}\ln(2x+1)$ : penser au coefficient.
\item \textbf{Mal lever les formes indéterminées.} « $+\infty-\infty$ » ne donne jamais $0$ automatiquement : il faut factoriser par le terme dominant et utiliser $\displaystyle\lim_{x\to+\infty}\dfrac{\ln x}{x}=0$. Écrire la factorisation explicitement sur la copie.
\item \textbf{Changer le sens d'une inégalité à tort.} La fonction $\ln$ est \emph{croissante} : elle conserve le sens des inégalités. En revanche, multiplier ou diviser par un nombre \emph{négatif} (comme $-0{,}5$ ou $-1$) change le sens. Bien distinguer les deux situations et justifier chaque étape.
\end{enumerate}
\end{attention}

\newpage

\coursec{Exercices}

\courssub{Je vérifie mes connaissances}

\begin{exercice}[QCM : Propriétés du logarithme népérien]
Parmi les affirmations suivantes, une seule est exacte pour tout $x>0$ et $y>0$. Laquelle ?
\begin{enumerate}
    \item $\ln(x+y) = \ln x + \ln y$
    \item $\ln(x-y) = \ln x - \ln y$
    \item $\ln\left(\dfrac{x}{y}\right) = \ln x - \ln y$
    \item $\ln(xy) = \ln x \times \ln y$
\end{enumerate}
\end{exercice}

\begin{exercice}[Vrai ou Faux ? Justifiez]
Déterminez si les affirmations suivantes sont vraies ou fausses. Justifiez votre réponse par un calcul ou un contre-exemple.
\begin{enumerate}
    \item $\ln 1 = 0$.
    \item $\ln e^2 = 2$.
    \item $\ln(-1)$ est un nombre réel.
    \item Pour tout $x>0$, $\ln(x^3) = 3\ln x$.
    \item La fonction $x \mapsto \ln x$ est décroissante sur $]0;+\infty[$.
\end{enumerate}
\end{exercice}

\begin{exercice}[Ensemble de définition et calculs directs]
\begin{enumerate}
    \item Déterminez l'ensemble de définition $D_f$ de la fonction $f$ définie par $f(x) = \ln(x^2 - 4)$.
    \item Calculez les valeurs exactes des expressions suivantes :
    \begin{itemize}
        \item $A = \ln 50 - 2\ln 5 + \ln 2$ (écrire le résultat sous la forme $k\ln 2$).
        \item $B = \ln(\sqrt{e}) + \ln(e^3) - \ln(e^2)$.
        \item $C = \ln\left(\dfrac{1}{e^4}\right)$.
    \end{itemize}
\end{enumerate}
\end{exercice}

\begin{exercice}[Limites et dérivées élémentaires]
\begin{enumerate}
    \item Calculez les limites suivantes :
    \begin{itemize}
        \item $\displaystyle\lim_{x\to +\infty} (x^2 - \ln x)$
        \item $\displaystyle\lim_{x\to 1} \dfrac{\ln x}{x-1}$
        \item $\displaystyle\lim_{x\to 0^+} x^2 \ln x$
    \end{itemize}
    \item Soit $g(x) = \ln(3x+1)$ sur $]-\frac{1}{3};+\infty[$. Calculez $g'(x)$.
\end{enumerate}
\end{exercice}

\courssub{Je m'entraîne}

\begin{exercice}[Équation et inéquation logarithmiques]
\begin{enumerate}
    \item Résolvez dans $\mathbb{R}$ l'équation : $\ln(x+1) + \ln(x-1) = \ln 8$.
    \item Résolvez dans $\mathbb{R}$ l'inéquation : $\ln(2x-3) \le 1$.
\end{enumerate}
\end{exercice}

\begin{exercice}[Étude de fonction : $f(x) = \ln(x^2+x+1)$]
Soit la fonction $f$ définie sur $\mathbb{R}$ par $f(x) = \ln(x^2+x+1)$.
\begin{enumerate}
    \item Calculez $f'(x)$.
    \item Étudiez le signe de $f'(x)$ et dressez le tableau de variations de $f$.
    \item Calculez les limites de $f$ en $-\infty$ et $+\infty$.
    \item La courbe $\mathcal{C}_f$ admet-elle des asymptotes ? Justifiez.
\end{enumerate}
\end{exercice}

\begin{exercice}[Primitives de la forme $u'/u$]
Déterminez l'ensemble des primitives sur l'intervalle où le dénominateur est strictement positif des fonctions suivantes :
\begin{enumerate}
    \item $f(x) = \dfrac{3x^2+2}{x^3+2x+1}$
    \item $g(x) = \dfrac{2x-1}{x^2-x+5}$
    \item $h(x) = \dfrac{e^x}{e^x+2}$
\end{enumerate}
\end{exercice}

\begin{exercice}[Application : Modèle de décroissance radioactive]
La masse $m(t)$ (en grammes) d'un échantillon radioactif au temps $t$ (en années) est modélisée par $m(t) = m_0 e^{-kt}$ avec $m_0=100$ g et $k=0,02$ an$^{-1}$.
\begin{enumerate}
    \item Exprimez $t$ en fonction de $m(t)$ à l'aide du logarithme népérien.
    \item Calculez le temps nécessaire pour que la masse passe de 100 g à 10 g (demi-vie non demandée, temps pour 10\% de la masse initiale).
    \item Vérifiez que la demi-vie $T_{1/2}$ (temps pour que la masse soit divisée par 2) vaut $\dfrac{\ln 2}{k}$ et donnez sa valeur approchée au dixième d'année près.
\end{enumerate}
\end{exercice}

\courssub{Je me prépare au Probatoire}

\begin{exercice}[Étude complète : $f(x) = x - 2\ln x$ (Type Probatoire C/E)]
Soit la fonction $f$ définie sur $]0;+\infty[$ par $f(x) = x - 2\ln x$. On note $\mathcal{C}_f$ sa courbe représentative dans un repère orthonormé.
\begin{enumerate}
    \item Calculez les limites de $f$ en $0^+$ et en $+\infty$. En déduire l'existence d'éventuelles asymptotes.
    \item Calculez la dérivée $f'(x)$ et étudiez son signe sur $]0;+\infty[$.
    \item Dressez le tableau de variations complet de $f$.
    \item Déterminez l'équation de la tangente $T$ à $\mathcal{C}_f$ au point d'abscisse $1$.
    \item Montrez que la courbe $\mathcal{C}_f$ se situe au-dessus de sa tangente $T$ en $1$.
    \item Tracez $\mathcal{C}_f$ et $T$ dans le même repère (échelle : 2 cm pour 1 unité sur les axes).
\end{enumerate}
\end{exercice}

\begin{exercice}[Problème : Croissance démographique (Type Probatoire)]
La population $P(t)$ (en habitants) d'une ville $t$ années après 2020 est modélisée par la fonction $P(t) = 50\,000 \, e^{0,03t}$.
\begin{enumerate}
    \item Calculez la population en 2020 et en 2030.
    \item Déterminez le temps $t_d$ (en années) nécessaire pour que la population double. Donnez la valeur exacte puis une valeur approchée au dixième d'année près.
    \item En quelle année la population dépassera-t-elle 200 000 habitants ? (Rédigez la démarche complète : inéquation, résolution, conclusion).
    \item On définit la fonction $f(t) = \ln P(t)$. Montrer que $f$ est une fonction affine et donnez son expression. Quelle est l'interprétation du coefficient directeur ?
\end{enumerate}
\end{exercice}

\begin{exercice}[Inégalité fondamentale et suites (Type Probatoire)]
\begin{enumerate}
    \item Soit la fonction $h$ définie sur $]0;+\infty[$ par $h(x) = x - 1 - \ln x$.
    \begin{itemize}
        \item Étudiez les variations de $h$ et dressez son tableau de variations.
        \item En déduire que pour tout $x>0$, $\ln x \le x-1$.
        \item Interprétez géométriquement cette inégalité par rapport à la courbe de la fonction logarithme népérien et à sa tangente au point d'abscisse 1.
    \end{itemize}
    \item Soit $(u_n)$ une suite géométrique de premier terme $u_0 = 3$ et de raison $q = 1,5$.
    \begin{itemize}
        \item Montrer que la suite $(v_n)$ définie par $v_n = \ln u_n$ est une suite arithmétique. Précisez sa raison.
        \item Déterminer le plus petit entier naturel $n$ tel que $u_n > 100$.
    \end{itemize}
\end{enumerate}
\end{exercice}

\newpage

\coursec{Corrigés des exercices}

\begin{corrige}[Exercice 1 — QCM : Propriétés du logarithme népérien]
L'unique affirmation exacte pour tout $x>0$ et $y>0$ est la \textbf{proposition 3} :
\[
\ln\left(\dfrac{x}{y}\right) = \ln x - \ln y.
\]
\par\medskip
\textbf{Justifications :}
\begin{itemize}
    \item \textbf{1. Faux :} $\ln(x+y) \neq \ln x + \ln y$. Contre-exemple : $x=y=1$, $\ln(2) \neq 0+0$.
    \item \textbf{2. Faux :} $\ln(x-y)$ n'est même pas toujours défini ($x>y$ nécessaire) et $\ln(x-y) \neq \ln x - \ln y$. Contre-exemple : $x=3, y=1$, $\ln 2 \neq \ln 3 - 0$.
    \item \textbf{3. Vrai :} Propriété fondamentale du logarithme népérien : le logarithme d'un quotient est la différence des logarithmes.
    \item \textbf{4. Faux :} $\ln(xy) = \ln x + \ln y$ (somme) et non le produit $\ln x \times \ln y$. Contre-exemple : $x=y=e$, $\ln(e^2)=2$ alors que $\ln e \times \ln e = 1$.
\end{itemize}
\end{corrige}

\begin{corrige}[Exercice 2 — Vrai ou Faux ? Justifiez]
\begin{enumerate}
    \item \textbf{Vrai.} Par définition, $\ln 1 = 0$ car $e^0 = 1$.
    \item \textbf{Vrai.} $\ln e^2 = 2 \ln e = 2 \times 1 = 2$.
    \item \textbf{Faux.} La fonction logarithme népérien n'est définie que sur $]0;+\infty[$. $-1 \notin D_{\ln}$, donc $\ln(-1)$ n'est pas un nombre réel (c'est un nombre complexe).
    \item \textbf{Vrai.} Pour tout $x>0$, $\ln(x^3) = 3\ln x$ (propriété $\ln(x^a)=a\ln x$).
    \item \textbf{Faux.} La fonction $x \mapsto \ln x$ est \textbf{strictement croissante} sur $]0;+\infty[$ (sa dérivée $1/x > 0$).
\end{enumerate}
\end{corrige}

\begin{corrige}[Exercice 3 — Ensemble de définition et calculs directs]
\begin{enumerate}
    \item $f(x) = \ln(x^2-4)$. La condition d'existence est $x^2-4 > 0 \iff (x-2)(x+2) > 0$.
    Le trinôme est positif en dehors de ses racines :
    \[
    D_f = ]-\infty;-2[ \;\cup\; ]2;+\infty[.
    \]
    \item \textbf{Calculs exacts :}
    \begin{itemize}
        \item $A = \ln 50 - 2\ln 5 + \ln 2 = \ln 50 - \ln(5^2) + \ln 2 = \ln 50 - \ln 25 + \ln 2$.
        \[
        A = \ln\left(\frac{50}{25}\right) + \ln 2 = \ln 2 + \ln 2 = 2\ln 2.
        \]
        Donc $k=2$.
        \item $B = \ln(\sqrt{e}) + \ln(e^3) - \ln(e^2) = \ln(e^{1/2}) + 3\ln e - 2\ln e = \frac{1}{2} + 3 - 2 = \frac{3}{2}$.
        \item $C = \ln\left(\dfrac{1}{e^4}\right) = \ln(e^{-4}) = -4 \ln e = -4$.
    \end{itemize}
\end{enumerate}
\end{corrige}

\begin{corrige}[Exercice 4 — Limites et dérivées élémentaires]
\begin{enumerate}
    \item \textbf{Limites :}
    \begin{itemize}
        \item $\displaystyle\lim_{x\to +\infty} (x^2 - \ln x)$.
        On factorise par $x^2$ : $x^2\left(1 - \frac{\ln x}{x^2}\right)$. Comme $\lim_{+\infty} \frac{\ln x}{x^2} = 0$, la limite est $+\infty \times 1 = +\infty$.
        \item $\displaystyle\lim_{x\to 1} \dfrac{\ln x}{x-1}$.
        C'est le taux d'accroissement de $\ln$ en $1$, ou la dérivée de $\ln$ en $1$ : $(\ln)'(1) = \frac{1}{1} = 1$.
        \item $\displaystyle\lim_{x\to 0^+} x^2 \ln x$.
        Forme indéterminée $0 \times (-\infty)$. On pose $X = \ln x$, alors $x = e^X$ et $x\to 0^+ \Rightarrow X\to -\infty$.
        $x^2 \ln x = e^{2X} \times X = \frac{X}{e^{-2X}}$. Comme l'exponentielle domine tout polynôme, la limite est $0$.
        (Règle connue : $\forall \alpha>0,\; \lim_{0^+} x^\alpha \ln x = 0$).
    \end{itemize}
    \item $g(x) = \ln(3x+1)$ sur $]-\frac{1}{3};+\infty[$.
    Dérivée : $g'(x) = \dfrac{(3x+1)'}{3x+1} = \dfrac{3}{3x+1}$.
\end{enumerate}
\end{corrige}

\begin{corrige}[Exercice 5 — Équation et inéquation logarithmiques]
\begin{enumerate}
    \item $\ln(x+1) + \ln(x-1) = \ln 8$.
    \textbf{Conditions :} $x+1>0$ et $x-1>0 \iff x>1$.
    \textbf{Résolution :}
    \[
    \ln\big((x+1)(x-1)\big) = \ln 8 \iff \ln(x^2-1) = \ln 8 \iff x^2-1 = 8 \iff x^2 = 9 \iff x = \pm 3.
    \]
    Seule $x=3$ satisfait la condition $x>1$.
    \textbf{Solution :} $S = \{3\}$.
    \item $\ln(2x-3) \le 1$.
    \textbf{Conditions :} $2x-3 > 0 \iff x > \frac{3}{2}$.
    \textbf{Résolution :} La fonction $\ln$ est croissante, donc :
    \[
    2x-3 \le e^1 \iff 2x \le e+3 \iff x \le \frac{e+3}{2}.
    \]
    Intersection avec la condition : $S = \left]\frac{3}{2}\,;\,\frac{e+3}{2}\right]$.
\end{enumerate}
\end{corrige}

\begin{corrige}[Exercice 6 — Étude de fonction : $f(x) = \ln(x^2+x+1)$]
\begin{enumerate}
    \item $f(x) = \ln(u(x))$ avec $u(x)=x^2+x+1$. $u'(x)=2x+1$.
    $u(x) > 0$ pour tout $x\in\mathbb{R}$ ($\Delta = 1-4 = -3 < 0$, coefficient directeur $>0$).
    \[
    f'(x) = \frac{u'(x)}{u(x)} = \frac{2x+1}{x^2+x+1}.
    \]
    \item Le dénominateur $x^2+x+1$ est strictement positif sur $\mathbb{R}$.
    Le signe de $f'(x)$ est donc le signe du numérateur $2x+1$.
    \begin{itemize}
        \item $f'(x) = 0 \iff 2x+1=0 \iff x = -\frac{1}{2}$.
        \item $f'(x) < 0$ sur $]-\infty;-\frac{1}{2}[$ (décroissante).
        \item $f'(x) > 0$ sur $]-\frac{1}{2};+\infty[$ (croissante).
    \end{itemize}
    \textbf{Tableau de variations :}
    \[
    \begin{tabular}{|c|ccccc|}\hline
    $x$ & $-\infty$ & & $-\frac{1}{2}$ & & $+\infty$ \\ \hline
    $f'(x)$ & & $-$ & $0$ & $+$ & \\ \hline
    $f$ & $+\infty$ & $\searrow$ & $\ln\frac{3}{4}$ & $\nearrow$ & $+\infty$ \\ \hline
    \end{tabular}
    \]
    Valeur minimale : $f(-\frac{1}{2}) = \ln\left(\frac{1}{4}-\frac{1}{2}+1\right) = \ln\frac{3}{4} = \ln 3 - 2\ln 2$.
    \item \textbf{Limites :}
    $\lim_{-\infty} (x^2+x+1) = +\infty \implies \lim_{-\infty} f(x) = +\infty$.
    $\lim_{+\infty} (x^2+x+1) = +\infty \implies \lim_{+\infty} f(x) = +\infty$.
    \item \textbf{Asymptotes :}
    \begin{itemize}
        \item Pas d'asymptote verticale (domaine $\mathbb{R}$).
        \item Asymptote oblique ? $\lim_{\pm\infty} \frac{f(x)}{x} = \lim_{\pm\infty} \frac{\ln(x^2+x+1)}{x} = 0$ (croissance logarithmique vs linéaire).
        Si asymptote oblique $y=ax+b$, alors $a=0$. Or $\lim_{\pm\infty} (f(x)-0) = +\infty \neq b$ (fini).
        Donc \textbf{pas d'asymptote oblique}.
    \end{itemize}
\end{enumerate}
\end{corrige}

\begin{corrige}[Exercice 7 — Primitives de la forme $u'/u$]
On cherche une primitive de la forme $\ln|u(x)| + C$ (valable sur chaque intervalle où $u$ ne s'annule pas).
\begin{enumerate}
    \item $f(x) = \dfrac{3x^2+2}{x^3+2x+1}$.
    Posons $u(x) = x^3+2x+1$. Alors $u'(x) = 3x^2+2$.
    $f(x) = \frac{u'(x)}{u(x)}$.
    Le polynôme $u$ est strictement croissant ($u'>0$) et admet une unique racine réelle $\alpha \in ]-1;0[$.
    Une primitive sur l'ensemble de définition $D_f = \mathbb{R}\setminus\{\alpha\}$ est :
    \[
    F(x) = \ln|x^3+2x+1| + C \quad (C \in \mathbb{R}).
    \]
    \item $g(x) = \dfrac{2x-1}{x^2-x+5}$.
    Posons $u(x) = x^2-x+5$. Alors $u'(x) = 2x-1$.
    $\Delta = 1-20 = -19 < 0$, donc $u(x) > 0$ sur $\mathbb{R}$.
    Une primitive sur $\mathbb{R}$ est :
    \[
    G(x) = \ln(x^2-x+5) + C \quad (C \in \mathbb{R}).
    \]
    \item $h(x) = \dfrac{e^x}{e^x+2}$.
    Posons $u(x) = e^x+2$. Alors $u'(x) = e^x$.
    $u(x) > 2 > 0$ sur $\mathbb{R}$.
    Une primitive sur $\mathbb{R}$ est :
    \[
    H(x) = \ln(e^x+2) + C \quad (C \in \mathbb{R}).
    \]
\end{enumerate}
\end{corrige}

\begin{corrige}[Exercice 8 — Application : Modèle de décroissance radioactive]
Données : $m(t) = 100 e^{-0,02t}$.
\begin{enumerate}
    \item On isole $t$ :
    \[
    \frac{m(t)}{100} = e^{-0,02t} \iff \ln\left(\frac{m(t)}{100}\right) = -0,02t \iff t = \frac{\ln\left(\frac{100}{m(t)}\right)}{0,02} = 50\,\ln\left(\frac{100}{m(t)}\right).
    \]
    \item Pour $m(t) = 10$ g :
    \[
    t = 50\,\ln\left(\frac{100}{10}\right) = 50\,\ln 10 \approx 50 \times 2,302585 \approx 115,1 \text{ ans}.
    \]
    \item Demi-vie $T_{1/2}$ : $m(T_{1/2}) = \frac{m_0}{2} = 50$.
    \[
    50 = 100 e^{-k T_{1/2}} \iff \frac{1}{2} = e^{-k T_{1/2}} \iff \ln\left(\frac{1}{2}\right) = -k T_{1/2} \iff -\ln 2 = -k T_{1/2}.
    \]
    Donc $\boxed{T_{1/2} = \dfrac{\ln 2}{k}}$.
    Valeur numérique : $T_{1/2} = \dfrac{\ln 2}{0,02} = 50 \ln 2 \approx 50 \times 0,693147 \approx 34,657 \approx \boxed{34,7 \text{ ans}}$ (au dixième près).
\end{enumerate}
\end{corrige}

\begin{corrige}[Exercice 9 — Étude complète : $f(x) = x - 2\ln x$ (Type Probatoire C/E)]
\begin{enumerate}
    \item \textbf{Limites et asymptotes :}
    \begin{itemize}
        \item $\lim_{0^+} f(x) = 0 - 2(-\infty) = +\infty$. La droite $x=0$ (axe des ordonnées) est une \textbf{asymptote verticale}.
        \item $\lim_{+\infty} f(x) = +\infty - 2(+\infty)$. Forme indéterminée. $f(x) = x\left(1 - 2\frac{\ln x}{x}\right)$. Comme $\lim_{+\infty} \frac{\ln x}{x} = 0$, $\lim_{+\infty} f(x) = +\infty \times 1 = +\infty$.
        \item Asymptote oblique ? $a = \lim_{+\infty} \frac{f(x)}{x} = \lim_{+\infty} \left(1 - 2\frac{\ln x}{x}\right) = 1$.
        $b = \lim_{+\infty} (f(x) - x) = \lim_{+\infty} (-2\ln x) = -\infty$.
        Pas d'asymptote oblique (ni horizontale).
    \end{itemize}
    \item \textbf{Dérivée et signe :}
    $f'(x) = 1 - \frac{2}{x} = \frac{x-2}{x}$ sur $]0;+\infty[$.
    \begin{itemize}
        \item $f'(x) = 0 \iff x=2$.
        \item $f'(x) < 0$ sur $]0;2[$ ($f$ décroissante).
        \item $f'(x) > 0$ sur $]2;+\infty[$ ($f$ croissante).
    \end{itemize}
    \item \textbf{Tableau de variations :}
    \[
    \begin{tabular}{|c|ccccc|}\hline
    $x$ & $0^+$ & & $2$ & & $+\infty$ \\ \hline
    $f'(x)$ & & $-$ & $0$ & $+$ & \\ \hline
    $f$ & $+\infty$ & $\searrow$ & $2-2\ln 2$ & $\nearrow$ & $+\infty$ \\ \hline
    \end{tabular}
    \]
    Minimum absolu : $m = f(2) = 2 - 2\ln 2 = 2(1-\ln 2) > 0$ (car $\ln 2 \approx 0,69 < 1$).
    \item \textbf{Tangente en 1 :}
    $f(1) = 1 - 0 = 1$. $f'(1) = 1 - 2 = -1$.
    Équation : $y - 1 = -1(x-1) \iff \boxed{y = -x + 2}$.
    \item \textbf{Position relative courbe/tangente :}
    Soit $\varphi(x) = f(x) - (-x+2) = x - 2\ln x + x - 2 = 2x - 2 - 2\ln x = 2(x-1-\ln x)$.
    Posons $h(x) = x-1-\ln x$ sur $]0;+\infty[$.
    $h'(x) = 1 - \frac{1}{x} = \frac{x-1}{x}$.
    $h$ décroît sur $]0;1]$, croît sur $[1;+\infty[$. Minimum $h(1)=0$.
    Donc $\forall x>0,\; h(x) \ge 0 \implies \varphi(x) \ge 0 \implies f(x) \ge -x+2$.
    Égalité seulement en $x=1$.
    \textbf{Conclusion :} La courbe $\mathcal{C}_f$ se situe \textbf{au-dessus} de sa tangente $T$ en $1$ (contact en $x=1$).
    \item \textbf{Trace (description) :} Repère orthonormé (2 cm/unité). Asymptote verticale $x=0$. Courbe passant par $(1;1)$ avec tangente $y=-x+2$ (passant par $(0;2)$ et $(2;0)$). Minimum en $(2; 2-2\ln 2 \approx 0,61)$. Branche gauche descendant de $+\infty$ vers le min, branche droite montant vers $+\infty$.
\end{enumerate}
\par\medskip
\textbf{Barème indicatif (sur 20) :} Limites/Asymptotes (4) ; Dérivée/Signe (4) ; Tableau variations (3) ; Tangente (3) ; Position relative (4) ; Tracé/Soin (2).
\par\textbf{Points de vigilance :} Citer la condition $x>0$ pour $\ln x$. Justifier le signe de $f'$ par le numérateur/dénominateur. Pour la position relative, définir explicitement la fonction différence et étudier ses variations (ne pas dire « c'est évident »). Tracer l'asymptote $x=0$ en pointillés.
\end{corrige}

\begin{corrige}[Exercice 10 — Problème : Croissance démographique (Type Probatoire)]
Données : $P(t) = 50\,000 \, e^{0,03t}$ ($t$ en années après 2020).
\begin{enumerate}
    \item \textbf{Population :}
    2020 ($t=0$) : $P(0) = 50\,000 \, e^0 = \boxed{50\,000 \text{ habitants}}$.
    2030 ($t=10$) : $P(10) = 50\,000 \, e^{0,3} \approx 50\,000 \times 1,34986 = \boxed{67\,493 \text{ habitants}}$ (arrondi unité).
    \item \textbf{Temps de doublement $t_d$ :}
    $P(t_d) = 2 \times 50\,000 = 100\,000$.
    $50\,000 e^{0,03 t_d} = 100\,000 \iff e^{0,03 t_d} = 2 \iff 0,03 t_d = \ln 2$.
    Valeur exacte : $\boxed{t_d = \dfrac{\ln 2}{0,03}}$ années.
    Valeur approchée : $t_d \approx \dfrac{0,693147}{0,03} \approx 23,1049 \approx \boxed{23,1 \text{ ans}}$ (dixième près).
    \item \textbf{Année dépassement 200 000 :}
    Inéquation : $P(t) > 200\,000 \iff 50\,000 e^{0,03t} > 200\,000 \iff e^{0,03t} > 4$.
    $\ln$ est croissante : $0,03t > \ln 4 = 2\ln 2 \iff t > \dfrac{2\ln 2}{0,03}$.
    $t > \dfrac{2 \times 0,693147}{0,03} \approx 46,21$ ans.
    Le plus petit entier $t$ est $47$.
    Année : $2020 + 47 = \boxed{2067}$.
    \item \textbf{Fonction $f(t) = \ln P(t)$ :}
    $f(t) = \ln(50\,000 \, e^{0,03t}) = \ln 50\,000 + \ln(e^{0,03t}) = \ln 50\,000 + 0,03t$.
    $f$ est une \textbf{fonction affine} de la forme $f(t) = at + b$ avec $a = 0,03$ et $b = \ln 50\,000$.
    \textbf{Interprétation :} Le coefficient directeur $0,03$ (soit $3\%$) représente le \textbf{taux de croissance instantané (ou relatif)} de la population. C'est la dérivée logarithmique : $\frac{P'(t)}{P(t)} = 0,03$.
\end{enumerate}
\par\medskip
\textbf{Barème indicatif (sur 20) :} Calculs populations (3) ; Doublement exact/approché (4) ; Inéquation/Année (5) ; Fonction affine/Interprétation (4) ; Rédaction/Présentation (4).
\par\textbf{Points de vigilance :} Ne pas confondre $t$ (durée) et année civile. Pour l'inéquation, bien écrire les étapes : inéquation $\to$ $\ln$ (croissante) $\to$ isolation $t$ $\to$ conclusion en année. Préciser que $f$ est affine car somme d'une constante et d'un terme proportionnel à $t$.
\end{corrige}

\begin{corrige}[Exercice 11 — Inégalité fondamentale et suites (Type Probatoire)]
\begin{enumerate}
    \item \textbf{Étude de $h(x) = x - 1 - \ln x$ sur $]0;+\infty[$ :}
    $h'(x) = 1 - \frac{1}{x} = \frac{x-1}{x}$.
    \begin{itemize}
        \item $h'(x) < 0$ sur $]0;1[$ $\implies$ $h$ décroissante.
        \item $h'(x) = 0$ en $x=1$.
        \item $h'(x) > 0$ sur $]1;+\infty[$ $\implies$ $h$ croissante.
    \end{itemize}
    \textbf{Tableau de variations :}
    \[
    \begin{tabular}{|c|ccccc|}\hline
    $x$ & $0^+$ & & $1$ & & $+\infty$ \\ \hline
    $h'(x)$ & & $-$ & $0$ & $+$ & \\ \hline
    $h$ & $+\infty$ & $\searrow$ & $0$ & $\nearrow$ & $+\infty$ \\ \hline
    \end{tabular}
    \]
    \textbf{Inégalité :} Le minimum de $h$ est $h(1)=0$. Donc $\forall x>0,\; h(x) \ge 0 \iff x-1-\ln x \ge 0 \iff \boxed{\ln x \le x-1}$.
    \textbf{Interprétation géométrique :} La courbe $\mathcal{C} : y = \ln x$ est \textbf{située en dessous} (ou au plus sur) de sa tangente au point d'abscisse $1$, d'équation $y = x-1$. Elles se confondent seulement au point $(1;0)$.
    \item \textbf{Suite $(u_n)$ géométrique :} $u_0=3$, $q=1,5$. $u_n = 3 \times (1,5)^n$.
    \begin{itemize}
        \item $v_n = \ln u_n = \ln(3 \times 1,5^n) = \ln 3 + n \ln 1,5$.
        $v_n$ s'écrit sous la forme $v_n = n \times \ln 1,5 + \ln 3$.
        C'est une \textbf{suite arithmétique} de raison $\boxed{r = \ln 1,5}$ (et premier terme $v_0 = \ln 3$).
        \item On cherche le plus petit $n \in \mathbb{N}$ tel que $u_n > 100$.
        $3 \times 1,5^n > 100 \iff 1,5^n > \frac{100}{3}$.
        $\ln$ croissante : $n \ln 1,5 > \ln(100/3) \iff n > \dfrac{\ln(100/3)}{\ln 1,5}$.
        Calcul : $\ln(100/3) \approx \ln 33,33 \approx 3,5066$ ; $\ln 1,5 \approx 0,4055$.
        $n > \frac{3,5066}{0,4055} \approx 8,647$.
        Le plus petit entier naturel est $\boxed{n = 9}$.
        (Vérification : $u_8 = 3 \times 1,5^8 \approx 76,9$ ; $u_9 \approx 115,3 > 100$).
    \end{itemize}
\end{enumerate}
\par\medskip
\textbf{Barème indicatif (sur 20) :} Étude $h$/Tableau (5) ; Inégalité déduite (2) ; Interprétation géométrique (3) ; Suite arithmétique (4) ; Inéquation suite/Calcul $n$ (4) ; Rigueur rédactionnelle (2).
\par\textbf{Points de vigilance :} Pour $h$, bien préciser le domaine $]0;+\infty[$ et le signe du dénominateur $x$ pour le signe de $h'$. L'inégalité $\ln x \le x-1$ est fondamentale, la déduire correctement du minimum. Pour la suite, utiliser les propriétés $\ln(ab)=\ln a+\ln b$ et $\ln(a^n)=n\ln a$. Pour l'inéquation, ne pas oublier que $\ln 1,5 > 0$ donc le sens de l'inéquation est conservé.
\end{corrige}

\newpage

\coursec{Fiche bilan}

\coursec{Fonction logarithme népérien}

\begin{popfigure}
\begin{tikzpicture}[
  mindmap,
  grow cyclic,
  every node/.style={concept, execute at begin node=\small\bfseries},
  root concept/.append style={concept color=popblue, minimum size=3.5cm, font=\large\bfseries},
  level 1 concept/.append style={minimum size=2.8cm, sibling angle=60, font=\normalsize\bfseries},
  level 2 concept/.append style={minimum size=2.2cm, font=\small},
  text=popdark
]
\node[root concept] {Fonction $\ln$}
  child[concept color=poporange] { node[level 1 concept, align=center]{Définition\\ \& Propriétés\\ de base}
    child { node[level 2 concept] {$\ln x = \int_1^x \frac{dt}{t}$} }
    child { node[level 2 concept] {Domaine $]0,+\infty[$} }
    child { node[level 2 concept] {Bijection vers $\mathbb{R}$} }
    child { node[level 2 concept] {$\ln 1=0,\ \ln e=1$} }
  }
  child[concept color=popgreen] { node[level 1 concept, align=center]{Propriétés\\ algébriques}
    child { node[level 2 concept] {$\ln(ab)=\ln a+\ln b$} }
    child { node[level 2 concept] {$\ln(\frac{a}{b})=\ln a-\ln b$} }
    child { node[level 2 concept] {$\ln(a^r)=r\ln a$} }
    child { node[level 2 concept] {Inverse : $\exp$} }
  }
  child[concept color=poppurple] { node[level 1 concept] {Limites}
    child { node[level 2 concept] {$\lim_{x\to 0^+}\ln x=-\infty$} }
    child { node[level 2 concept] {$\lim_{x\to+\infty}\ln x=+\infty$} }
    child { node[level 2 concept] {$\lim_{x\to+\infty}\frac{\ln x}{x}=0$} }
    child { node[level 2 concept] {$\lim_{x\to 0^+}x\ln x=0$} }
  }
  child[concept color=poppink] { node[level 1 concept] {Dérivation}
    child { node[level 2 concept] {$(\ln u)'=\frac{u'}{u}$} }
    child { node[level 2 concept] {$(\ln|u|)'=\frac{u'}{u}$} }
    child { node[level 2 concept] {Dérivée de $\ln x = \frac{1}{x}$} }
    child { node[level 2 concept] {Croissance stricte} }
  }
  child[concept color=popgold] { node[level 1 concept] {Primitives}
    child { node[level 2 concept] {$\int \frac{u'}{u}=\ln|u|+C$} }
    child { node[level 2 concept] {$\int \frac{1}{x}=\ln|x|+C$} }
    child { node[level 2 concept] {Recherche forme $u'/u$} }
    child { node[level 2 concept] {Intégration par parties} }
  }
  child[concept color=popblue] { node[level 1 concept, align=center]{Étude\\ complète}
    child { node[level 2 concept] {Domaine, Parité, Périodicité} }
    child { node[level 2 concept] {Limites, Asymptotes} }
    child { node[level 2 concept] {Dérivée, Variations} }
    child { node[level 2 concept] {Tableau, Courbe, Applis.} }
  };
\end{tikzpicture}
\legende{Carte mentale : Fonction logarithme népérien (Terminale Technique)}
\end{popfigure}

\begin{definition}[Définitions clés]
\begin{itemize}
  \item $\ln : ]0,+\infty[ \to \mathbb{R}$ définie par $\ln x = \int_1^x \frac{dt}{t}$.
  \item $e$ unique réel tel que $\ln e = 1$ ($e \approx 2,718$).
  \item Fonction réciproque : l'exponentielle $\exp$ (ou $e^x$).
  \item $\forall x>0,\ \exp(\ln x)=x$ ; $\forall x\in\mathbb{R},\ \ln(e^x)=x$.
\end{itemize}
\end{definition}

\begin{propriete}[Théorèmes et formules à connaître par cœur]
\begin{center}
\begin{tabularx}{\linewidth}{|l|X|}\hline
\rowcolor{popblueL} \textbf{Propriété / Formule} & \textbf{Énoncé / Condition} \\ \hline
Propriétés algébriques & $\ln(ab)=\ln a+\ln b$ ; $\ln(\frac{a}{b})=\ln a-\ln b$ ; $\ln(a^r)=r\ln a$ ($a,b>0, r\in\mathbb{R}$) \\ \hline
Limites usuelles & $\lim_{x\to 0^+}\ln x=-\infty$ ; $\lim_{x\to+\infty}\ln x=+\infty$ ; $\lim_{x\to+\infty}\frac{\ln x}{x}=0$ ; $\lim_{x\to 0^+}x\ln x=0$ \\ \hline
Dérivée de $\ln u$ & Si $u>0$ dérivable : $(\ln u)' = \frac{u'}{u}$ \\ \hline
Dérivée de $\ln|u|$ & Si $u\neq 0$ dérivable : $(\ln|u|)' = \frac{u'}{u}$ (étend le domaine) \\ \hline
Primitive de $\frac{u'}{u}$ & $\int \frac{u'(x)}{u(x)}dx = \ln|u(x)| + C$ \\ \hline
Primitive de $\frac{1}{x}$ & $\int \frac{dx}{x} = \ln|x| + C$ \\ \hline
\end{tabularx}
\end{center}
\end{propriete}

\begin{methode}[Méthodes express]
\begin{enumerate}
  \item \textbf{Simplifier une expression} : Factoriser les coefficients devant les $\ln$ en puissances ($\ln a^r=r\ln a$), regrouper sommes/différences en produit/quotient.
  \item \textbf{Résoudre $\ln(f(x))=k$} : Vérifier $f(x)>0$, puis $f(x)=e^k$. Résoudre l'équation résultante.
  \item \textbf{Étudier $f(x)=\ln(u(x))$} :
    \begin{itemize}
      \item Domaine : $u(x)>0$.
      \item Dérivée : $f'(x)=\frac{u'(x)}{u(x)}$.
      \item Signe de $f'$ = signe de $u'$ (car $u>0$ sur le domaine).
      \item Limites aux bornes du domaine.
    \end{itemize}
  \item \textbf{Calculer $\int \frac{P'(x)}{P(x)}dx$} : Identifier $u=P(x)$, résultat $\ln|P(x)|+C$.
  \item \textbf{Intégration par parties avec $\ln$} : Poser $u=\ln(\dots)$, $dv=dx$ (ou $x^n dx$). $du$ devient rationnel.
\end{enumerate}
\end{methode}

\begin{savaistu}[Applications concrètes du logarithme népérien]
\begin{itemize}
  \item \textbf{Chimie (pH)} : $\mathrm{pH} = -\log_{10}[\mathrm{H}^+] = -\frac{\ln[\mathrm{H}^+]}{\ln 10}$.
  \item \textbf{Sismologie (échelle de Richter)} : Magnitude $M = \log_{10}(A/A_0) = \frac{\ln(A/A_0)}{\ln 10}$.
  \item \textbf{Physique (radioactivité)} : Loi de décroissance $N(t)=N_0 e^{-\lambda t} \Rightarrow t = \frac{1}{\lambda}\ln\frac{N_0}{N}$.
  \item \textbf{Économie / Démographie} : Croissance continue $P(t)=P_0 e^{rt} \Rightarrow t = \frac{1}{r}\ln\frac{P}{P_0}$.
  \item \textbf{Électricité (charge/décharge condensateur)} : $u(t)=U_0(1-e^{-t/RC})$ ou $u(t)=U_0 e^{-t/RC}$.
\end{itemize}
\end{savaistu}

\begin{aretenir}[Checklist avant le Probatoire]
$\square$ Savoir écrire la définition intégrale de $\ln$ et en déduire $\ln 1=0$, $\ln e=1$.\par
$\square$ Appliquer sans hésiter les 3 propriétés algébriques (produit, quotient, puissance).\par
$\square$ Connaître les 4 limites de référence ($0^+$, $+\infty$, $\frac{\ln x}{x}$, $x\ln x$) et les justifier.\par
$\square$ Dériver $\ln(u(x))$ et $\ln|u(x)|$ : écrire explicitement $u'$ et $u$ avant de conclure.\par
$\square$ Poser correctement le domaine de définition AVANT de dériver ou d'étudier une fonction $\ln$.\par
$\square$ Construire un tableau de variations complet : domaine, limites, dérivée, signe, variations, valeurs remarquables.\par
$\square$ Identifier la forme $\frac{u'}{u}$ pour calculer une primitive (penser au $\ln|u|$).\par
$\square$ Maîtriser l'intégration par parties avec $u=\ln(g(x))$ (choix standard).\par
$\square$ Résoudre équations et inéquations logarithmiques : condition de positivité $\to$ passage à l'exponentielle $\to$ vérification.\par
$\square$ Interpréter graphiquement : asymptote verticale à la borne exclue du domaine, concavité (souvent convexe).\par
$\square$ Modéliser un problème concret (croissance, décroissance radioactive, pH, échelle de Richter) avec $\ln$ ou $\exp$.\par
$\square$ Rédiger une démonstration rigoureuse (ex: monotonie de $\ln$, inégalité $\ln x \le x-1$).
\end{aretenir}

\findecours

\end{document}
